% concatenated from main.tex
\documentclass[11pt,reqno]{amsart}
\usepackage{amssymb,mathrsfs,color,mathtools,empheq, verbatim, epstopdf}
\usepackage{pinlabel}
\mathtoolsset{showonlyrefs}
%\usepackage[notref,notcite]{showkeys}
%\usepackage{fullpage} 
\usepackage{hyperref} 

\usepackage{enumerate, bbm}

\usepackage{graphicx}
%\usepackage{xcolor} % A package to add color.
\usepackage{tensor}
\usepackage{slashed}
\usepackage{cite}
\usepackage[dvipsnames]{xcolor}

\usepackage[shortlabels]{enumitem}

\usepackage{geometry}\geometry{margin=1.1in}


\numberwithin{equation}{section}
\setlength{\parindent}{0pt}
\newtheorem{mainthm}{Theorem}
\newtheorem{thm}{Theorem}[section]
\newtheorem{cor}[thm]{Corollary}
\newtheorem{lem}[thm]{Lemma}
\newtheorem{prop}[thm]{Proposition}
\newtheorem{ques}{Question}
\newtheorem{con}{Conjecture}

\theoremstyle{definition} 
\newtheorem{rem}[thm]{Remark}
\newtheorem{defn}[thm]{Definition}

%\newtheorem{rem}[thm]{Remark}
\theoremstyle{remark}
\newtheorem{claim}[thm]{Claim}
\newtheorem{ex}[thm]{Example}



%\usepackage{empheq}

%\usepackage[T1]{fontenc}
%\usepackage[shortlabels]{enumitem}
%\usepackage[notref,notcite]{showkeys}



\def\bC {\mathbb{C}}
\def\bN {\mathbb{N}}
\def\bQ {\mathbb{Q}}
\def\bR {\mathbb{R}}
\def\bS {\mathbb{S}}
\def\bT {\mathbb{T}}
\def\bZ {\mathbb{Z}}

\def\fA {\mathfrak{A}}
\def\fB {\mathfrak{B}}
\def\fH {\mathfrak{H}}
\def\fv {\mathfrak{v}}
\def\frakh{\mathfrak{h}}

\def\cA {\mathcal{A}}
\def\cB {\mathcal{B}}
%\def\cC {\mathcal{C}}
\def\cD {\mathcal{D}}
\def\cE {\mathcal{E}}
\def\cF {\mathcal{F}}
\def\cG {\mathcal{G}}
\def\cH {\mathcal{H}}
\def\cJ {\mathcal{J}}
\def\cK {\mathcal{K}}
\def\cL {\mathcal{L}}
\def\cM {\mathcal{M}}
\def\cN {\mathcal{N}}
\def\cP {\mathcal{P}}
\def\cQ {\mathcal{Q}}
\def\cR {\mathcal{R}}
%\def\cS {\mathcal{S}}
\def\cT {\mathcal{T}}
\def\cW {\mathcal{W}}
\def\cX {\mathcal{X}}
\def\cY {\mathcal{Y}}
\def\cZ {\mathcal{Z}}

\def\scrL{\mathscr{L}}

\def\grad {{\nabla}}

\def\rstr {{\big |}}

\newcommand{\la}{\langle}
\newcommand{\ra}{\rangle}
\newcommand{\La}{\big\langle}
\newcommand{\Ra}{\big\rangle}
\def \lA {\big\langle \! \! \big\langle}
\def \rA {\big\rangle \! \! \big\rangle}
\def \LA {\bigg\langle \! \! \! \! \! \; \bigg\langle}
\def \RA {\bigg\rangle \! \! \! \! \! \; \bigg\rangle}

\def\pp {{\parallel}}

\def\vide {{\varnothing}}


%%%%%%%%%%%%%%%%%%%%%%%%%%%%%%%%%%%%COMMANDES
\newcommand{\red}[1]{\textcolor{red}{#1}}

\newcommand{\tx}[1]{\mathrm{#1}}
\newcommand{\wto}{\rightharpoonup}
\newcommand{\sto}[1]{\xrightarrow[#1]{}}

\newcommand{\wh}[1]{\widehat{#1}}
\newcommand{\wt}[1]{\widetilde{#1}}
\newcommand{\bs}[1]{\boldsymbol{#1}}
%\newcommand{\conj}[1]{\overline{#1}}

\newcommand{\vvec}{\overrightarrow}
\renewcommand{\div}{\operatorname{div}}
\newcommand{\curl}{\operatorname{curl}}
\newcommand{\Dom}{\operatorname{Dom}}
\newcommand{\sign}{\operatorname{sign}}
\newcommand{\spn}{\operatorname{span}}
\newcommand{\supp}{\operatorname{supp}}
\newcommand{\Tr}{\operatorname{Tr}}
\newcommand{\codim}{\operatorname{codim}}
\newcommand{\dist}{\operatorname{dist}}
\renewcommand{\ker}{\operatorname{ker}}
\newcommand{\im}{\operatorname{im}}
\newcommand{\Id}{\operatorname{Id}}
\newcommand{\Lip}{\operatorname{Lip}}
\newcommand{\eee}{\mathrm e}
\newcommand{\iii}{\mathrm i}

\DeclareMathOperator{\sech}{sech}

%\newcommand{\ud}{\mathrm{\,d}}
\newcommand{\vd}{\mathrm{d}}
\newcommand{\udr}{\,r\vd r}
\newcommand{\vD}{\mathrm{D}}
%\newcommand{\dd}[1]{{\frac{\vd}{\vd{#1}}}}
\newcommand{\pd}[1]{{\frac{\partial}{\partial{#1}}}}

%\renewcommand{\dot}[1]{\accentset{\mbox{\large .}}{#1}}

%%%%%%%%%%%%%%%%%%%%%%%%% AD-HOC MACROS
\newcommand{\enorm}{{\dot H^1 \times L^2}}
\newcommand{\edual}{{L^2 \times \dot H^{-1}}}
\newcommand{\uln}[1]{{\underline{ #1 }}}
\newcommand{\lin}{_{\textsc{l}}}
\newcommand{\nl}{\textsc{nl}}
\newcommand{\linl}{_{\textsc{lin}, \lambda}}
\newcommand{\yp}{{\cY}^+}
\newcommand{\ym}{{\cY}^-}
\newcommand{\ypl}{{\cY}^+_{\lambda}}
\newcommand{\yml}{{\cY}^-_{\lambda}}
\newcommand{\fb}{\mathfrak b}
\newcommand{\Hloc}{\dot H^1_{\tx{loc}}}
\newcommand{\Lin}{\textsc{l}} 
%%%%%%%%%%%%%%%%%%%%%%%%%%%%%%%%%%%%%%%%%%%%%%%%%%%%%%%%%%%%%%%%%%%%%%%%%%%%%%%%%%%%%%%%%%%%%%%%%%%%%%%%%

\newcommand{\conj}[1]{\overline{#1}}
\newcommand{\sh}[1]{#1^\sharp}
\newcommand{\fl}[1]{#1^\flat}
\newcommand{\nt}[1]{#1^\natural}


\newcommand{\dd}[1]{\frac{\ud}{\ud{#1}}}






%------------------------------------------------
%AMS LaTeX Paper:  Andy's preamble
%------------------------------------------------

%Pictures
%\usepackage{wrapfig}
%\usepackage{tikz}
%\usetikzlibrary{arrows,calc,decorations.pathreplacing}
%\definecolor{light-gray1}{gray}{0.90}
%\definecolor{light-gray2}{gray}{0.80}

%\parindent=0pt
%\allowdisplaybreaks

%colors
\definecolor{deepgreen}{cmyk}{1,0,1,0.5}
\newcommand{\Blue}[1]{{\color{blue} #1}}
\newcommand{\Red}[1]{{\color{red} #1}}
\newcommand{\Green}[1]{{\color{deepgreen} #1}}

%greenoff
%\renewcommand{\Green}[1]{#1}

%mathcal letters
\newcommand{\A}{\mathcal{A}}
\newcommand{\B}{\mathcal{B}}
%\newcommand{\calD}{\mathcal{D}}
\newcommand{\E}{\mathcal{E}}
\newcommand{\F}{\mathcal{F}}
\newcommand{\G}{\mathcal{G}}
\newcommand{\HH}{\mathcal{H}}
\newcommand{\J}{\mathcal{J}}
\newcommand{\K}{\mathcal{K}}
\newcommand{\LL}{\mathcal{L}}
\newcommand{\NN}{\mathcal{N}}
\newcommand{\M}{\mathcal{M}}
\newcommand{\cS}{\mathcal{S}}
\newcommand{\T}{\mathcal{T}}
\newcommand{\QQ}{\mathcal{Q}}
\newcommand{\cC}{\mathcal{C}}
\newcommand{\RR}{\mathcal{R}}


%\mathscri letters
\newcommand{\CC}{\mathscr{C}}
\newcommand{\Dm}{\mathscr{D}}
\newcommand{\Es}{\mathscr{E}}
\newcommand{\FF}{\mathscr{F}}
\newcommand{\JJ}{\mathscr{J}}
\newcommand{\sS}{\mathscr{S}}
\newcommand{\Hs}{\mathscr{H}}
\newcommand{\Vs}{\mathscr{V}}
\newcommand{\Ks}{\mathcal{K}}
\newcommand{\Ls}{\mathscr{L}}


%mathbb letters
\newcommand{\C}{\mathbb{C}}
\newcommand{\Hp}{\mathbb{H}}
\newcommand{\N}{\mathbb{N}}
\newcommand{\R}{\mathbb{R}}
\newcommand{\Sp}{\mathbb{S}}
\newcommand{\Z}{\mathbb{Z}}
\newcommand{\D}{\mathbb{D}}

\newcommand{\HL}{\mathfrak{H}}
\newcommand{\Sg}{\mathfrak{S}}


%bold lowercase
\newcommand{\g}{\mathbf{g}}
\newcommand{\h}{\mathbf{h}}
\newcommand{\etab}{\boldsymbol{\eta}}
\newcommand{\m}{\mathbf{m}}

%bold uppercase
\newcommand{\Db}{\mathbf{D}}


%greek lowercase
\newcommand{\al}{\alpha}
\newcommand{\be}{\beta}
\newcommand{\ga}{\gamma}
\newcommand{\de}{\delta}
\newcommand{\e}{\varepsilon}
\newcommand{\fy}{\varphi}
\newcommand{\om}{\omega}
\newcommand{\lam}{\lambda}
%\newcommand{\lam}{\lambda} 
\newcommand{\te}{\theta}
\newcommand{\vt}{\vartheta}
\newcommand{\s}{\sigma}
\newcommand{\ta}{\tau}
\newcommand{\ka}{\kappa}
\newcommand{\x}{\xi}
\newcommand{\y}{\eta}
\newcommand{\z}{\zeta}
\newcommand{\rh}{\rho}
\newcommand{\io}{\iota}
\newcommand{\si}{\varsigma}


%greek uppercase
\newcommand{\De}{\Delta}
\newcommand{\Om}{\Omega}
\newcommand{\Ga}{\Gamma}
%\newcommand{\La}{\Lambda}
\newcommand{\Lam}{\Lambda} 
\newcommand{\Sig}{\Sigma}
\newcommand{\Te}{\Theta} 
\newcommand{\Fi}{\Phi} 


%misc symbols
\newcommand{\p}{\partial}
\newcommand{\na}{\nabla}
\newcommand{\re}{\mathop{\mathrm{Re}}}
%\newcommand{\im}{\mathop{\mathrm{Im}}}
\newcommand{\weak}{\operatorname{w-}}
%\newcommand{\supp}{\operatorname{supp}}
\newcommand{\Cu}{\bigcup}
\newcommand{\Ca}{\bigcap}
\newcommand{\loc}{\operatorname{loc}}
\newcommand{\ext}{\operatorname{ext}}
\DeclareMathOperator*{\esssup}{ess\,sup}

\makeatletter
\newcommand{\rmnum}[1]{\romannumeral #1}
\newcommand{\Rmnum}[1]{\expandafter\@slowromancap\romannumeral #1@}
\makeatother

\newcommand{\gr}[1]{\lceil #1 \rfloor}
\newcommand{\lec}{\lesssim}
\newcommand{\gec}{\gtrsim}
\newcommand{\slec}{\,\lesssim\,}
\newcommand{\sgec}{\,\gtrsim\,}
\newcommand{\precc}{\preccurlyeq}
\newcommand{\IN}[1]{\text{ in }#1}
\newcommand{\Bdd}[1]{\text{ bounded in }#1}
\newcommand{\etc}{,\ldots,}
\newcommand{\I}{\infty}
\newcommand{\da}{\dagger}

\newcommand{\ti}{\widetilde}
\newcommand{\ba}{\overline}
\newcommand{\ha}{\widehat}
\newcommand{\U}{\underline}

%bracketing
\newcommand{\LR}[1]{{\langle #1 \rangle}}
\newcommand{\ang}[1]{\left\langle{#1}\right\rangle}
\newcommand{\abs}[1]{\left\lvert{#1}\right\rvert}
\newcommand{\ds}{\displaystyle}

 %math environment shortcuts
\newcommand{\ali}[1]{\begin{align}\begin{split} #1 \end{split}\end{align}}
\newcommand{\ant}[1]{\begin{align*}\begin{split} #1 \end{split}\end{align*}}
\newcommand{\EQ}[1]{\begin{equation}\begin{split} #1 \end{split}\end{equation}}
\newcommand{\pmat}[1]{\begin{pmatrix} #1 \end{pmatrix}}

\setlength{\marginparwidth}{2cm}
\newcommand{\sidenote}[1]{\marginpar[\raggedleft\tiny #1]{\raggedright\tiny #1}}
\newcommand{\Del}[1]{}
\newcommand{\CAS}[1]{\begin{cases} #1 \end{cases}}
\newcommand{\pt}{&}
\newcommand{\pr}{\\ &}
\newcommand{\pq}{\quad}
\newcommand{\pn}{}
\newcommand{\prq}{\\ &\quad}
\newcommand{\prQ}{\\ &\qquad}
\newcommand{\prQQ}{\\ &\qquad\qquad}




%\renewcommand{\theenumi}{\roman{enumi}}

\newcommand{\mb}{\mathfrak{m}_\phi}
\newcommand{\tm}{\tilde{\mathfrak{m}}_\de\,}
\newcommand{\diff}[1]{{\triangleleft #1}}
\newcommand{\pa}{\triangleright}

%text shortcuts
\newcommand{\mwhere}{{\ \ \text{where} \ \ }}
\newcommand{\mand}{{\ \ \text{and} \ \  }}
\newcommand{\mor}{{\ \ \text{or} \ \ }}
\newcommand{\mif}{{\ \ \text{if} \ \ }}
\newcommand{\mfor}{{\ \ \text{for} \ \ }}
\newcommand{\mas}{{\ \ \text{as} \ \ }}
\newcommand{\inn}{{\ \ \text{in} \ \ }}
\newcommand{\euc}{\mathrm{euc}}
\newcommand{\Lline}{\LL_{V_{\la,k}}}
\newcommand{\Qu}{Q_{\la,k}}
\newcommand{\Lv}{\ti{\LL}_{V_{\la,k}}}



%\newcommand{\sign}{\operatorname{sign}}
\newcommand{\Ext}{E_{\operatorname{ext}}}
\newcommand{\V}{\vec}
% \newcommand{\dvol}{\operatorname{dVol}} 
\newcommand{\uD}{\operatorname{D}}

\newcommand{\VV}{\mathcal{V}}
\newcommand{\WW}{\mathcal{W}}
\renewcommand\Re{\mathrm{Re}\,}
\renewcommand\Im{\mathrm{Im}\,}
\def\glei{\mathrm{eq}}

\newcommand{\atanh}{\mathrm{arctanh}}
\newcommand{\atan}{\mathrm{arctan}}
\newcommand{\YM}{\mathrm{YM}}
\newcounter{parts}

\newcommand{\so}{\mathfrak{so}}
\newcommand{\ad}{\mathrm{ad}}
\newcommand{\at}{\big|_}

\newcommand{\rrho}{\rho}
\newcommand{\ssigma}{\sigma}
\newcommand{\ttau}{\tau}
%\usepackage{showkeys}

\newcommand{\rst}{\!\upharpoonright}



\newcommand{\XP}{{}^{(X)} \! P }
\newcommand{\oP}{{}^{(X_0)} \! P }
\newcommand{\eP}{{}^{(X_{\eps})} \! P }
\newcommand{\enP}{{}^{(X_{\eps_n})} \! P }
\newcommand{\XJ}{{}^{(X)} \! J }
\newcommand{\tJ}{{}^{(\p_t)} \! J }
\newcommand{\oJ}{{}^{(X_0)} \! J }
\newcommand{\Xp}{{}^{(X)} \! \pi }
\newcommand{\xop}{{}^{(X_0)} \! \pi }
\newcommand{\di}{\textrm{div}}
\newcommand{\UL}{\underline{L}}
\newcommand{\slp}{\slashed{\p}}
\newcommand{\ven}{v_{\eps_n}}
\newcommand{\uen}{u_{\eps_n}}

%\newcommand{\euc}{\mathrm{euc}}
\newcommand{\rdr}{\, r\mathrm{d}r}
\newcommand{\dr}{\, \mathrm{d}r}
\newcommand{\rest}{\!\!\restriction}


\newcommand{\ED}{\operatorname{ED}}



% COLORS ------------------------------------------------------------
\definecolor{green}{rgb}{0.3,0.8,0} % Redefines the color green.
%%%% Annotations %%%%
\newcommand{\texthl}[1]{\begingroup\setlength\fboxsep{1pt}\colorbox{yellow}{#1}\endgroup} % Defines the command "\texthl{<text>}" to be the text of its argument highlighted in yellow.
\newcommand{\displayhl}[1]{\begingroup\setlength\fboxsep{1pt}\colorbox{yellow}{$\displaystyle #1$}\endgroup} % Defines the command "\displayhl{<displayed math>}" to be the displayed mathematics of its argument highlighted in yellow.
\newcommand{\scripthl}[1]{{\setlength\fboxsep{1pt}\colorbox{yellow}{$\scriptstyle #1$}}} % Defines the command "\scripthl{<superscript or subscript>}" to be the superscript or subscript of its argument highlighted in yellow.
%\newcommand{\comment}[1]{\begingroup\color{red} #1\endgroup} % Defines the command "\comment{}" to be its argument written in red font.
\newcommand{\green}[1]{\begingroup\color{green} #1\endgroup}

% THEOREMS -------------------------------------------------------
%\newtheorem{theorem}{Theorem}[section]
%\newtheorem{corollary}[theorem]{Corollary}
%\newtheorem{lemma}[theorem]{Lemma}
%\newtheorem{proposition}[theorem]{Proposition}
%\theoremstyle{definition}
%\newtheorem{definition}[theorem]{Definition}
%\newtheorem{example}[theorem]{Example}
%\theoremstyle{remark}
%\newtheorem{remark}[theorem]{Remark}
%\numberwithin{equation}{section}
%\numberwithin{equation}{subsection}

% MATH -----------------------------------------------------------
\newcommand{\nrm}[1]{\Vert#1\Vert}
%\newcommand{\abs}[1]{\vert#1\vert}
\newcommand{\brk}[1]{\langle#1\rangle}
\newcommand{\set}[1]{\{#1\}}
\newcommand{\tld}[1]{\widetilde{#1}}
\newcommand{\br}[1]{\overline{#1}}
\newcommand{\ubr}[1]{\underline{#1}}
\newcommand{\wht}[1]{\widehat{#1}}

%\newcommand{\dist}{\mathrm{dist}}
\newcommand{\sgn}{{\mathrm{sgn}}}
\newcommand{\tr}{\textrm{tr}}

\renewcommand{\Re}{\mathrm{Re}}
\renewcommand{\Im}{\mathrm{Im}}

\newcommand{\aeq}{\sim}
\newcommand{\aleq}{\lesssim}
\newcommand{\ageq}{\gtrsim}

\newcommand{\lap}{\triangle}
\newcommand{\bx}{\Box}

\newcommand{\ud}{\mathrm{d}}
\newcommand{\rd}{\partial}
\newcommand{\nb}{\nabla}

\newcommand{\imp}{\Rightarrow}
\newcommand{\pmi}{\Leftarrow}
\newcommand{\impmi}{\Leftrightarrow}

\newcommand{\bb}{\Big}

\newcommand{\0}{\emptyset}

%Greek Characters
\newcommand{\alp}{\alpha}
\newcommand{\bt}{\beta}
\newcommand{\gmm}{\gamma}
\newcommand{\Gmm}{\Gamma}
\newcommand{\dlt}{\delta}
\newcommand{\Dlt}{\Delta}
\newcommand{\eps}{\epsilon}
\newcommand{\veps}{\varepsilon}
\newcommand{\kpp}{\kappa}
\newcommand{\lmb}{\lambda}
\newcommand{\Lmb}{\Lambda}
\newcommand{\sgm}{\sigma}
\newcommand{\Sgm}{\Sigma}
\newcommand{\tht}{\theta}
\newcommand{\Tht}{\Theta}
\newcommand{\vtht}{\vartheta}
\newcommand{\omg}{\omega}
\newcommand{\Omg}{\Omega}
\newcommand{\ups}{\upsilon}
\newcommand{\zt}{\zeta}
\newcommand{\Zt}{\Zeta}

%Bold Characters
\newcommand{\bfa}{{\bf a}}
\newcommand{\bfb}{{\bf b}}
\newcommand{\bfc}{{\bf c}}
\newcommand{\bfd}{{\bf d}}
\newcommand{\bfe}{{\bf e}}
\newcommand{\bff}{{\bf f}}
\newcommand{\bfg}{{\bf g}}
\newcommand{\bfh}{{\bf h}}
\newcommand{\bfi}{{\bf i}}
\newcommand{\bfj}{{\bf j}}
\newcommand{\bfk}{{\bf k}}
\newcommand{\bfl}{{\bf l}}
\newcommand{\bfm}{{\bf m}}
\newcommand{\bfn}{{\bf n}}
\newcommand{\bfo}{{\bf o}}
\newcommand{\bfp}{{\bf p}}
\newcommand{\bfq}{{\bf q}}
\newcommand{\bfr}{{\bf r}}
\newcommand{\bfs}{{\bf s}}
\newcommand{\bft}{{\bf t}}
\newcommand{\bfu}{{\bf u}}
\newcommand{\bfv}{{\bf v}}
\newcommand{\bfw}{{\bf w}}
\newcommand{\bfx}{{\bf x}}
\newcommand{\bfy}{{\bf y}}
\newcommand{\bfz}{{\bf z}}

\newcommand{\hbfd}{\widehat{\bfd}}

\newcommand{\bfA}{{\bf A}}
\newcommand{\bfB}{{\bf B}}
\newcommand{\bfC}{{\bf C}}
\newcommand{\bfD}{{\bf D}}
\newcommand{\bfE}{{\bf E}}
\newcommand{\bfF}{{\bf F}}
\newcommand{\bfG}{{\bf G}}
\newcommand{\bfH}{{\bf H}}
\newcommand{\bfI}{{\bf I}}
\newcommand{\bfJ}{{\bf J}}
\newcommand{\bfK}{{\bf K}}
\newcommand{\bfL}{{\bf L}}
\newcommand{\bfM}{{\bf M}}
\newcommand{\bfN}{{\bf N}}
\newcommand{\bfO}{{\bf O}}
\newcommand{\bfP}{{\bf P}}
\newcommand{\bfQ}{{\bf Q}}
\newcommand{\bfR}{{\bf R}}
\newcommand{\bfS}{{\bf S}}
\newcommand{\bfT}{{\bf T}}
\newcommand{\bfU}{{\bf U}}
\newcommand{\bfV}{{\bf V}}
\newcommand{\bfW}{{\bf W}}
\newcommand{\bfX}{{\bf X}}
\newcommand{\bfY}{{\bf Y}}
\newcommand{\bfZ}{{\bf Z}}

%Blackboard Bold Characters
\newcommand{\bbA}{\mathbb A}
\newcommand{\bbB}{\mathbb B}
\newcommand{\bbC}{\mathbb C}
\newcommand{\bbD}{\mathbb D}
\newcommand{\bbE}{\mathbb E}
\newcommand{\bbF}{\mathbb F}
\newcommand{\bbG}{\mathbb G}
\newcommand{\bbH}{\mathbb H}
\newcommand{\bbI}{\mathbb I}
\newcommand{\bbJ}{\mathbb J}
\newcommand{\bbK}{\mathbb K}
\newcommand{\bbL}{\mathbb L}
\newcommand{\bbM}{\mathbb M}
\newcommand{\bbN}{\mathbb N}
\newcommand{\bbO}{\mathbb O}
\newcommand{\bbP}{\mathbb P}
\newcommand{\bbQ}{\mathbb Q}
\newcommand{\bbR}{\mathbb R}
\newcommand{\bbS}{\mathbb S}
\newcommand{\bbT}{\mathbb T}
\newcommand{\bbU}{\mathbb U}
\newcommand{\bbV}{\mathbb V}
\newcommand{\bbW}{\mathbb W}
\newcommand{\bbX}{\mathbb X}
\newcommand{\bbY}{\mathbb Y}
\newcommand{\bbZ}{\mathbb Z}

%MathCal Characters
\newcommand{\calA}{\mathcal A}
\newcommand{\calB}{\mathcal B}
\newcommand{\calC}{\mathcal C}
\newcommand{\calD}{\mathcal D}
\newcommand{\calE}{\mathcal E}
\newcommand{\calF}{\mathcal F}
\newcommand{\calG}{\mathcal G}
\newcommand{\calH}{\mathcal H}
\newcommand{\calI}{\mathcal I}
\newcommand{\calJ}{\mathcal J}
\newcommand{\calK}{\mathcal K}
\newcommand{\calL}{\mathcal L}
\newcommand{\calM}{\mathcal M}
\newcommand{\calN}{\mathcal N}
\newcommand{\calO}{\mathcal O}
\newcommand{\calP}{\mathcal P}
\newcommand{\calQ}{\mathcal Q}
\newcommand{\calR}{\mathcal R}
\newcommand{\calS}{\mathcal S}
\newcommand{\calT}{\mathcal T}
\newcommand{\calU}{\mathcal U}
\newcommand{\calV}{\mathcal V}
\newcommand{\calW}{\mathcal W}
\newcommand{\calX}{\mathcal X}
\newcommand{\calY}{\mathcal Y}
\newcommand{\calZ}{\mathcal Z}

\newcommand{\To}{\longrightarrow}
%\newcommand{\BX}{\mathbf{B}(X)}
%\newcommand{\A}{\mathcal{A}}

\newcommand{\arctanh}{\mathrm{arctanh} \, }
\newcommand{\vartht}{\vartheta}

\newcommand{\ztinf}{\zeta_{\infty}}
\newcommand{\ztz}{\zeta_{0}}
\newcommand{\gapE}{\mu^{2}_{\lmb,k}}
\newcommand{\rnL}{\widetilde{\calL}}



% ----------------------------------------------------------------

\vfuzz2pt % Don't report over-full v-boxes if over-edge is small
\hfuzz2pt % Don't report over-full h-boxes if over-edge is small

%-----------------------------------------------------------------



\newcommand{\vs}{v^*}
\newcommand{\us}{u^*}
\newcommand{\bvs}{\bs v^*}
\newcommand{\bus}{\bs u^*}


\newcommand{\ula}{\underline{\lambda}}
\newcommand{\umu}{\underline{\mu}}


%\newcommand{\uln}[1]{{\underline{ #1 }}}
%\newcommand{\lin}{_{\textsc{l}}}
%\newcommand{\linl}{_{\textsc{lin}, \lambda}}


%---------------------- Local commands ------------------------------%
\newcommand{\met}{G} 
\newcommand{\phii}{f_i} 
\newcommand{\phic}{f_c} 
\newcommand{\phia}{F_a}


\newcommand{\ULam}{\U{\Lam}} 
\newcommand{\ulam}{\U{\lam}} 

\newcommand{\app}{\operatorname{app}}
\newcommand{\form}{\operatorname{form}}

\newcommand{\aaa}{\operatorname{a}} 
\newcommand{\Fa}{F_{\bfa}} 


\newcommand{\lama}{\lam_{\bfa}}
\newcommand{\mua}{\mu_{\bfa}} 
\newcommand{\lamo}{\lam_{\bfa, t_0}}
\newcommand{\muo}{\mu_{\bfa, t_0}}

%\newcommand{\app}{\textrm{\textrm{app}}

\newcommand{\inte}{\operatorname{int}}
\newcommand{\exte}{\operatorname{ext}}

\newcommand{\uprd}{\overline{\vd_\tx{r}}}
\newcommand{\dold}{\underline{\vd_\tx{l}}}
\DeclareMathOperator*{\argmax}{arg\,max}

\newcommand{\oo}{\textrm{o}}
\newcommand{\ii}{\textrm{i}} 





\newcommand{\bfalp}{\boldsymbol{\alpha}}
\newcommand{\bfbt}{\boldsymbol{\beta}}
\newcommand{\bfgmm}{\boldsymbol{\gamma}}
\newcommand{\bfGmm}{\boldsymbol{\Gamma}}
\newcommand{\bfdlt}{\boldsymbol{\delta}}
\newcommand{\bfDlt}{\boldsymbol{\Delta}}
\newcommand{\bfeps}{\boldsymbol{\epsilon}}
\newcommand{\bfveps}{\boldsymbol{\varepsilon}}
%\newcommand{\bfpi}{\boldsymbol{\pi}}
%\newcommand{\bfmu}{\boldsymbol{\mu}}
%\newcommand{\bfnu}{\boldsymbol{\nu}}
\newcommand{\bfkpp}{\boldsymbol{\kappa}}
\newcommand{\bflmb}{\boldsymbol{\lambda}}
\newcommand{\bfLmb}{\boldsymbol{\Lambda}}
\newcommand{\bfvphi}{\boldsymbol{\varphi}}
\newcommand{\bfvpi}{\boldsymbol{\varpi}}
\newcommand{\bfvrho}{\boldsymbol{\varrho}}
\newcommand{\bfsgm}{\boldsymbol{\sigma}}
\newcommand{\bfSgm}{\boldsymbol{\Sigma}}
\newcommand{\bftht}{\boldsymbol{\theta}}
\newcommand{\bfTht}{\boldsymbol{\Theta}}
\newcommand{\bfvtht}{\boldsymbol{\vartheta}}
\newcommand{\bfomg}{\boldsymbol{\omega}}
\newcommand{\bfOmg}{\boldsymbol{\Omega}}
\newcommand{\bfups}{\boldsymbol{\upsilon}}
\newcommand{\bfzt}{\boldsymbol{\zeta}}
\newcommand{\bfZt}{\boldsymbol{\Zeta}}


\newcommand{\dlthf}{\dlt_{h}}							

\newcommand{\epsN}{\eps_{1}}							% small parameter for finding homotopy
\newcommand{\epshf}{\eps_{\ast}}							% small parameter for well-posedness near $Q$
\newcommand{\epslhf}{\eps_{\ast\ast}}						% small parameter for the size of $\ringA$ and $\ringPsi$ for the linearized harmonic map heat flow
\newcommand{\tDb}{\widetilde{\Db}}
\newcommand{\tgmfd}{\calN}							% Target manifold
\newcommand{\tg}{\tilde{\calN}}							% Target
\newcommand{\tcv}{\bfkpp}								% the Gauss curvature of the target manifold
\newcommand{\tsgn}{\bfsgm}								% the sign of the sectional curvature of the target manifold

%\newcommand{\omg}{\omega}
%\newcommand{\brk}[1]{\langle #1 \rangle}

\newcommand{\pfstep}[1]{\smallskip \noindent {\bf #1.}}		

				


%%---------------------- Local commands ------------------------------%
%\newcommand{\met}{G} 
%\newcommand{\phii}{f_i} 
%\newcommand{\phic}{f_c} 
%\newcommand{\phia}{F_a}
%
%
%\newcommand{\ULam}{\U{\Lam}} 
%\newcommand{\ulam}{\U{\lam}} 
%
%\newcommand{\app}{\operatorname{app}}
%\newcommand{\form}{\operatorname{form}}
%
%\newcommand{\aaa}{\operatorname{a}} 
%\newcommand{\Fa}{F_{\bfa}} 
%
%
%\newcommand{\lama}{\lam_{\bfa}}
%\newcommand{\mua}{\mu_{\bfa}} 
%\newcommand{\lamo}{\lam_{\bfa, t_0}}
%\newcommand{\muo}{\mu_{\bfa, t_0}}
%
%%\newcommand{\app}{\textrm{\textrm{app}}

\newcommand{\abl}{\partial}
\newcommand{\les}{\lesssim}
\newcommand{\dbar}{\abl_{\bar z}}
\newcommand{\dz}{\abl_z}
\newcommand{\ol}{\overline}
\def\loc{\mathrm{loc}}
 \def\Res{\mathrm{Res}}
 \def\PSL{\mathrm{PSL}}
\newcommand{\tom}{\tilde{\omega}}
\newcommand{\dom}{\mathrm{dom}}
\newcommand{\const}{\mathrm{const}}
\newcommand{\var}{\mathcal{V}}
\newcommand{\Xdil}{X_{\mathrm{dil}}}

\newcommand{\np}{_{\mathrm{np}}}
\newcommand{\spo}{_{\mathrm{sp}}}

\newcommand{\pz}{\partial_{z}} 
\newcommand{\pzb}{\partial_{\ba{z}}} 

\newcommand{\CP}{\mathbb{CP}}
\newcommand{\Ric}{\operatorname{Ric}}
\newcommand{\can}{\mathrm{can}}
\newcommand{\vol}{\operatorname{vol}}
\newcommand{\dvol}{\,\mathrm{dvol}}
\newcommand{\CD}{\operatorname{CD}}

\begin{document}
\title[Spectral obstructions to contracting transport maps on curved spaces]{Spectral obstructions to contracting transport maps on curved spaces}
\author{Shrey Aryan}


\begin{abstract}
Caffarelli's contraction theorem states that the Brenier optimal transport
map from the standard Gaussian measure to a more log-concave probability
measure is $1$-Lipschitz. Motivated by this theorem,
Milman~\cite{milman2018spectral} formulated several conjectures for the
round sphere and for weighted manifolds satisfying the curvature-dimension
condition $\CD(\rho,\infty)$. More recently, Beck and Jerison
\cite{beck2021friedland} raised related questions on the round hemisphere.
A contracting transport map as in these conjectures implies the
corresponding spectral comparison. In the spherical setting, this
comparison was also conjectured by Colding and Minicozzi
\cite{colding1998weyl} for compact manifolds with Ricci curvature lower
bounds. We construct counterexamples to the corresponding spectral
comparisons on spheres in dimensions $d\geq4$ and on weighted manifolds
satisfying the $\CD(1,\infty)$ condition in dimensions $d\geq4$, thereby
obtaining obstructions to contracting transport maps. In dimensions
$d\geq5$, the weighted counterexamples may additionally be chosen to
satisfy $\operatorname{Ric}_g\geq 0$ and $\nabla_g^2V \geq g$ separately. In dimension two, we use inverse mean curvature flow to construct a contracting transport map from the appropriately rescaled round sphere to every closed connected Riemannian surface satisfying the corresponding positive Ricci curvature lower bound. This implies the spectral comparison in dimension two. Together with the recent counterexample in dimension three constructed by Lin, Wang, and Xu
\cite{lin2026eigenvalues} and our counterexamples in dimensions $d\geq4$, this settles the spherical spectral comparison in every dimension $d\geq2$. Using the same method, we also construct a contracting transport from the uniform probability measure on a hemisphere onto the
normalized uniform measure on any geodesically convex subset of positive volume, thereby answering affirmatively the remaining case of a conjecture by Beck and Jerison \cite{beck2021friedland} following the work of Fathi, Fradelizi, Gozlan, and Zugmeyer \cite{fathi2023some}.
\end{abstract}


\keywords{Caffarelli contraction theorem; Lipschitz transport maps}
\subjclass[2010]{53C21 (primary); 58J50, 35P15, 53C20}

\thanks{}



\maketitle
\section{Introduction} 
In \cite{caffarelli2000monotonicity}, Caffarelli proved the following theorem:
\begin{thm}[cf. Theorem 11 \cite{caffarelli2000monotonicity}]\label{thm:caff}
Let $d\geq 1$, and let $\gamma^d$ denote the standard Gaussian measure on $\R^d$. If $\mu(\ud x)=\gamma^d(\ud x)$ and $\nu(\ud x) = e^{-W(x)} \gamma^d(\ud x)$ are two probability measures where $W$ is a convex function, then there exists a convex function $\phi : \mathbb{R}^d \to \mathbb{R}$ such that $T=\nabla \phi$ is $1$-Lipschitz, and transports $\mu$ to $\nu$, i.e. $T_{\#}\mu =\nu.$
\end{thm}
% The existence of the transport map $T=\nabla \phi$ follows from solving optimal transport for the quadratic distance cost, i.e.
% \begin{align}
%   T:= \operatorname{argmin}_{\nu = \tilde{T}_{\#} \mu} \int_{\R^d} |x-\tilde{T}(x)|^2 \mu(\ud x)
% \end{align}
% and Brenier's theorem \cite{brenier1991polar}. 
% however from this variation problem it is not clear why the transport map is $1$-Lipschitz. The key insight of Caffarelli was to use the Monge-Ampère equation satisfied by $T=\nabla \phi$,
% \begin{align}
% \det (D^2\phi) = e^{W} 
% \end{align}
% and use the log-concavity of the determinant, log-convexity of $W$ along with a clever maximum principle to deduce that a bound on $|D^2\phi|\leq 1$, implying that $|\nabla T|\leq 1$. In general, as long as the $\nu$ is more log-concave than $\mu$, one can obtain transport maps with dimension independent Lipschitz constants. 
The Gaussian space \((\R^d,|\cdot|,\gamma^d)\) can be regarded as a weighted
manifold with synthetic dimension \(\infty\) and constant Bakry--\'Emery
curvature \(1\). In particular, it satisfies the curvature-dimension $\operatorname{CD}(1,\infty)$ condition (cf. Theorem 0.12 in \cite{lott2009ricci}) and therefore one can view Theorem~\ref{thm:caff} as a statement that asks for contracting volume-preserving maps from the Gaussian space to a possibly more curved space (cf. p.~514 in \cite{villani2009optimal}). In \cite{milman2018spectral}, Milman formulated precise conjectures to generalize this observation on finite and infinite dimensional spaces:
\begin{con}[cf. Conjecture $3$ \cite{milman2018spectral}]\label{conj:sphere-spectral}
For any $(\bS^d,g,\vol_g)$ satisfying $\operatorname{Ric}_g\geq \rho g$ with $\rho>0$, we have
\begin{align}
        \lambda_k(\bS^d,g,\vol_g)
        \geq
        \lambda_k(\bS^d,g^\rho_{\can},\vol_{g^\rho_{\can}})
        \qquad
        \forall k\geq 1,
\end{align}
where $g^\rho_{\can}$ denotes the canonical metric on $\bS^d$ rescaled to have
\begin{align}
        \operatorname{Ric}_{g^\rho_{\can}}
        =
        \rho g^\rho_{\can},
\end{align}
where $g^\rho_{\can}$ denotes the rescaled canonical metric satisfying
\begin{align}
        \Ric_{g^\rho_{\can}}
        =
        \rho g^\rho_{\can}.
\end{align}
Equivalently,
\begin{align}
        g^\rho_{\can}
        =
        \frac{d-1}{\rho}g_{\can}.
\end{align}
\end{con}
For $\rho>0$, let
\begin{align}
        \gamma_\rho^d(\ud x)
        :=
        \left(
        \frac{\rho}{2\pi}
        \right)^{d/2}
        e^{-\rho|x|^2/2}\,\ud x.
\end{align}
\begin{con}[cf. Conjecture $4$ \cite{milman2018spectral}]\label{conj:sphere-transport}
For any $(\bS^d,g,\vol_g)$ satisfying $\operatorname{Ric}_g\geq \rho g$ with $\rho>0$, there exists a map
\begin{align}
        T:(\bS^d,g^\rho_{\can},\vol_{g^\rho_{\can}})
        \to
        (\bS^d,g,\vol_g),
\end{align}
pushing forward $\vol_{g^\rho_{\can}}$ onto $\vol_g$ up to a finite constant and contracting the corresponding metrics.
\end{con}
\begin{con}[cf. Conjecture $1^*$ \cite{milman2018spectral}]\label{conj:gaussian-spectral}
For any $(\bR^d,g,\mu)$ satisfying $\operatorname{CD}(\rho,\infty)$ with $\rho>0$, we have
\begin{align}
        \lambda_k(\bR^d,g,\mu)
        \geq
        \lambda_k(\bR^d,|\cdot|,\gamma_\rho^d)
        \qquad
        \forall k\geq 1.
\end{align}
\end{con}

\begin{con}[cf. Conjecture $2^*$ \cite{milman2018spectral}]\label{conj:gaussian-transport}
For any $(\bR^d,g,\mu)$ satisfying $\operatorname{CD}(\rho,\infty)$ with $\rho>0$, there exists a map
\begin{align}
        T:(\bR^d,|\cdot|,\gamma_\rho^d)
        \to
        (\bR^d,g,\mu),
\end{align}
pushing forward $\gamma_\rho^d$ onto $\mu$ up to a finite constant and contracting the corresponding metrics.
\end{con}
\begin{con}[cf. Conjectures $1$ and $2$ \cite{milman2018spectral}]\label{conj:gaussian-combined}
Conjectures~\ref{conj:gaussian-spectral} and
\ref{conj:gaussian-transport} hold when restricted to weighted manifolds diffeomorphic to $\R^d$ and satisfying
\begin{align}
        \Ric_g
        \geq
        0,
        \qquad
        \nabla_g^2V
        \geq
        \rho g.
\end{align}
Here
\begin{align}
        \mu
        =
        Z^{-1}e^{-V}\dvol_g,
\end{align}
where $V:\R^d\to\R$ is smooth and $Z$ is a normalizing constant.
\end{con}
After rescaling the Ricci curvature lower bound, Conjecture
\ref{conj:sphere-spectral} is a subcase of a long-standing
eigenvalue comparison problem stated by Colding and Minicozzi in \cite[p.~261]{colding1998weyl}. In particular, they asked whether every closed Riemannian manifold $(N^d,h)$ satisfying $\operatorname{Ric}_h\geq(d-1)h$ also satisfies
\begin{align}\label{eqn:CM-long-standing-problem}
        \lambda_k(N^d,h)
        \geq
        \lambda_k(\mathbb S^d,g_{\can})
        \qquad
        \forall k\geq1,
\end{align}
where the eigenvalues are counted with multiplicity. They also asked whether equality in \eqref{eqn:CM-long-standing-problem} implies that $(N^d,h)$ is isometric to the round sphere. Colding and Minicozzi observed that this problem is closely related, through the metric cone, to Yau's conjecture asking whether Euclidean space has the maximal dimension of polynomial growth harmonic functions among complete manifolds with nonnegative Ricci curvature, see \cite[Problem 48]{yau1992open}.

The main results of this work show that all the above conjectures are false in general when the dimension is large enough.
% \begin{mainthm}\label{thm:sphere-main}
% Let \(d\geq 4\). Then there exists a smooth Riemannian metric \(g\) on
% \(\mathbb S^d\), a number \(\rho>0\), and an integer \(K\geq 1\) such that
% \begin{align}
% \operatorname{Ric}_{g}\geq \rho g    
% \end{align}
% but
% \begin{align}\label{eqn:index-K}
% \lambda_K(\mathbb S^d,g) < \lambda_K(\mathbb S^d,g_{\mathrm{can}}^\rho).    
% \end{align}
% \end{mainthm}
\begin{mainthm}\label{thm:sphere-main}
Let $d\geq4$. Then there exist an integer $K\geq1$, a number
$\veps_d>0$, and a smooth family of Riemannian metrics
$\{\widehat g_\veps\}_{0\leq\veps<\veps_d}$ on $\bbS^d$ such that
$\widehat g_0=g_\can$ and
\begin{align}\label{eqn:sphere-local-convergence}
        \widehat g_\veps
        \rightarrow
        g_\can
        \qquad
        \text{in }C^\infty
        \text{ as }\veps\to0.
\end{align}
Moreover, for every $0<\veps<\veps_d$, we have
\begin{align}\label{eqn:sphere-local-Ricci}
        \Ric_{\widehat g_\veps}
        \geq
        (d-1)\widehat g_\veps,
\end{align}
but
\begin{align}\label{eqn:index-K}
        \lam_K(\bbS^d,\widehat g_\veps)
        <
        \lam_K(\bbS^d,g_\can).
\end{align}
\end{mainthm}
Thus, Theorem~\ref{thm:sphere-main} gives a negative answer to
Conjecture~\ref{conj:sphere-spectral} and to the eigenvalue comparison problem of Colding--Minicozzi in dimensions $d\geq4$. After the first version of this work appeared, Lin, Wang, and Xu \cite{lin2026eigenvalues} proved Conjecture~\ref{conj:sphere-spectral} when $d=2$ and constructed a counterexample on $\bbS^3$. Therefore, together with their work, Theorem~\ref{thm:sphere-main} settles the ordered eigenvalue comparison on spheres.
In contrast, we prove that Conjecture~\ref{conj:sphere-transport}
holds when $d=2$.
\begin{thm}\label{thm:sphere-transport-dimension-two}
Let $(N^2,g)$ be a closed connected Riemannian surface satisfying
\begin{align}\label{eq:sphere-transport-dimension-two-Ricci}
        \Ric_g
        \geq
        \rho g
\end{align}
for some $\rho>0$. Then there exists a continuous surjective map
\begin{align}
        T:
        (\bS^2,g^\rho_{\can})
        \to
        (N^2,g)
\end{align}
such that
\begin{align}\label{eq:sphere-transport-dimension-two-contraction}
        d_g(T(x),T(y))
        \leq
        d_{g^\rho_{\can}}(x,y)
        \qquad
        \text{for every }x,y\in\bS^2,
\end{align}
and
\begin{align}\label{eq:sphere-transport-dimension-two-pushforward}
        T_\#\dvol_{g^\rho_{\can}}
        =
        \frac{4\pi}{\rho\,\vol_g(N^2)}
        \dvol_g.
\end{align}
Here $g^\rho_{\can}=\rho^{-1}g_{\can}$.
\end{thm}
As a consequence, Conjecture~\ref{conj:sphere-spectral} holds when $d=2$. In the case in which $N^2$ is diffeomorphic to $\bS^2$, this result was also obtained at the same time independently by Han and Zhu \cite{han-zhu}, using a related inverse mean curvature flow argument.

Using the same idea, we also answer affirmatively the remaining question posed by Fathi, Fradelizi, Gozlan, and Zugmeyer \cite{fathi2023some}, arising from a question of Beck and Jerison \cite{beck2021friedland}.
\begin{thm}\label{thm:hemisphere-transport}
Let $d\geq2$, let $\bS^d_+$ be a closed hemisphere of the canonical round
sphere $(\bS^d,g_{\can})$, and let
\begin{align}\label{eq:hemisphere-transport-sigma-plus}
        \sigma_+
        :=
        \frac{1}{\vol_{g_{\can}}(\bS_+^d)}
        \mathbf{1}_{\bS^d_+}\dvol_{g_{\can}}
\end{align}
be the uniform probability measure on $\bS^d_+$. Let
$K\subset\bS^d_+$ be a closed geodesically convex subset satisfying $\sigma_+(K)>0$. Then there exists a continuous surjective map
\begin{align}
        T:\bS^d_+\to K
\end{align}
such that
\begin{align}\label{eq:hemisphere-transport-contraction}
        d_{g_{\can}}(T(x),T(y)) \leq d_{g_{\can}}(x,y)
        \quad
        \text{for every }x,y\in\bS^d_+,
\end{align}
and
\begin{align}\label{eq:hemisphere-transport-pushforward}
        T_\#\sigma_+=\frac{1}{\sigma_+(K)}\mathbf{1}_K\sigma_+.
\end{align}
\end{thm}
We note that the transport map constructed above need not coincide with
the optimal transport map for the quadratic distance cost. Indeed, take
\begin{align}
        \bS^2_+
        :=
        \left\{
        (x_1,x_2,x_3)\in\bS^2:
        x_3\geq0
        \right\},
        \quad 
        K
        :=
        \left\{
        (x_1,x_2,x_3)\in\bS^2:
        x_3\geq0,\ x_1\geq0
        \right\}.
\end{align}
In geodesic polar coordinates $(\theta,\varphi)$ centered at
$e_1=(1,0,0)$, the quadratic optimal transport map from $\sigma_+$ onto
the normalized uniform measure on $K$ is given $\sigma_+$-almost
everywhere by
\begin{align}
        T_{\mathrm{opt}}(\theta,\varphi)
        =
        (r(\theta),\varphi),
        \qquad
        \cos r(\theta)
        =
        \frac{1+\cos\theta}{2}.
\end{align}
Indeed, since $x\mapsto d_{g_{\can}}(e_1,x)$ is $1$-Lipschitz, every
transport plan has quadratic cost at least the one-dimensional quadratic
transport cost between the radial marginals, and $T_{\mathrm{opt}}$
attains this lower bound because $r$ is their monotone rearrangement and
$T_{\mathrm{opt}}$ moves each point along the same radial geodesic emanating from $e_1$ and ending at $-e_1.$ Moreover,
if $ e_\varphi :=\frac{1}{\sin\theta}\partial_\varphi,$ then the operator norm of the differential satisfies
\begin{align}
        \left|
        d(T_{\mathrm{opt}})_{(\theta,\varphi)}
        \right|_{\mathrm{op}}
        \geq
        \left|
        d(T_{\mathrm{opt}})_{(\theta,\varphi)}
        (e_\varphi)
        \right|_{g_{\can}}
        =
        \frac{\sin r(\theta)}{\sin\theta}
        \rightarrow
        \infty
\end{align}
as $\theta\to\pi$. Thus $T_{\mathrm{opt}}$ is not Lipschitz.

Next, in the $\CD(\rho,\infty)$ setting, we prove the following result.
\begin{mainthm}\label{thm:gaussian-main}
For every $d\geq4$, there exists a smooth complete weighted manifold $(\bR^d,g,\mu)$ satisfying the $\CD(1,\infty)$ condition and $\Ric_g\geq0$ such that
\begin{align}\label{eqn:index-gaussian-K}
\lambda_{d+2}(\bR^d,g,\mu)<2=
\lambda_{d+2}(\bR^d,|\cdot|,
\gamma^d).    
\end{align}
\end{mainthm}
Finally, we can strengthen the above result to show that
\begin{mainthm}\label{thm:gaussian-split-main}
For every $d\geq5$, there exist a smooth complete weighted manifold $(\bR^d,g,\mu)$ and an integer $K\geq1$ with the following properties. The measure $\mu$ has the form
\begin{align}
        \mu
        =
        Z^{-1}e^{-V}\dvol_g,
        \qquad
        Z
        :=
        \int_{\bR^d}e^{-V}\dvol_g
        <
        \infty,
\end{align}
where $V:\bR^d\to\bR$ is smooth. Moreover,
$(\bR^d,g,\mu)$ satisfies the $\CD(1,\infty)$ condition,
\begin{align}
        \Ric_g
        \geq
        0,
        \qquad
        \nabla_g^2V
        \geq
        g,
\end{align}
but
\begin{align}\label{eqn:index-gaussian-split-K}
        \lambda_K(\bR^d,g,\mu)
        <
        \lambda_K(\bR^d,|\cdot|,\gamma^d).
\end{align}
\end{mainthm}
Thus, Conjectures \ref{conj:sphere-spectral} and \ref{conj:gaussian-spectral} are false in general. Combining Theorems \ref{thm:sphere-main}, \ref{thm:gaussian-main} and \ref{thm:gaussian-split-main} with Theorem~\ref{thm:contraction}, we also deduce that Conjectures \ref{conj:sphere-transport}, \ref{conj:gaussian-transport} and \ref{conj:gaussian-combined} are false in general; see Corollaries \ref{cor:sphere}, \ref{cor:gaussian} and \ref{cor:gaussian-split-main}.

Although the conjectures above are false in general, they were supported by considerably more than the formal analogy with Caffarelli's theorem in Euclidean space. In the compact setting, Conjecture \ref{conj:sphere-spectral} is consistent with the Lichnerowicz theorem, the Bérard--Gallot heat-trace comparison, Weyl asymptotics, and a comparison up to a dimension-independent multiplicative constant for sufficiently large eigenvalues. Similarly, the existence of a contracting transport map in Conjecture \ref{conj:sphere-transport} is consistent with many fundamental results in comparison geometry, including the Bonnet--Myers diameter bound, Bishop--Gromov volume comparison, the Bakry--\'Emery log-Sobolev inequality, the sharp Sobolev inequality, and the L\'evy--Gromov isoperimetric inequality. We refer to \cite[Sections 1 and 2]{milman2018spectral} for a detailed discussion. 

There are also several positive generalizations, as well as new proofs, of Caffarelli's contraction theorem. These include results in infinite-dimensional spaces \cite{feyel2004monge,mikulincer2024brownian}, results for $1/d$-concave densities \cite{carlier2024optimal}, and results in curved settings \cite{fathi2024transportation,lopez2025bakry,ge2025generalization}. Kim and Milman \cite{kim2012generalization} constructed Lipschitz transports using the heat flow, while entropic optimal transport has led to further proofs and extensions \cite{chewi2023entropic,fathi2020proof}. In a related direction, Beck and Jerison \cite{beck2021friedland} asked whether there exists an injective $1$-Lipschitz transport from the uniform measure on the hemisphere to a
log-concave probability measure supported on a geodesically convex subset, and whether the corresponding $W_2$-optimal transport map is $1$-Lipschitz. Fathi, Fradelizi, Gozlan, and Zugmeyer \cite{fathi2023some} gave counterexamples to both assertions in general. They left open whether, for every geodesically convex subset $K$ of positive volume, there exists a $1$-Lipschitz transport from $\sigma_+$
onto the normalized uniform measure on $K$. This case is resolved by Theorem~\ref{thm:hemisphere-transport}.

% Besides Caffarelli's contraction theorem, our interest in contracting
% transport maps on curved spaces was also motivated by a conjecture of Yau, which states that if
% \(\Sigma^2\subset\mathbb S^{3}\) is a closed embedded minimal hypersurface,
% then its first nonzero Laplace--Beltrami eigenvalue is \(2\) (cf. \cite{yau1982seminar}). Using the convention $0=\lambda_1(\Sigma)\leq\lambda_2(\Sigma)\leq\cdots,$ one can show that this is equivalent to $\lam_2(\Sigma)\geq 2$ since the coordinate functions restricted to $\Sigma$ are eigenfunctions with eigenvalue $2$.
% For \(d=2\), a contracting measure-preserving map
% \(T:\mathbb S^2\to\Sigma\) would imply, by Theorem~\ref{thm:contraction}, that
% \begin{align}
%         \lambda_k(\Sigma)\geq\lambda_k(\mathbb S^2)
%         \qquad\forall k\geq1,
% \end{align}
% which in particular would imply the desired estimate when $k=2.$ However, this stronger conclusion is false in general. For instance, for the Clifford
% torus
% \begin{align}
%         \Sigma
%         =
%         \mathbb S^1\!\left(\frac1{\sqrt2}\right)
%         \times
%         \mathbb S^1\!\left(\frac1{\sqrt2}\right)
%         \subset\mathbb S^3,    
% \end{align}
% we have (cf. Remark 2 $\mathrm{(ii)}$ in \cite{muto1984homogeneous})
% \begin{align}
%  \lambda_5(\Sigma)=2<6=\lambda_5(\mathbb S^2).    
% \end{align}
% The fact that a product structure as above can give rise to high multiplicity of small eigenvalues is the key observation behind the construction of the counterexamples in this work.
% \section{Introduction} 
% In \cite{caffarelli2000monotonicity}, Caffarelli proved the following theorem:
% \begin{thm}[cf. Theorem 11 \cite{caffarelli2000monotonicity}]\label{thm:caff}
% Let $d\geq 1$, and let $\gamma_d$ denote the standard Gaussian measure on $\R^d$. If $\mu(\ud x)=\gamma_d(\ud x)$ and $\nu(\ud x) = e^{-W(x)} \gamma_d(\ud x)$ are two probability measures where $W$ is a convex function, then there exists a convex function $\phi : \mathbb{R}^d \to \mathbb{R}$ such that $T=\nabla \phi$ is $1$-Lipschitz, and transports $\mu$ to $\nu$, i.e. $T_{\#}\mu =\nu.$
% \end{thm}
% % The existence of the transport map $T=\nabla \phi$ follows from solving optimal transport for the quadratic distance cost, i.e.
% % \begin{align}
% %   T:= \operatorname{argmin}_{\nu = \tilde{T}_{\#} \mu} \int_{\R^d} |x-\tilde{T}(x)|^2 \mu(\ud x)
% % \end{align}
% % and Brenier's theorem \cite{brenier1991polar}. 
% % however from this variation problem it is not clear why the transport map is $1$-Lipschitz. The key insight of Caffarelli was to use the Monge-Ampère equation satisfied by $T=\nabla \phi$,
% % \begin{align}
% % \det (D^2\phi) = e^{W} 
% % \end{align}
% % and use the log-concavity of the determinant, log-convexity of $W$ along with a clever maximum principle to deduce that a bound on $|D^2\phi|\leq 1$, implying that $|\nabla T|\leq 1$. In general, as long as the $\nu$ is more log-concave than $\mu$, one can obtain transport maps with dimension independent Lipschitz constants. 
% The Gaussian space \((\R^d,|\cdot|,\gamma^d)\) can be regarded as a weighted
% manifold with synthetic dimension \(\infty\) and constant Bakry--\'Emery
% curvature \(1\). In particular, it satisfies the curvature-dimension $\operatorname{CD}(1,\infty)$ condition (cf. Theorem 0.12 in \cite{lott2009ricci}) and therefore one can view Theorem~\ref{thm:caff} as a statement that asks for contracting volume-preserving maps from the Gaussian space to a possibly more curved space (cf. p.~514 in \cite{villani2009optimal}). In \cite{milman2018spectral}, Milman formulated precise conjectures to generalize this observation on finite and infinite dimensional spaces:
% \begin{con}[cf. Conjecture $3$ \cite{milman2018spectral}]\label{conj:sphere-spectral}
% For any $(\bS^d,g,\vol_g)$ satisfying $\operatorname{Ric}_g\geq \rho g$ with $\rho>0$, we have
% \begin{align}
%         \lambda_k(\bS^d,g,\vol_g)
%         \geq
%         \lambda_k(\bS^d,g^\rho_{\can},\vol_{g^\rho_{\can}})
%         \qquad
%         \forall k\geq 1,
% \end{align}
% where $g^\rho_{\can}$ denotes the canonical metric on $\bS^d$ rescaled to have
% \begin{align}
%         \operatorname{Ric}_{g^\rho_{\can}}
%         =
%         \rho g^\rho_{\can}.
% \end{align}
% \end{con}
% \begin{con}[cf. Conjecture $4$ \cite{milman2018spectral}]\label{conj:sphere-transport}
% For any $(\bS^d,g,\vol_g)$ satisfying $\operatorname{Ric}_g\geq \rho g$ with $\rho>0$, there exists a map
% \begin{align}
%         T:(\bS^d,g^\rho_{\can},\vol_{g^\rho_{\can}})
%         \to
%         (\bS^d,g,\vol_g),
% \end{align}
% pushing forward $\vol_{g^\rho_{\can}}$ onto $\vol_g$ up to a finite constant and contracting the corresponding metrics.
% \end{con}
% \begin{con}[cf. Conjecture $1^*$ \cite{milman2018spectral}]\label{conj:gaussian-spectral}
% For any $(\bR^d,g,\mu)$ satisfying $\operatorname{CD}(\rho,\infty)$ with $\rho>0$, we have
% \begin{align}
%         \lambda_k(\bR^d,g,\mu)
%         \geq
%         \lambda_k(\bR^d,|\cdot|,\gamma_\rho^d)
%         \qquad
%         \forall k\geq 1,
% \end{align}
% where \(\gamma_\rho^d\) denotes the centered Gaussian probability measure
% \begin{align}
% \gamma_\rho^d(\ud x)
%         :=
%         c_{\rho}^d
%         e^{-\rho |x|^2/2}\,\ud x,    
% \end{align}
% where $c_\rho^d>0$ is a normalization constant. 

% \end{con}

% \begin{con}[cf. Conjecture $2^*$ \cite{milman2018spectral}]\label{conj:gaussian-transport}
% For any $(\bR^d,g,\mu)$ satisfying $\operatorname{CD}(\rho,\infty)$ with $\rho>0$, there exists a map
% \begin{align}
%         T:(\bR^d,|\cdot|,\gamma_\rho^d)
%         \to
%         (\bR^d,g,\mu),
% \end{align}
% pushing forward $\gamma_\rho^d$ onto $\mu$ up to a finite constant and contracting the corresponding metrics. Here \(\gamma_\rho^d\) denotes the centered Gaussian probability measure
% \begin{align}
% \gamma_\rho^d(\ud x)
%         :=
%         c_{\rho}^d
%         e^{-\rho |x|^2/2}\,\ud x,    
% \end{align}
% where $c_\rho^d>0$ is a normalization constant. 
% \end{con}
% \begin{con}[cf. Conjectures $1$ and $2$ \cite{milman2018spectral}]\label{conj:gaussian-combined}
% Conjectures $1^*$ and $2^*$ hold true when restricted to weighted manifolds satisfying $\operatorname{Ric}_g \geq 0$ and $\nabla_g^2 V \geq \rho g$ (and still diffeomorphic to $\mathbb{R}^d$). Here $\mu= Z^{-1}e^{-V}\dvol_g$, where $V:\R^d\to \R$ is a smooth function and $Z$ is a normalizing constant.
% \end{con}
% The main result of this work shows that all the above conjectures are false in general. In particular, 
% \begin{mainthm}\label{thm:sphere-main}
% Let \(d\geq 4\). Then there exists a smooth Riemannian metric \(g\) on
% \(\mathbb S^d\), a number \(\rho>0\), and an integer \(K\geq 1\) such that
% \begin{align}
% \operatorname{Ric}_{g}\geq \rho g    
% \end{align}
% but
% \begin{align}\label{eqn:index-K}
% \lambda_K(\mathbb S^d,g) < \lambda_K(\mathbb S^d,g_{\mathrm{can}}^\rho).    
% \end{align}
% \end{mainthm}
% Next, in the \(\operatorname{CD}(\rho,\infty)\) setting, we show
% \begin{mainthm}\label{thm:gaussian-main}
% For every $d\geq 4$, there is a smooth complete weighted manifold $(\bR^d,g,\mu)$ satisfying $\CD(1,\infty)$ such that $\Ric_g\geq 0$, 
% \begin{align}\label{eqn:index-gaussian-K}
% \lambda_{d+2}(\bR^d,g,\mu)<2=
% \lambda_{d+2}(\bR^d,|\cdot|,
% \gamma^d).    
% \end{align}
% \end{mainthm}
% Finally, we can strengthen the above result to show that
% \begin{mainthm}\label{thm:gaussian-split-main}
% For every $d\geq 5$, there exists a smooth complete weighted manifold
% $(\bR^d,g,\mu)$ satisfying $\CD(1,\infty)$ condition, with $\mu=Z^{-1}e^{-V}\dvol_g$, where $Z=\int_{\bR^d}e^{-V}\dvol_g<\infty,$ $V:\bR^d\to\bR$ is a smooth function, and an integer $K\geq1$ such that 
% \begin{align}
%         \Ric_g\geq0,
%         \qquad
%         \nabla_g^2V\geq g,
% \end{align}
% but
% \begin{align}\label{eqn:index-gaussian-split-K}
%         \lambda_K(\bR^d,g,\mu)
%         <
%         \lambda_K(\bR^d,|\cdot|,\gamma^d).
% \end{align}
% \end{mainthm}
% Thus, Conjectures \ref{conj:sphere-spectral} and \ref{conj:gaussian-spectral} are false in general. Combining Theorems \ref{thm:sphere-main}, \ref{thm:gaussian-main} and \ref{thm:gaussian-split-main} with Theorem~\ref{thm:contraction}, we also deduce that Conjectures \ref{conj:sphere-spectral}, \ref{conj:gaussian-transport} and \ref{conj:gaussian-combined} are false in general; see Corollaries \ref{cor:sphere}, \ref{cor:gaussian} and \ref{cor:gaussian-split-main}.

% Although the conjectures above are false in general, they were supported by considerably more than the formal analogy with Caffarelli's theorem in Euclidean space. In the compact setting, Conjecture \ref{conj:sphere-spectral} is consistent with the Lichnerowicz theorem, the Bérard--Gallot heat-trace comparison, Weyl asymptotics, and a comparison up to a dimension-independent multiplicative constant for sufficiently large eigenvalues. Similarly, the existence of a contracting transport map in Conjecture \ref{conj:sphere-transport} is consistent with many fundamental results in comparison geometry, including the Bonnet--Myers diameter bound, Bishop--Gromov volume comparison, the Bakry--\'Emery log-Sobolev inequality, the sharp Sobolev inequality, and the L\'evy--Gromov isoperimetric inequality. We refer to \cite[Sections 1 and 2]{milman2018spectral} for a detailed discussion. 

% There are also several positive generalizations, as well as new proofs, of Caffarelli's contraction theorem. These include results in infinite-dimensional spaces \cite{feyel2004monge,mikulincer2024brownian}, results for $1/d$-concave densities \cite{carlier2024optimal}, and results in curved settings \cite{fathi2024transportation,lopez2025bakry,ge2025generalization}. Kim and Milman \cite{kim2012generalization} constructed Lipschitz transports using the heat flow, while entropic optimal transport has led to further proofs and extensions \cite{chewi2023entropic,fathi2020proof}. The question of generalizing Theorem~\ref{thm:caff} was also raised by Jerison--Beck \cite{beck2021friedland} in the setting of measures supported on the hemisphere. Later, Fathi \cite{fathi2023some} constructed counterexamples to the existence of contracting Brenier maps from the uniform measure on the hemisphere to log-concave measures on the hemisphere. 


% Besides Caffarelli's contraction theorem, our interest in contracting
% transport maps on curved spaces was also motivated by a conjecture of Yau, which states that if
% \(\Sigma^2\subset\mathbb S^{3}\) is a closed embedded minimal hypersurface,
% then its first nonzero Laplace--Beltrami eigenvalue is \(2\) (cf. \cite{yau1982seminar}). Using the convention $0=\lambda_1(\Sigma)\leq\lambda_2(\Sigma)\leq\cdots,$ one can show that this is equivalent to $\lam_2(\Sigma)\geq 2$ since the coordinate functions restricted to $\Sigma$ are eigenfunctions with eigenvalue $2$.
% For \(n=2\), a contracting measure-preserving map
% \(T:\mathbb S^2\to\Sigma\) would imply, by Theorem~\ref{thm:contraction}, that
% \begin{align}
%         \lambda_k(\Sigma)\geq\lambda_k(\mathbb S^2)
%         \qquad\forall k\geq1,
% \end{align}
% which in particular would imply the desired estimate when $k=2.$ However, this stronger conclusion is false in general. For instance, for the Clifford
% torus
% \begin{align}
%         \Sigma
%         =
%         \mathbb S^1\!\left(\frac1{\sqrt2}\right)
%         \times
%         \mathbb S^1\!\left(\frac1{\sqrt2}\right)
%         \subset\mathbb S^3,    
% \end{align}
% we have (cf. Remark 2 $\mathrm{(ii)}$ in \cite{muto1984homogeneous})
% \begin{align}
%  \lambda_5(\Sigma)=2<6=\lambda_5(\mathbb S^2).    
% \end{align}
% The fact that a product structure as above can give rise to high multiplicity of small eigenvalues is the key observation behind the construction of the counterexamples in this work.
% \subsection{Acknowledgements}
% The author would like to thank his advisor, Tobias Colding, for his support over the years and for many insightful conversations. The author is also grateful to David Jerison for introducing him to the Caffarelli contraction theorem and for several illuminating discussions about transport maps. Finally, the author would also like to thank Max Fathi, Promit Ghosal, Dan Mikulincer, and Yair Shenfeld for many helpful discussions. This work was supported in part by a Simons Dissertation Fellowship and NSF DMS Grant 2405393.
\section{Preliminaries}
A weighted manifold is a triple \((M^d,g,\mu)\), where \((M,g)\) is a complete
smooth Riemannian manifold,
\[
        \ud\mu=e^{-W}\dvol_g,
\]
and \(W\) is smooth. The weighted Laplacian is defined as:
\begin{align}
\Delta_{g,\mu}f= \Delta_g f-\langle\nabla_g W,\nabla_g f\rangle_g.    
\end{align}
We use the nonnegative operator $-\Delta_{g,\mu}$ and for $k\geq1$, define the $k$-th variational eigenvalue by
\begin{align}
        \lambda_k(M,g,\mu)
        :=
        \inf_{\substack{
        E\subset W^{1,2}(M,g,\mu)\\
        \dim E=k
        }}
        \sup_{0\neq f\in E}
        \frac{
        \int_M|\nabla f|_g^2\,\ud\mu
        }{
        \int_Mf^2\,\ud\mu
        }.
\end{align}
When the corresponding weighted Laplacian has compact resolvent, finite volume $\mu(M)<\infty$ and $M$ is connected then these variational eigenvalues are the usual eigenvalues, counted with multiplicity. We can write these eigenvalues in non-decreasing order as
\[
0=\lambda_1(M,g,\mu)\leq \lambda_2(M,g,
\mu)\leq\cdots,
\]
including multiplicity. The Bakry--\'Emery curvature tensor is
\[
\Ric_{g,
\mu}:=\Ric_g+\nabla_g^2W.
\]
From Theorem 0.12 \cite{lott2009ricci}, a weighted manifold $(M^d,g,\mu)$ is said to satisfy curvature-dimension condition $\CD(\rho,
\infty)$ for $\rho \in \R$ if 
\begin{align}
\Ric_g+\nabla_g^2W\geq \rho g.    
\end{align}
Let $\gamma^d(\ud x)=(2\pi)^{-d/2}e^{-|x|^2/2}\,\ud x$ denote the standard Gaussian measure on $\bR^d$, with the weighted Laplacian being the Ornstein--Uhlenbeck operator whose first few eigenvalues satisfy (cf. Section 2.1 in \cite{milman2018spectral})
\begin{align}\label{eqn:gaussian-spectrum}
\lambda_1(\R^d,|\cdot|,\gamma^d)=0,
\qquad
\lambda_2(\R^d,|\cdot|,\gamma^d)=\cdots=\lambda_{d+1}(\R^d,|\cdot|,\gamma^d)=1,
\qquad
\lambda_{d+2}(\R^d,|\cdot|,\gamma^d)=2.    
\end{align}
Note that $(\R^d,|\cdot|,\gamma^d)$ is $\CD(1,\infty).$  Next, we recall the following theorem due to Milman \cite{milman2018spectral} that relates spectral comparison between weighted manifolds to the Lipschitz transport maps.
\begin{thm}[Contraction Principle]\label{thm:contraction}
Let
\[
    T : (M_1, g_1, \mu_1) \to (M_2, g_2, \mu_2)
\]
denote an $L$-Lipschitz map between two complete weighted manifolds
pushing forward $\mu_1$ onto $\mu_2$ up to a finite constant. Then
\[
    \lambda_k(M_2, g_2, \mu_2)
    \geq
    \frac{1}{L^2}\lambda_k(M_1, g_1, \mu_1)
    \qquad \forall k \geq 1.
\]
\end{thm}
In particular, a $1$-Lipschitz transport map from a model space to a target space implies that the spectrum of the target space dominates that of the source. 



\section{Proof of Theorem~\ref{thm:sphere-main}}\label{sec:sphere-thm-main}
Let $d\geq 4$. We denote the round sphere $\bbS^d=\{(z,y)\in \mathbb C\times \mathbb R^{d-1}: |z|^2+|y|^2=1\}$ and will write \(z=x_1+\mathrm i x_2\). Consider the map
\begin{align}
    F:(0,\pi/2)\times \bbS^1\times \bbS^{d-2} \to \mathbb{S}^d \subset \mathbb{C}\times \bbR^{d-1}
\end{align}
such that 
\begin{align}
    F(t,\theta,\omega) = (e^{i\theta}\sin t, (\cos t) \omega).
\end{align}
Then $F$ maps into the sphere $\mathbb{S}^d$ and is a diffeomorphism onto the open set $U:= \{(z,y)\in\mathbb S^d\subset\mathbb C\times\mathbb R^{d-1}:z\neq0,\ y\neq0\}.$ Let $g_0$ denote the induced Euclidean metric on $\mathbb{S}^d$, then in local coordinates on $U$, we can write the metric $g_0$ using the pull-back metric $F^*g_0$ as follows:
\begin{align}\label{eqn:expression-for-g_can-Sn}
    g_0  &= |\ud z|^2 + |\ud y|^2\\
    &= |e^{i\theta} (\cos t)  \ud t+ ie^{i\theta} (\sin t)  \ud \theta|^2 + |-(\sin t) \omega \ud t + (\cos t) \ud \omega|^2\\
    &= \cos^2 t \ud t^2 + \sin^2 t \ud \theta^2 + \sin^2 t \ud t^2 + (\cos^2 t) g_{\mathbb{S}^{d-2}}\\
    &= \ud t^2+\sin^2t\,\ud \theta^2+\cos^2t\,g_{\mathbb S^{d-2}},
\end{align}
where in the above expression we used the orthogonality of $e^{i\theta}$ and $ie^{i\theta}$ with respect to the real Euclidean inner product in $\C\simeq \R^2$ and the fact $|\omega|^2=1$ which after differentiating implies that $\la \omega, \ud \omega\ra=0$.

The idea is to perturb this metric slightly in the \(\bbS^1\)-direction. As
expected, this lowers the Ricci curvature lower bound from \(d-1\) to some value $\rho$, but it will also drop the eigenvalues of the perturbed manifold below that of the sphere with the round metric and $\Ric=\rho g.$ Note that, $g_0=g_{\can}$ where $g_{\can}$ denotes the canonical round metric on the unit sphere $\bbS^d$ and $\Ric_{g_0}=(d-1)g_0.$

Before proceeding with the proof of Theorem~\ref{thm:sphere-main}, we record a
useful lemma (cf. Lemma 2.3 in \cite{karaca2018gradient}) which we will use to
compute the Ricci curvature of warped product metrics.
\begin{lem}[Ricci tensor of a multiple warped product]\label{lem:multiply-warped-ricci}
Let \((B,g_B)\) and \((F_i,g_i)\), \(1\leq i\leq m\), be Riemannian
manifolds, and let $b_i:B\to (0,\infty)$ be smooth positive functions. Let $M=B\times_{b_1}F_1\times\cdots\times_{b_m}F_m$ be the multiply warped product with metric
\[
        g
        =
        g_B+\sum_{i=1}^m b_i^2 g_i.
\]
Write $s_i=\dim F_i.$ Let \(X,Y\) be vector fields tangent to the base \(B\), and let \(V_i,W_i\)
be vector fields tangent to the fiber \(F_i\), all lifted to \(M\). Then the
Ricci tensor of \(g\) is given as follows:
\begin{enumerate}
    \item $\operatorname{Ric}_g(X,Y)
        =
        \operatorname{Ric}_B(X,Y)
        -
        \sum_{i=1}^m
        \frac{s_i}{b_i}\operatorname{Hess}_B b_i(X,Y).$
    \item $\operatorname{Ric}_g(X,V_i)=0.$
    \item if
\(i\neq j\), then
$\operatorname{Ric}_g(V_i,V_j)=0.$
    \item for each $1\leq i\leq m$ we have
    \begin{align}
      \operatorname{Ric}_g(V_i,W_i)
        &=
        \operatorname{Ric}_{F_i}(V_i,W_i)
        \\
        &\quad
        -
        g(V_i,W_i)
        \left(
        \frac{\Delta_B b_i}{b_i}
        +
        (s_i-1)\frac{|\nabla_B b_i|^2}{b_i^2}
        +
        \sum_{\substack{k=1\\ k\neq i}}^m
        s_k
        \frac{\langle\nabla_B b_i,\nabla_B b_k\rangle_B}{b_i b_k}
        \right).  
    \end{align}
    
\end{enumerate}
Here all gradients, Hessians, Laplacians, and inner products appearing on
the right-hand side are computed with respect to the base metric \(g_B\).
\end{lem}

\begin{proof}[Proof of Theorem~\ref{thm:sphere-main}]
Let \(d\geq 4\). Our proof proceeds in four main steps. 
\paragraph{$\mathrm{(i)}$ \textit{Construction of the metric}} For \(0<\varepsilon\ll 1\), define a metric on the coordinate chart $U$ as 
\begin{equation}\label{eq:geps-coordinate}
g_\varepsilon=
        \ud t^2 +
        \sin^2t\bigl(1+\varepsilon\sin^4t\bigr)\ud \theta^2 + \cos^2t\,g_{\mathbb S^{d-2}}.
\end{equation}
This is a smooth, well-defined metric on $\bbS^d$ as can be seen by rewriting 
\begin{align}\label{eq:geps-global}
g_\varepsilon= g_0+\varepsilon |z|^2\alpha\otimes\alpha, 
\end{align}
where $\alpha=x_1\,dx_2-x_2\,dx_1=\sin^2 t \ud \theta$ and we used the fact that $|z|^2=\sin^2 t.$
\paragraph{$\mathrm{(ii)}$ \textit{Ricci curvature lower bound}} We now estimate its Ricci tensor. Set 
\begin{align}
a_{\veps}(t):= (\sin t)\sqrt{1+\veps \sin^4 t},\quad c(t):=\cos t 
\end{align}
Then, working in a local orthonormal frame with basis vectors $e_t=\partial_t$, $e_\theta=\frac{1}{a_\veps(t)} \partial_\theta$ and $e_j=\frac{1}{\cos t} v_j$ where $\{v_j\}_{j=1}^{d-2}$ is a local $g_{\bbS^{d-2}}$-orthonormal frame on the $\bbS^{d-2}$-factor, we can denote 
\begin{align}
R_t=\Ric_{g_\veps}(e_t,e_t),\quad R_\theta=\Ric_{g_{\veps}}(e_\theta,e_\theta),\quad R_{S}=\Ric_{g_\veps}(e_j,e_j) \text{ for }1\leq j\leq d-2.
\end{align}
Note that the last term is independent of any $1\leq j\leq d-2$ since the Ricci curvature of $\bbS^{d-2}$ is a constant multiple of the $g_{\bbS^{d-2}}.$ Using Lemma \ref{lem:multiply-warped-ricci} we can compute each term as follows. First,
\begin{align}
        -\frac{a_\veps''}{a_\veps}
        =
        1-\veps\left(\frac{\sin^2 t(10-12\sin^2 t)+2\veps \sin^6 t (3-4\sin^2 t)}{(1+\veps \sin^4 t)^2}\right),
\end{align}
and
\begin{align}\label{eq:tan-a-prime-over-a}
        \tan t\frac{a_\veps'}{a_\veps}
        =
        1+
        \frac{2\veps \sin^4t}{1+\veps\sin^4t}.
\end{align}
We shall also use the elementary estimate
\begin{align}\label{eq:ricci-elementary-bound}
        f(q):=
        \frac{q (10-12q )+2\veps q^3  (3-4q )}{(1+\veps q^2)^2}
        \leq \frac{25}{12}
\end{align}
for $q\in [0,1]$ and $\veps\geq0$. 
% Indeed,
% \begin{align}
%         \frac{25}{12}(1+\veps q^2)^2
%         &-
%         \left(q(10-12q)+2\veps q^3(3-4q)\right)\notag\\
%         &=
%         12\left(q-\frac{5}{12}\right)^2
%         +
%         \veps q^2\left(8q^2-6q+\frac{25}{6}\right)
%         +
%         \frac{25}{12}\veps^2q^4
%         \geq0,
% \end{align}
% because $8q^2-6q+\frac{25}{6}>0$ for all $q\in[0,1]$.
Thus
\begin{align}
 R_t &= \operatorname{Ric}_{g_{\veps}}(\partial_t, \partial_t) \\
 &= \Ric_{B}(\partial_t,\partial_t) - \sum_{i=1}^2 \frac{s_i}{b_i}\operatorname{Hess}_{B}b_i(\partial_t,\partial_t)\\
 &= 0 - \frac{a_\veps''}{a_\veps}-\frac{(d-2)c''}{c} \\
 &= 1-\veps f(\sin^2 t) + (d-2) \\
 &\geq (d-1) -\frac{25\veps}{12}, 
\end{align}
where $B=(0,\pi/2)$, $s_1=\operatorname{dim} \bbS^1=1$, $s_2=\dim \bbS^{d-2}=(d-2)$, $b_1=a_\veps$, and $b_2=c$. Similarly, applying Lemma \ref{lem:multiply-warped-ricci}  to the
\(\bbS^1\)-factor with \(B=(0,\pi/2)\), \(F_1=\bbS^1\), \(F_2=\bbS^{d-2}\),
\(b_1=a_\veps\), \(b_2=c\), \(s_1=\dim\bbS^1=1\), and
\(s_2=\dim\bbS^{d-2}=d-2\) and using \eqref{eq:tan-a-prime-over-a}, we have
\begin{align}
    R_\theta
    &=
    \Ric_{g_\veps}(e_\theta,e_\theta) =
    \frac{1}{a_\veps^2}
    \Ric_{g_\veps}(\partial_\theta,\partial_\theta) \notag\\
    &=
    \frac{1}{a_\veps^2}
    \left[
    \Ric_{\bbS^1}(\partial_\theta,\partial_\theta)
    -
    g_\veps(\partial_\theta,\partial_\theta)
    \left(
    \frac{a_\veps''}{a_\veps}
    +
    (d-2)\frac{a_\veps'c'}{a_\veps c}
    \right)
    \right] \notag\\
    &=
    -\frac{a_\veps''}{a_\veps}
    -(d-2)\frac{a_\veps'c'}{a_\veps c} \notag\\
    &=
    -\frac{a_\veps''}{a_\veps}
    +(d-2)\tan t \frac{a_\veps'}{a_\veps} \notag\\
    &=
    (d-2)
    +
    \frac{2(d-2)\veps \sin^4 t}{1+\veps \sin^4 t}
    +
    1
    -
    \veps f(\sin^2 t) \notag\\
    &\geq d-1-\frac{25\veps}{12}.
\end{align}
Lastly, applying Lemma
\ref{lem:multiply-warped-ricci} to the \(\bbS^{d-2}\)-factor gives
\begin{align}
    R_S
    &=
    \Ric_{g_\veps}(e_j,e_j) =
    \frac{1}{c^2}\Ric_{g_\veps}(v_j,v_j) \notag\\
    &=
    \frac{1}{c^2}
    \left[
    \Ric_{\bbS^{d-2}}(v_j,v_j)
    -
    g_\veps(v_j,v_j)
    \left(
    \frac{c''}{c}
    +
    (d-3)\frac{(c')^2}{c^2}
    +
    \frac{a_\veps'c'}{a_\veps c}
    \right)
    \right] \notag\\
    &=
    \frac{d-3}{c^2}
    -\frac{c''}{c}
    -(d-3)\frac{(c')^2}{c^2}
    -\frac{a_\veps'c'}{a_\veps c} \notag\\
    &=
    d-2+\tan t\frac{a_\veps'}{a_\veps} \notag\\
    &=
    d-2+1+\frac{2\veps\sin^4t}{1+\veps\sin^4t}
    \geq d-1.
\end{align}
% Similarly, using \eqref{eq:tan-a-prime-over-a}, we have 
% \begin{align}
%     R_\theta &=  -\frac{a_\veps''}{a_\veps} + (d-2)\tan t \frac{a_\veps'}{a_\veps}\\
%     &= (d-2) + \frac{2(d-2)\veps \sin^4 t}{1+\veps \sin^4 t} + 1  \\
%     &\quad  -\veps f(\sin^2 t)\\
%     &\geq d-1 -\frac{25\veps}{12}.
% \end{align}
% For $v_j$ with $|v_j|_{g_{\bbS^{d-2}}}=1$, the corresponding
% $g_\veps$-unit vector is $e_j=c^{-1}v_j$. Applying Lemma
% \ref{lem:multiply-warped-ricci} to the $\bbS^{d-2}$-factor gives
% \begin{align}
%     R_S
%     &=\operatorname{Ric}_{g_\veps}(e_j,e_j)\notag\\
%     &=
%     \frac{d-3}{c^2}
%     -\frac{c''}{c}
%     -(d-3)\frac{(c')^2}{c^2}
%     -\frac{a_\veps'c'}{a_\veps c}\notag\\
%     &=
%     d-2+\tan t\frac{a_\veps'}{a_\veps}\notag\\
%     &=
%     d-2+1+\frac{2\veps\sin^4t}{1+\veps\sin^4t}
%     \geq d-1.
% \end{align}
Combining the above lower bounds on $R_t$, $R_\theta$ and $R_S$ implies that
\begin{align}
        \operatorname{Ric}_{g_\varepsilon}
        \geq
        \left(
        d-1-\frac{25 \varepsilon}{12}
        \right)g_\varepsilon
        \qquad\text{on }U.
\end{align}
By continuity the above inequality holds on all of $\bbS^d$, i.e.
\begin{equation}\label{eq:ricci-lower}
        \operatorname{Ric}_{g_\varepsilon}
        \geq
        \left(
        d-1-\frac{25 \varepsilon}{12}
        \right)g_\varepsilon.
\end{equation}
Set $\rho_\veps:=d-1-\frac{25 \varepsilon}{12}.$ For $d\geq 2$, $\rho_\veps>0$ for $0<\veps\ll 1$.  

\paragraph{$\mathrm{(iii)}$ \textit{Construction of test functions}} We next show the desired spectral inequality. To achieve this goal, we will construct suitable test functions.  Consider the Rayleigh quotient,
\EQ{
\calR_\varepsilon(f)
    = \frac{\int_{\mathbb S^d}
        |\nabla f|_{g_\varepsilon}^2\,\dvol_{g_\varepsilon}}
        {\int_{\mathbb S^d}
        |f|^2\,\dvol_{g_\varepsilon}}.
}
The simplest test functions to try would be to take functions of the form $|z|^k$, however to take advantage of the $\bbS^1$ factor in the metric, we will consider for $k\geq 1$, the function
\begin{align}
     f_k=\operatorname{Re}(z^k) = q^{k/2}\cos(k\theta),
        \text{ where } q=\sin^2t.   
\end{align}
The $\cos(k\theta)$ term in $f_k$ will turn out to be advantageous. Observe that
\begin{align}\label{eqn:denom-lb}
\int_{\mathbb S^d} |f_k|^2\,\dvol_{g_\varepsilon} &= \int_{\bbS^d} q^{k} \cos^2(k\theta) \sqrt{1+\veps q^2} \dvol_{g_0}\\
&= \frac{1}{2}\int_{\bbS^d} q^{k}  \sqrt{1+\veps q^2} \dvol_{g_0}
\end{align}
where we used the fact that $\dvol_{g_\varepsilon}=\sqrt{1+\varepsilon q^2}\,\dvol_{g_0},$ and the fact that $\int_0^{2\pi} \cos^2(k\theta) =\pi=\frac{1}{2}\int_{0}^{2\pi}\ud\theta$. Then denote $Q_j=\int_{\mathbb{S}^d} q^j \dvol_{g_0}.$ Using $\sqrt{1+x}\geq 1+\frac{x}{2}-\frac{x^2}{8}$ and the above expression, we get 
\begin{align}
  \int_{\mathbb S^d} |f_k|^2\,\dvol_{g_\varepsilon} &\geq \frac{Q_k}{2} +\frac{\veps Q_{k+2}}{4}-\frac{\veps^2 Q_{k+4}}{16}.
\end{align}
Next, we estimate the numerator. Since, $q=\sin^2 t$ we have
\[
        \partial_t f_k
        =
        kq^{(k-1)/2}\sqrt{1-q}\,\cos(k\theta),
        \qquad
        \partial_\theta f_k
        =
        -kq^{k/2}\sin(k\theta),
\]
we have 
\[
|\nabla f_k|_{g_\varepsilon}^2
=
k^2q^{k-1}(1-q)\cos^2(k\theta)
+
\frac{k^2q^{k-1}}{1+\varepsilon q^2}\sin^2(k\theta).
\]
Thus,
\begin{align}\label{eqn:numerator-ub}
\int_{\mathbb S^d} |\nabla f_k|_{g_\veps}^2\,\dvol_{g_\varepsilon} &= \int_{\bbS^d} k^2 q^{k-1}(1-q)\cos^2(k\theta)\sqrt{1+\veps q^2} \dvol_{g_0}  \\
&\quad + \int_{\bbS^d}\frac{k^2 q^{k-1}}{1+\veps q^2} \sin^2 (k\theta) \sqrt{1+\veps q^2} \dvol_{g_0}\\
&= \frac{k^2}{2}\int_{\bbS^d} q^{k-1}(1-q)\sqrt{1+\veps q^2} \dvol_{g_0}  + \frac{k^2}{2}\int_{\bbS^d}\frac{q^{k-1}}{\sqrt{1+\veps q^2}} \dvol_{g_0}\\
&\leq \frac{k^2}{2} \left(2Q_{k-1}-Q_k -\frac{\veps}{2}Q_{k+2}+\frac{3\veps^2}{8}Q_{k+3}\right)\\
&\leq \frac{1}{2}k(k+d-1)Q_k - \frac{\veps k^2}{4} Q_{k+2} + \frac{3\veps^2 k^2}{16} Q_{k+3}
\end{align}
where we used the elementary inequality for $q\in [0,1]$,
\begin{align}
    (1-q)\sqrt{1+\veps q^2} + \frac{1}{\sqrt{1+\veps q^2}}\leq 2-q-\frac{\veps q^3}{2} + \frac{3\veps^2q^4}{8}
\end{align}
and the fact that the sum of the squares of the first two coordinate functions on \(\bbS^d\),
namely \(q=|z|^2\), follows a beta distribution $\operatorname{Beta}(1,\frac{d-1}{2})$ (cf. Corollary 1.1 \cite{frankl1990some}). This implies, in particular, that (cf. Equation~(2.2) in \cite{szablowski2021moment})
\begin{align}
    \frac{Q_k}{Q_{k-1}} =\frac{2k}{2k+d-1},\quad 2Q_{k-1}-Q_k  =\left(\frac{k+d-1}{k}\right) Q_k.
\end{align}
Define
\begin{align}
        \beta_{d,k}
        :=
        \frac{k^2+k(k+d-1)}{2k(k+d-1)}
        \frac{Q_{k+2}}{Q_k}.
\end{align}
Since
\begin{align}
        \frac{Q_{k+2}}{Q_k}
        =
        \frac{2(k+1)}{2(k+1)+d-1}\cdot
        \frac{2(k+2)}{2(k+2)+d-1}
        \to 1, \text{ as }k\to \infty
\end{align}
we see that $\lim_{k\to\infty}\beta_{d,k}=1$. Therefore, since $d\geq4$, we may choose $k\geq k_0$ for some $k_0\gg 1$ such that
\begin{align}
        \eta_{d,k}:=\beta_{d,k}-\frac{25}{12(d-1)}>0.
\end{align}
We fix this value of $k$ for the remainder of the proof. In particular, the integer $K$ defined below depends only on $d$ and not on $\veps$. For some fixed \(k\geq k_0\), define
\begin{align}
        \Phi_{d,k}(\veps)
        :=
        \frac{
        \frac{1}{2}k(k+d-1)Q_k-\frac{\veps k^2}{4}Q_{k+2}
        +\frac{3\veps^2k^2}{16}Q_{k+3}}
        {\frac{Q_k}{2}+\frac{\veps Q_{k+2}}{4}-\frac{\veps^2Q_{k+4}}{16}}.
\end{align}
For $\veps>0$ small enough, the denominator is positive and
\eqref{eqn:denom-lb} and \eqref{eqn:numerator-ub} imply
\begin{align}
        \cR_\veps(f_k)\leq \Phi_{d,k}(\veps).
\end{align}
Moreover, $\Phi_{d,k}$ is smooth near $\veps=0$ and
\begin{align}
        \Phi_{d,k}(0)&=k(k+d-1),\\
        \Phi_{d,k}'(0)&=
        -\frac{k^2+k(k+d-1)}{2}\frac{Q_{k+2}}{Q_k}
        =-k(k+d-1)\beta_{d,k}.
\end{align}
Thus there is a constant $C_{d,k}<\infty$ such that, for all sufficiently small
$\veps>0$,
\begin{align}
        \cR_\veps(f_k)
        \leq
        k(k+d-1)(1-\beta_{d,k}\veps)+C_{d,k}\veps^2.
\end{align}
Choosing $\veps>0$ further so that
\begin{align}
        C_{d,k}\veps^2
        <
        \frac12 k(k+d-1)\eta_{d,k}\veps,
\end{align}
we get
\begin{align}\label{eqn:f_k-ub}
        \mathcal{R}_\veps(f_k)
        &<
        k(k+d-1)\left(1-\frac{25\veps}{12(d-1)}\right)\\
        &=
        \frac{\rho_\veps}{d-1}k(k+d-1),
\end{align}
where we recall that $\rho_\veps:=d-1-\frac{25\veps}{12}.$ 
Define the following real vector spaces
\begin{align}
  E_{k-1}:=\bigoplus_{\ell=0}^{k-1}H_{\ell},\quad F_k:= E_{k-1} \oplus \operatorname{span}\{f_k\}  
\end{align}
where $H_\ell$ is the space of spherical harmonics of degree $0\leq \ell\leq k-1$. Note that $f_k$ is orthogonal to $H_\ell$ for any $0\leq \ell\leq k-1$ since, after complexifying $H_\ell$, any $h_{\ell}\in H_\ell$ can be expressed as a linear combination of $z^\alp \bar{z}^\beta y^\gamma$ for $\alp,\beta\in \bN_0$ and $\gamma \in \bN_0^{d-1}$ and some constants $c_{\alp\beta\gamma}\in \bC$ as follows:
\begin{align}
    \int_{\mathbb{S}^d} h_{\ell} f_k \dvol_{g_\veps} &= \sum_{\alpha+\beta+|\gamma|=\ell}c_{\alpha\beta\gamma}\int_{\mathbb{S}^d} z^\alp \bar{z}^\beta y^\gamma  \Re(z^k) \dvol_{g_\veps} \\
    &= \sum_{\alpha+\beta+|\gamma|=\ell}\frac{c_{\alpha\beta\gamma}}{2}\int_{\mathbb{S}^d} z^\alp \bar{z}^\beta y^\gamma (z^k+\bar{z}^k)   \dvol_{g_\veps}\\
    &= \frac12
    \sum_{\alpha+\beta+|\gamma|=\ell}
    c_{\alpha\beta\gamma}
    \left[
    \int_{\mathbb S^d}
    z^{\alpha+k}\bar z^\beta y^\gamma
    \,d\operatorname{vol}_{g_\veps}
    +
    \int_{\mathbb S^d}
    z^\alpha\bar z^{\beta+k}y^\gamma
    \,d\operatorname{vol}_{g_\veps}
    \right].
\end{align}
Both of the above integrals vanish. This can be seen by making the change of variables
\((z,y)\mapsto(e^{i\tau}z,y)\) which preserves the volume form $\dvol_{g_{\veps}}$
\begin{align}\label{eqn:change-of-variables}
    \int_{\mathbb S^d}
    z^{\alpha+k}\bar z^\beta y^\gamma
    \,\dvol_{g_\veps}
    &=
    e^{i(\alpha+k-\beta)\tau}
    \int_{\mathbb S^d}
    z^{\alpha+k}\bar z^\beta y^\gamma
    \,\dvol_{g_\veps}.
\end{align}
Since $\alpha+\beta+|\gamma|=\ell\leq k-1,$  we get $\alpha+k-\beta\neq 0.$ Choosing
\(\tau\) so that \(e^{i(\alpha+k-\beta)\tau}\neq 1\), we obtain
\[
        \int_{\mathbb S^d}
        z^{\alpha+k}\bar z^\beta y^\gamma
        \,\dvol_{g_\veps}
        =
        0.
\]
Similarly,
\begin{align}
    \int_{\mathbb S^d}
    z^\alpha\bar z^{\beta+k}y^\gamma
    \,\dvol_{g_\veps}
    &= 0.
\end{align}
Thus, for any $0\leq \ell \leq k-1$, we have
\begin{align}
\int_{\mathbb{S}^d} h_{\ell} f_k \dvol_{g_\veps} = 0.
\end{align}
We now claim that 
\begin{align}\label{claim:F_k}
    \sup_{f\neq 0,f\in F_k}\mathcal{R}_\veps(f)<\frac{\rho_\veps}{d-1}k(k+d-1).
\end{align}
To this end, take any \(f\in F_k\) and write
\(f=h_{k-1}+a f_k\), where \(a\in\R\) and \(h_{k-1}\in E_{k-1}\). Then
\begin{align}
  \cR_\veps(f) &= \mathcal{R}_\veps(h_{k-1} + a f_k)\\
  &=  \frac{\int_{\mathbb S^d}
        |\nabla h_{k-1}|_{g_\varepsilon}^2\,\dvol_{g_\varepsilon}+a^2\int_{\mathbb S^d}
        |\nabla f_{k}|_{g_\varepsilon}^2\,\dvol_{g_\varepsilon}}
        {\int_{\mathbb S^d}
        |h_{k-1}|^2\,\dvol_{g_\varepsilon}+a^2\int_{\mathbb S^d}
        |f_{k}|^2\,\dvol_{g_\varepsilon}}\\
        &\leq \max \{\cR_\veps(h_{k-1}),\cR_\veps(f_k)\},
\end{align}
where we used the fact that for any $h_{k-1}\in E_{k-1}$,
\begin{align}
    \int_{\mathbb{S}^d} \la \nabla h_{k-1}, \nabla f_k \ra_{g_\veps} \dvol_{g_\veps} &= \int_{\mathbb{S}^d} (\partial_t h_{k-1})(\partial_t f_k) \dvol_{g_\veps}\\
    &\quad + \frac{1}{q(1+\veps q^2)}(\partial_\theta h_{k-1})(\partial_\theta f_k) + \frac{1}{\cos^2 t}\la \nabla_\om h_{k-1}, \nabla_\om f_k\ra  \dvol_{g_\veps}\\
    &= 0.
\end{align}
Indeed, $f_k$ is independent of $\omega$, so the last term vanishes. For a monomial
$z^\alpha\bar z^\beta y^\gamma$ of degree $\ell\leq k-1$, the $\partial_t$- and
$\partial_\theta$-terms above are finite sums of terms with $\theta$-weights
$e^{i(\alpha-\beta+k)\theta}$ or $e^{i(\alpha-\beta-k)\theta}$, multiplied by
functions of $t$ and $\omega$ only. Since $\alpha+\beta+|\gamma|=\ell\leq k-1$, both
$\alpha-\beta+k$ and $\alpha-\beta-k$ are nonzero. Their $\theta$-integrals vanish,
and the claim follows by linearity. We already have an upper bound for $\cR_\veps(f_k)$ so we focus on $\cR_\veps(h_{k-1}).$ To this end, first observe that for the Rayleigh quotient for the canonical round metric on the sphere $\bbS^d$, we have 
\begin{align}
    \cR_0(h)\leq (k-1)(k+d-2)
\end{align}
for any $h\in E_{k-1}$ since 
\begin{align}
    -\Delta_{g_0} h_{\ell} = \ell(\ell + d-1) h_{\ell}
\end{align}
for any $h_{\ell} \in H_\ell$ and $0\leq \ell\leq k-1.$ Furthermore, for any function $h\in E_{k-1}$ and for sufficiently small $0<\veps\ll 1$ we have
\begin{align}\label{eqn:h-ub}
    \cR_\veps(h) &= \frac{\int_{\bbS^d} |\nabla h|_{g_\veps}^2\dvol_{g_\veps}}{\int_{\bbS^d} h^2 \dvol_{g_\veps}}  \\
    &\leq \sqrt{1+\veps} \frac{\int_{\bbS^d} |\nabla h|_{g_0}^2\dvol_{g_0}}{\int_{\bbS^d} h^2 \dvol_{g_0}}  \\
    &= \sqrt{1+\veps}\cR_0(h)\\
    &\leq \left(1+\frac{\veps}{2}\right) (k-1)(k+d-2)\\
    &< \frac{\rho_\veps}{d-1}k(k+d-1),
\end{align}
where the last inequality follows by choosing $\veps<\veps_0$ where $\veps_0:=\frac{2k+d-2}{\frac{1}{2}(k-1)(k+d-2)+\frac{25}{12(d-1)}k(k+d-1)}$. Therefore, by \eqref{eqn:f_k-ub} and \eqref{eqn:h-ub} we have 
\begin{align}
    \sup_{f\neq 0,f\in F_k}\mathcal{R}_\veps(f)<\frac{\rho_\veps}{d-1}k(k+d-1).
\end{align}
\paragraph{$\mathrm{(iv)}$ \textit{Rescaling and convergence to the round metric}} Define
\begin{align}
        K
        :=
        \operatorname{dim}E_{k-1}+1,
        \quad
        c_\veps
        :=
        \frac{\rho_\veps}{d-1}
        =
        1-\frac{25\veps}{12(d-1)},\quad 
        \widehat g_\veps
        :=
        c_\veps g_\veps.
\end{align}
Since $E_{k-1}$ is the direct sum of the eigenspaces of the round sphere of degrees $0,\ldots,k-1$, the integer $K$ is the first index in the degree-$k$ eigenvalue cluster. Hence
\begin{align}\label{eqn:sphere-round-K}
        \lam_K(\bbS^d,g_0)
        =
        k(k+d-1).
\end{align}
It follows from the min--max characterization and the estimate above on
$F_k$ that
\begin{align}\label{eqn:sphere-unscaled-local-gap}
        \lam_K(\bbS^d,g_\veps)
        \leq
        \sup_{0\neq f\in F_k}\cR_\veps(f)
        <
        \frac{\rho_\veps}{d-1}k(k+d-1)
        =
        c_\veps\lam_K(\bbS^d,g_0).
\end{align}

Choose $\veps_d>0$ sufficiently small so that all the preceding estimates
hold and $c_\veps>0$ for every $0\leq\veps<\veps_d$. Then the metrics
$\widehat g_\veps$ are well-defined. Since $c_\veps$ is constant on
$\bbS^d$, the Levi--Civita connections of $g_\veps$ and
$\widehat g_\veps$ agree. Therefore, using
\eqref{eq:ricci-lower}, we obtain
\begin{align}
        \Ric_{\widehat g_\veps}
        =
        \Ric_{g_\veps}\geq
        \rho_\veps g_\veps=
        (d-1)\widehat g_\veps.
\end{align}
This proves \eqref{eqn:sphere-local-Ricci}. Moreover, for every constant $c>0$, we have $ \cR_{cg}(f)=c^{-1}\cR_g(f).$ 
Consequently, the min--max characterization gives
$\lam_j(\bbS^d,cg)=c^{-1}\lam_j(\bbS^d,g)$ for every $j\geq1$. Thus,
using \eqref{eqn:sphere-unscaled-local-gap}, we obtain
\begin{align}
        \lam_K(\bbS^d,\widehat g_\veps)
        =
        c_\veps^{-1}\lam_K(\bbS^d,g_\veps)<
        \lam_K(\bbS^d,g_0).
\end{align}
This proves \eqref{eqn:index-K}. Finally, by \eqref{eq:geps-global}, we have
\begin{align}
        \widehat g_\veps-g_0
        &=
        \left(
        1-\frac{25\veps}{12(d-1)}
        \right)
        \left(
        g_0+\veps |z|^2\alpha\otimes\alpha
        \right)
        -g_0\\
        &=
        \veps
        \left(
        |z|^2\alpha\otimes\alpha
        -\frac{25}{12(d-1)}g_0
        \right)
        -
        \frac{25\veps^2}{12(d-1)}
        |z|^2\alpha\otimes\alpha.
\end{align}
Since $|z|^2\alpha\otimes\alpha$ is a smooth tensor on the compact sphere,
for every integer $m\geq0$ there exists a constant $C_{d,m}<\infty$ such
that $\left\|
        \widehat g_\veps-g_0
        \right\|_{C^m(g_0)}
        \leq
        C_{d,m}\veps.$
Consequently,
\begin{align}
        \widehat g_\veps
        \rightarrow
        g_0
        =
        g_\can
         \text{ in }C^\infty
        \text{ as }\veps\to0.
\end{align}
This proves \eqref{eqn:sphere-local-convergence} and completes the proof.
\end{proof}
% \begin{cor}\label{cor:sphere}
% Let $(\bbS^d,g)$ be the metric constructed in Theorem~\ref{thm:sphere-main}. Then there does not exist a map $T:(\bbS^d,g_\can,\dvol_{g_\can})\to (\bbS^d,g,\dvol_g)$ which is measure preserving and contracting.  
% \end{cor}
\begin{cor}\label{cor:sphere}
Let $\{\widehat g_\veps\}_{0\leq\veps<\veps_d}$ be the family of metrics
constructed in Theorem~\ref{thm:sphere-main}. Then, for every
$0<\veps<\veps_d$, there does not exist a $1$-Lipschitz map
\[
        T:
        (\bbS^d,g_\can,\vol_{g_\can})
        \to
        (\bbS^d,\widehat g_\veps,\vol_{\widehat g_\veps})
\]
pushing forward $\vol_{g_\can}$ onto $\vol_{\widehat g_\veps}$ up to a
finite constant.
\end{cor}

\begin{proof}
If such a map existed, then Theorem~\ref{thm:contraction} would imply that for every $k\geq1$ we have
\begin{align}
        \lam_k(\bbS^d,\widehat g_\veps)
        \geq
        \lam_k(\bbS^d,g_\can).
\end{align}
This contradicts \eqref{eqn:index-K}.
\end{proof}
Thus, Theorem~\ref{thm:sphere-main} and Corollary \ref{cor:sphere} imply that Conjectures 3 and 4 in \cite{milman2018spectral} are not true in general.
\section{Proof of Theorem~\ref{thm:sphere-transport-dimension-two}}
\begin{proof}[Proof of Theorem~\ref{thm:sphere-transport-dimension-two}]
By rescaling the metric it suffices to prove the theorem in the case when $\rho=1$. Our proof proceeds in four main steps. We first prove the theorem under the stronger assumption $\Ric_g>g$ and then remove the strict inequality by an approximation argument.
\paragraph{$\mathrm{(i)}$ \textit{Construction of the initial embedding and the flow}}
Assume that $\Ric_g>g$. Since $g$ is two-dimensional, we have
\begin{align}
        \Ric_g
        =
        K_g g,
        \qquad
        \operatorname{Scal}_g
        =
        2K_g,
\end{align}
where $K_g$ denotes the Gauss curvature. Thus $K_g>1$. By the Gauss--Bonnet theorem and the classification of closed
surfaces, $N^2$ is diffeomorphic to either $\bS^2$ or $\mathbb{RP}^2$. Let
\begin{align}
        \pi:(\widetilde N,\widetilde g)\to(N,g)
\end{align}
be the identity map in the first case and the orientable double covering in the
second case, where $\widetilde g:=\pi^*g$. Denote the degree of $\pi$ by
$m\in\{1,2\}$. Then $\pi$ is a Riemannian covering, $\widetilde N$ is
diffeomorphic to $\bS^2$, $K_{\widetilde g}=K_g\circ\pi>1$, $\Ric_{\widetilde g}>\widetilde g$, and
\begin{align}\label{eq:sphere-transport-dimension-two-covering-identities}
        \vol_{\widetilde g}(\widetilde N)
        =
        m\vol_g(N),
        \qquad
        \pi_\#\dvol_{\widetilde g}
        =
        m\dvol_g,
\end{align}
while
\begin{align}\label{eq:sphere-transport-dimension-two-covering-distance}
        d_g(\pi(p),\pi(q))
        \leq
        d_{\widetilde g}(p,q)
        \qquad
        \text{for every }p,q\in\widetilde N.
\end{align}
After identifying $\widetilde N$ with $\bS^2$, \cite[Theorem~1.3]{lu2020weyl} gives a $C^2$ isometric embedding $ X_0:(\widetilde N,\widetilde g) \to  (\bS^3,\bar g_{\can}).$ Since $\widetilde N$ is orientable, $X_0$ admits a global unit normal.
Let $A_0$ denote the corresponding shape operator. The Gauss equation,
which is valid for the $C^2$ embedding, gives
\begin{align}
        \det A_0
        =
        K_{\widetilde g}-1
        >
        0.
\end{align}
Thus $X_0$ is an elliptic isometric embedding. By the regularity theorem
of Dubrovin \cite{dubrovin1966regularity}, if the intrinsic
and ambient metrics are of class $C^k$, with $k>3$, then a $C^2$
elliptic isometric immersion is of class $C^{k-1,\alpha}$ for every
$0<\alpha<1$. Since $\widetilde g$ and $\bar g_{\can}$ are smooth, it
follows that $X_0$ is smooth.

Moreover, $A_0$ is definite at every point. Since $\widetilde N$ is
connected, its sign is constant. After reversing the global unit normal
if necessary, we may assume that the second fundamental form $h_0$ is
positive definite. Set $ \widetilde M_0:=X_0(\widetilde N),$ and let $M_0$ denote its polar hypersurface. By \cite[Section~4]{gerhardt2015curvature}, $M_0$ is smooth and strictly
convex. Define
\begin{align}
        F(\lambda_1,\lambda_2)
        &:=
        \frac{2}{
        \lambda_1^{-1}+\lambda_2^{-1}
        },
        \\
        \widetilde F(\kappa_1,\kappa_2)
        &:=
        \frac{\kappa_1+\kappa_2}{2}.
\end{align}
These curvature functions satisfy the hypotheses of
\cite[Theorem~1.1]{gerhardt2015curvature}. Applying that theorem to the
polar pair $(M_0,\widetilde M_0)$ gives the expanding flow
\begin{align}\label{eq:sphere-transport-dimension-two-flow}
        \partial_tX_t
        =
        \frac{2}{H_t}\nu_t,
        \qquad
        X(0,\cdot)
        =
        X_0,
        \qquad
        0\leq t<T_*.
\end{align}
The flow exists smoothly up to a finite time $T_*$, remains strictly
convex, and converges smoothly in geodesic graph coordinates to a round
equator as $t\to T_*$. Here, $\nu_t$ denotes the outward unit normal.

\paragraph{$\mathrm{(ii)}$ \textit{Evolution of the induced metric and its volume form}}
Let $g_t:=X_t^*\bar g_{\can}$ denote the induced metric on the fixed surface
$\widetilde N$. For $V,W\in T\widetilde N$, set
$h_t(V,W):=\la\bar\nabla_{dX_t(V)}\nu_t,dX_t(W)\ra_{\bar g_{\can}}$,
where $dX_t$ denotes the differential of $X_t$. Since strict convexity is
preserved by the flow, $h_t$ is positive definite and
$H_t=\operatorname{tr}_{g_t}h_t$. If $X_i:=\partial_iX_t$ and
$\bar\nabla_i:=\bar\nabla_{X_i}$, then
\begin{align}\label{eq:sphere-transport-dimension-two-metric-evolution}
        \partial_t(g_t)_{ij}
        =
        \la\bar\nabla_i(2H_t^{-1}\nu_t),X_j\ra_{\bar g_{\can}}
        +
        \la X_i,\bar\nabla_j(2H_t^{-1}\nu_t)\ra_{\bar g_{\can}}
        =
        \frac{4}{H_t}(h_t)_{ij}.
\end{align}
Since $h_t$ is positive definite and $H_t>0$, we have
$g_t\geq g_0=\widetilde g$ for every $0\leq t<T_*$. Taking the trace in
\eqref{eq:sphere-transport-dimension-two-metric-evolution} gives
$\operatorname{tr}_{g_t}(\partial_tg_t)=4$. Using
$\partial_t\dvol_{g_t}=\frac12\operatorname{tr}_{g_t}(\partial_tg_t)
\dvol_{g_t}$, we obtain $\partial_t\dvol_{g_t}=2\dvol_{g_t}$, and hence
\begin{align}\label{eq:sphere-transport-dimension-two-volume-scaling}
        \dvol_{g_t}
        =
        e^{2t}\dvol_{\widetilde g}.
\end{align}
Thus the induced metric is monotone increasing, while its volume form is
multiplied by the spatially constant factor $e^{2t}$.

\paragraph{$\mathrm{(iii)}$ \textit{Construction of the transport map}}
By the convergence statement in \cite[Theorem~1.1]{gerhardt2015curvature}, we
may choose coordinates $(r,\omega)$ on an open set containing the limiting
equator such that
$\bar g_{\can}=\ud r^2+\sin^2r\,g_{\can}$ and the limiting equator is given by
$r=\pi/2$. For $t$ sufficiently close to $T_*$, there are smooth functions
$u_t:\bS^2\to(0,\pi)$ and parametrizations
$Y_t(\omega):=(u_t(\omega),\omega)$ of $X_t(\widetilde N)$ such that
$u_t\to\pi/2$ in $C^\infty(\bS^2)$ as $t\to T_*$. Consequently,
\begin{align}\label{eq:sphere-transport-dimension-two-source-convergence}
        \widehat g_t
        &:=
        Y_t^*\bar g_{\can}
        =
        \ud u_t\otimes\ud u_t+\sin^2(u_t)g_{\can},
        \qquad
        \widehat g_t
        \rightarrow
        g_{\can}
        \quad
        \text{in }C^\infty(\bS^2)
        \text{ as }t\to T_*.
\end{align}

Define $\widetilde T_t:=X_t^{-1}\circ Y_t:\bS^2\to\widetilde N$. Since
$X_t\circ\widetilde T_t=Y_t$, we have
$\widetilde T_t^*g_t=\widehat g_t$. Combining this identity with
$g_t\geq\widetilde g$ gives
$\widetilde T_t^*\widetilde g\leq\widehat g_t$. In particular,
\begin{align}\label{eq:sphere-transport-dimension-two-Tt-distance}
        d_{\widetilde g}(\widetilde T_t(x),\widetilde T_t(y))
        \leq
        d_{\widehat g_t}(x,y)
        \qquad
        \text{for every }x,y\in\bS^2.
\end{align}
Moreover, \eqref{eq:sphere-transport-dimension-two-volume-scaling} gives
$\dvol_{\widehat g_t}=\widetilde T_t^*\dvol_{g_t}
=e^{2t}\widetilde T_t^*\dvol_{\widetilde g}$. Therefore,
\begin{align}\label{eq:sphere-transport-dimension-two-Tt-pushforward}
        (\widetilde T_t)_\#\dvol_{\widehat g_t}
        =
        e^{2t}\dvol_{\widetilde g}.
\end{align}
Since $\widetilde T_t$ is a diffeomorphism and
$\widetilde T_t^*g_t=\widehat g_t$, we have
\begin{align}
        \vol_{g_t}(\widetilde N)
        =
        \vol_{\widehat g_t}(\bS^2)
        \rightarrow
        4\pi
\end{align}
as $t\to T_*$. Combining this with
\eqref{eq:sphere-transport-dimension-two-volume-scaling}, we obtain
\begin{align}\label{eq:sphere-transport-dimension-two-terminal-factor}
        e^{2T_*}
        =
        \frac{4\pi}{\vol_{\widetilde g}(\widetilde N)}.
\end{align}

By \eqref{eq:sphere-transport-dimension-two-source-convergence}, the metrics
$\widehat g_t$ are uniformly bounded above by a fixed multiple of $g_{\can}$
for $t$ sufficiently close to $T_*$. Thus
\eqref{eq:sphere-transport-dimension-two-Tt-distance} implies that the maps
$\widetilde T_t$ are equicontinuous. By the Arzel\`a--Ascoli theorem, after
passing to a sequence $t_j\to T_*$, there exists a continuous map
$\widetilde T:\bS^2\to\widetilde N$ such that
$\widetilde T_{t_j}\to\widetilde T$ uniformly. Since
$d_{\widehat g_{t_j}}\to d_{g_{\can}}$ uniformly on
$\bS^2\times\bS^2$, we may pass to the limit in
\eqref{eq:sphere-transport-dimension-two-Tt-distance} to obtain
\begin{align}\label{eq:sphere-transport-dimension-two-cover-limit-contraction}
        d_{\widetilde g}(\widetilde T(x),\widetilde T(y))
        \leq
        d_{g_{\can}}(x,y)
        \qquad
        \text{for every }x,y\in\bS^2.
\end{align}
We next verify that $\widetilde T$ is a transport map. For every
$f\in C^0(\widetilde N)$,
\eqref{eq:sphere-transport-dimension-two-Tt-pushforward} gives
\begin{align}
        \int_{\bS^2}f(\widetilde T_{t_j}(x))
        \dvol_{\widehat g_{t_j}}(x)
        =
        e^{2t_j}\int_{\widetilde N}f\dvol_{\widetilde g}.
\end{align}
Using the uniform convergence of $\widetilde T_{t_j}$,
\eqref{eq:sphere-transport-dimension-two-source-convergence}, and
\eqref{eq:sphere-transport-dimension-two-terminal-factor}, we may let
$j\to\infty$ to obtain
\begin{align}\label{eq:sphere-transport-dimension-two-cover-pushforward}
        \widetilde T_\#\dvol_{g_{\can}}
        =
        \frac{4\pi}{\vol_{\widetilde g}(\widetilde N)}
        \dvol_{\widetilde g}.
\end{align}
To see that $\widetilde T$ is surjective, suppose that
$y\notin\widetilde T(\bS^2)$. Since $\widetilde T(\bS^2)$ is compact, there
exists a nonempty open set containing $y$ and disjoint from
$\widetilde T(\bS^2)$. The left-hand side of the preceding identity assigns
zero measure to this set, whereas the right-hand side assigns positive measure,
which is a contradiction.

Set $T:=\pi\circ\widetilde T:\bS^2\to N$. By
\eqref{eq:sphere-transport-dimension-two-covering-distance} and
\eqref{eq:sphere-transport-dimension-two-cover-limit-contraction},
\begin{align}\label{eq:sphere-transport-dimension-two-limit-contraction}
        d_g(T(x),T(y))
        \leq
        d_{g_{\can}}(x,y)
        \qquad
        \text{for every }x,y\in\bS^2.
\end{align}
Furthermore, using
\eqref{eq:sphere-transport-dimension-two-covering-identities} and
\eqref{eq:sphere-transport-dimension-two-cover-pushforward}, we obtain
\begin{align}\label{eq:sphere-transport-dimension-two-normalized-pushforward}
        T_\#\dvol_{g_{\can}}
        &=
        \pi_\#\left(
        \widetilde T_\#\dvol_{g_{\can}}
        \right)
        =
        \frac{4\pi}{\vol_{\widetilde g}(\widetilde N)}
        \pi_\#\dvol_{\widetilde g}
        =
        \frac{4\pi}{\vol_g(N)}\dvol_g.
\end{align}
Since both $\widetilde T$ and $\pi$ are surjective, $T$ is surjective. This
proves the theorem under the strict inequality $\Ric_g>g$.

\paragraph{$\mathrm{(iv)}$ \textit{Removal of the strict inequality}}
Suppose now that $\Ric_g\geq g$. For $0<\veps<1$, set
$g_\veps:=(1-\veps)g$. Since a constant rescaling does not change the Ricci
tensor, we have $\Ric_{g_\veps}=\Ric_g\geq g=(1-\veps)^{-1}g_\veps>g_\veps$. By the
previous step, there exists a continuous surjective map
$T_\veps:(\bS^2,g_{\can})\to(N,g_\veps)$ such that
\begin{align}
        d_g(T_\veps(x),T_\veps(y))
        &\leq
        \frac{1}{\sqrt{1-\veps}}d_{g_{\can}}(x,y),
        \notag\\
        (T_\veps)_\#\dvol_{g_{\can}}
        &=
        \frac{4\pi}{\vol_{g_\veps}(N)}\dvol_{g_\veps}
        =
        \frac{4\pi}{\vol_g(N)}\dvol_g.
\end{align}
Choose a sequence $\veps_j\to0$ with $0<\veps_j\leq1/2$. The maps $T_{\veps_j}$ are equicontinuous. After passing to a uniformly convergent subsequence, the same
limiting argument as above gives a continuous map
$T:(\bS^2,g_{\can})\to(N,g)$ satisfying
\eqref{eq:sphere-transport-dimension-two-limit-contraction} and
\eqref{eq:sphere-transport-dimension-two-normalized-pushforward}. The same
argument using the pushforward identity shows that $T$ is surjective. This completes the proof when $\rho=1$. The general case follows by applying the preceding result to $g_\rho:=\rho g$ and using
\begin{align}
        d_{g_\rho}
        =
        \sqrt{\rho}\,d_g,
        \quad
        \dvol_{g_\rho}
        =
        \rho\,\dvol_g,
        \quad
        g^\rho_{\can}
        =
        \rho^{-1}g_{\can}.
\end{align}
\end{proof}



\section{Proof of Theorem~\ref{thm:hemisphere-transport}}
The proof uses the expanding inverse mean curvature flow together with a smooth approximation of the convex set $K$ by the upper halves of symmetric strictly convex hypersurfaces.
\begin{lem}\label{lem:hemisphere-smooth-symmetric-transport}
Let $d\geq2$, identify $\bS^{d+1}$ with the unit sphere in
$\bR^{d+1}\times\bR$, and let
\begin{align}
        R(x,z):=(x,-z).
\end{align}
Let
\begin{align}
        \Sigma^d
        \subset
        (\bS^{d+1},\overline g_{\can})
\end{align}
be a smooth closed strictly convex hypersurface contained in an open
hemisphere and invariant under $R$. Assume that
$\Sigma\cap\{z>0\}$ is connected and set $\Sigma^+ := \Sigma\cap\{z\geq0\}.$ Let $g_\Sigma$ denote the metric induced on $\Sigma$ by
$\overline g_{\can}$. Then there exist an $R$-invariant round equator $E\subset\bS^{d+1}$, a closed hemisphere $E^+\subset E$, and a continuous surjective map
\begin{align}
        P:E^+\to\Sigma^+.
\end{align}
Denoting by $g_E$ the induced round metric on $E$, we have
\begin{align}\label{eq:hemisphere-smooth-half-contraction}
        d_{g_\Sigma}(P(x),P(y))
        \leq
        d_{g_E}(x,y)
        \qquad
        \text{for every }x,y\in E^+,
\end{align}
and
\begin{align}\label{eq:hemisphere-smooth-half-pushforward}
        P_\#\left(
        \frac{1}{\vol_{g_E}(E^+)}
        \dvol_{g_E}|_{E^+}
        \right)
        =
        \frac{1}{\vol_{g_\Sigma}(\Sigma^+)}
        \dvol_{g_\Sigma}|_{\Sigma^+}.
\end{align}
\end{lem}
\begin{proof}
Let $\widetilde\Sigma$ denote the polar hypersurface of $\Sigma$. By \cite[Section~4]{gerhardt2015curvature}, $\widetilde\Sigma$ is smooth and strictly convex. Since $R$ is an
isometry and $R(\Sigma)=\Sigma$, the polar hypersurface $\widetilde\Sigma$ is also $R$-invariant. Define
\begin{align}
        F(\lambda_1,\ldots,\lambda_d)
        &:=
        \frac{d}{
        \sum_{i=1}^d\lambda_i^{-1}
        },
        \\
        \widetilde F(\kappa_1,\ldots,\kappa_d)
        &:=
        \frac{1}{d}
        \sum_{i=1}^d\kappa_i.
\end{align}
These curvature functions satisfy the hypotheses of
\cite[Theorem~1.1]{gerhardt2015curvature}. Applying that theorem to the
polar pair $(\widetilde\Sigma,\Sigma)$, we obtain a smooth family of strictly convex
embeddings satisfying
\begin{align}\label{eq:hemisphere-smooth-expanding-flow}
        \partial_tX_t
        &=
        \frac{d}{H_t}\nu_t,
        \qquad
        X(0,p)
        =
        p
        \quad
        \text{for every }p\in\Sigma,
        \qquad
        0\leq t<T_*.
\end{align}
The flow converges smoothly in geodesic graph coordinates to a round
equator $E$ as $t\to T_*$. Set $r:=R|_\Sigma$. By uniqueness of the flow,
\begin{align}\label{eq:hemisphere-smooth-flow-equivariance}
        X_t\circ r
        =
        R\circ X_t.
\end{align}
Let $x_0\in\bS^{d+1}$ denote the extinction point of the polar
contracting flow. Since the polar flow is $R$-invariant and its extinction point is unique, we have
\begin{align}
        R(x_0)
        =
        x_0.
\end{align}
Thus $x_0\in\{z=0\}$. By
\cite[Theorem~1.1]{gerhardt2015curvature},
\begin{align}
        E
        =
        x_0^\perp\cap\bS^{d+1}.
\end{align}
Hence $E$ is $R$-invariant. Moreover, if
$e_z:=(0,1)\in\bR^{d+1}\times\bR$, then $e_z\in E$, so $R|_E$ is a
nontrivial reflection and
\begin{align}
        \operatorname{Fix}(R|_E)
        =
        E\cap\{z=0\}
\end{align}
is a round equator of $E$. For $t$ sufficiently close
to $T_*$, let
\begin{align}
        Y_t:E\to X_t(\Sigma)
\end{align}
be the geodesic graph parametrization. Its uniqueness and
\eqref{eq:hemisphere-smooth-flow-equivariance} imply that
\begin{align}
        Y_t\circ R|_E
        =
        R\circ Y_t.
\end{align}
Set
\begin{align}
        P_t
        :=
        X_t^{-1}\circ Y_t:
        E\to\Sigma.
\end{align}
Consequently, the diffeomorphisms satisfy
\begin{align}
        P_t\circ R|_E
        =
        R|_\Sigma\circ P_t.
\end{align}
Since $P_t$ is an $R$-equivariant diffeomorphism, we have $ P_t(\operatorname{Fix}(R|_E)) = \operatorname{Fix}(R|_\Sigma).$ Furthermore,
\begin{align}
        \Sigma\setminus\operatorname{Fix}(R|_\Sigma)
        =
        \bigl(\Sigma\cap\{z>0\}\bigr)
        \sqcup
        \bigl(\Sigma\cap\{z<0\}\bigr).
\end{align}
The first set on the right-hand side is connected by assumption, and the
second is connected by $R$-symmetry. Thus these are precisely the two
connected components of the complement of the fixed-point set. Let
$E^+$ be the closure of one component of
$E\setminus\operatorname{Fix}(R|_E)$. Since $P_t$ is an
$R$-equivariant diffeomorphism, after interchanging the two choices of
$E^+$ once, we have
\begin{align}\label{eq:hemisphere-smooth-half-image-prelimit}
        P_t(E^+)
        =
        \Sigma^+
\end{align}
for every $t$ sufficiently close to $T_*$. The choice is independent of $t$ because the family $P_t$ depends continuously on $t$ and the two possible component assignments form a
discrete set. Let
\begin{align}
        g_t:=X_t^*\overline g_{\can},
        \qquad
        \widehat g_t:=Y_t^*\overline g_{\can}.
\end{align}
Let $h_t$ denote the second fundamental form and set
\begin{align}
        H_t
        :=
        \operatorname{tr}_{g_t}h_t.
\end{align}
As in the previous section, we have
\begin{align}\label{eq:hemisphere-smooth-metric-volume-evolution}
        \partial_tg_t
        =
        \frac{2d}{H_t}h_t
        \geq0,
        \quad
        \partial_t\dvol_{g_t}
        =
        d\dvol_{g_t}.
\end{align}
Since $g_0=g_\Sigma$, it follows that
\begin{align}\label{eq:hemisphere-smooth-metric-volume-consequences}
        g_t\geq g_\Sigma,
        \quad
        \dvol_{g_t}
        =
        e^{dt}\dvol_{g_\Sigma}.
\end{align}
Since $X_t\circ P_t=Y_t$, we also have $P_t^*g_t=\widehat g_t$. Thus, setting $P_t^+:=P_t|_{E^+}$ and using
\eqref{eq:hemisphere-smooth-half-image-prelimit},
\begin{align}\label{eq:hemisphere-smooth-half-prelimit}
        (P_t^+)^*g_\Sigma
        \leq
        \widehat g_t|_{E^+},
        \quad 
        (P_t^+)_\#\left(
        \dvol_{\widehat g_t}|_{E^+}
        \right)
        =
        e^{dt}\dvol_{g_\Sigma}|_{\Sigma^+}.
\end{align}
Moreover,
\begin{align}
        \widehat g_t
        \rightarrow
        g_E
        \qquad
        \text{in }C^\infty(E),
\end{align}
and, by the second identity in \eqref{eq:hemisphere-smooth-half-prelimit} we get
\begin{align}\label{eq:hemisphere-smooth-half-terminal-factor}
        \vol_{g_E}(E^+)
        =
        e^{dT_*}\vol_{g_\Sigma}(\Sigma^+).
\end{align}
Since $\widehat g_t\to g_E$ smoothly, the metrics $\widehat g_t$ are
uniformly bounded above by a fixed multiple of $g_E$ for $t$ sufficiently
close to $T_*$. The first inequality in
\eqref{eq:hemisphere-smooth-half-prelimit} therefore gives equicontinuity
of $P_t^+$. Hence, after passing to a sequence $t_j\to T_*$,
we have
\begin{align}
        P_{t_j}^+
        \rightarrow P
        \quad
        \text{uniformly on }E^+.
\end{align}
Since $E^+$ is a geodesically convex round hemisphere, the intrinsic
distance induced by $g_E|_{E^+}$ agrees with the restriction of
$d_{g_E}$. Passing to the limit gives \eqref{eq:hemisphere-smooth-half-contraction}. If
$f\in C^0(\Sigma^+)$, then the second line of
\eqref{eq:hemisphere-smooth-half-prelimit} gives
\begin{align}
        \int_{E^+}f(P_{t_j}^+(x))\dvol_{\widehat g_{t_j}}(x)
        =
        e^{dt_j}\int_{\Sigma^+}f\dvol_{g_\Sigma}.
\end{align}
Using the uniform convergence, the smooth convergence of $\widehat g_{t_j}$,
and \eqref{eq:hemisphere-smooth-half-terminal-factor}, we obtain
\eqref{eq:hemisphere-smooth-half-pushforward}. Finally,
$P(E^+)\subset\Sigma^+$ is compact, while the measure on the right-hand side
of \eqref{eq:hemisphere-smooth-half-pushforward} has support equal to
$\Sigma^+$. Hence $P(E^+)=\Sigma^+$.
\end{proof}

\begin{lem}\label{lem:hemisphere-pancake-approximation}
Let $d\geq2$, let $\bS^d_+$ be a closed hemisphere of the equatorial copy
\begin{align}
        \bS^d
        =
        \bS^{d+1}\cap\{z=0\},
\end{align}
and let $K\subset\bS^d_+$ be a closed geodesically convex set satisfying $ \vol_{g_{\can}}(K)>0.$ Then there exist smooth closed strictly convex hypersurfaces $\Sigma_j \subset (\bS^{d+1},\overline g_{\can})$ contained in an open hemisphere and invariant under $R(x,z)=(x,-z)$ such that, setting $ \Sigma_j^+:=\Sigma_j\cap\{z\geq0\},$ we have
\begin{align}\label{eq:hemisphere-pancake-Hausdorff-limit}
        d_H(\Sigma_j^+,K)
        \rightarrow
        0,
\end{align}
and
\begin{align}\label{eq:hemisphere-pancake-measure-limit}
        \dvol_{g_{\Sigma_j}}|_{\Sigma_j^+}
        \wto
        \mathbf{1}_K\dvol_{g_{\can}}
\end{align}
as $j\to \infty$. Moreover, $ \Sigma_j\cap\{z>0\}$ is connected for every $j$. Here $d_H$ denotes the Hausdorff distance between compact subsets of $(\bS^{d+1},\overline g_{\can})$.
\end{lem}

\begin{proof}
We first construct an approximation of $K$ that is away from the boundary of the hemisphere. Let $N\in\bS^d$ denote the pole of the hemisphere $\bS^d_+$, so that
\begin{align}
        \bS^d_+
        =
        \{x\in\bS^d:x\cdot N\geq0\}.
\end{align}
Since $K$ is geodesically convex and has positive $d$-dimensional volume,
it has nonempty interior. Choose $p\in\operatorname{int}(K)$. Since $K\subset\bS^d_+$, we have
$p\cdot N>0$. For each $j\geq1$, set
\begin{align}
        \delta_j
        :=
        \frac{p\cdot N}{j+1},
        \qquad
        K_j
        :=
        K\cap\{x\cdot N\geq\delta_j\}.
\end{align}
Then the sets $K_j$ are compact, increasing, geodesically convex, and
contained in the open hemisphere $\{x\cdot N>0\}$. Moreover,
$p\in\operatorname{int}(K_j)$ for every $j\geq1$. We claim that
\begin{align}
        \overline{\bigcup_{j=1}^\infty K_j}
        =
        K.
\end{align}
Indeed, every $q\in K$ satisfying $q\cdot N>0$ belongs to $K_j$ for all
sufficiently large $j$. If instead $q\cdot N=0$, then $p$ and $q$ are not antipodal since $p\cdot N>0$. Thus, if we define, for $0\leq s<1$,
\begin{align}
        q_s
        :=
        \frac{(1-s)p+sq}{|(1-s)p+sq|},
\end{align}
then the point $q_s$ lies on the minimizing geodesic segment from $p$ to $q$ and
therefore belongs to $K$ by geodesic convexity. Furthermore,
\begin{align}
        q_s\cdot N
        =
        \frac{(1-s)p\cdot N}{|(1-s)p+sq|}
        >
        0,
\end{align}
so $q_s\in K_j$ for all sufficiently large $j$. Since $q_s\to q$ as
$s\to1$, the claim follows. Since $K_j\subset K$, the preceding claim and the compactness of $K$ imply
\begin{align}
        d_H(K_j,K)
        \rightarrow
        0
\end{align}
as $j\to \infty.$ Moreover, as $j\to \infty$
\begin{align}
        \mathbf{1}_{K_j}
        \rightarrow
        \mathbf{1}_K
        \quad
        \dvol_{g_{\can}}\text{-almost everywhere},
\end{align}
because the equator $\{x\cdot N=0\}$ has zero $d$-dimensional volume. Consequently, by the dominated convergence theorem,
\begin{align}\label{eq:hemisphere-pancake-truncation-limit}
        \mathbf{1}_{K_j}\dvol_{g_{\can}}
        \wto
        \mathbf{1}_K\dvol_{g_{\can}}.
\end{align}
Next, we further approximate the sets $K_j$ by convex sets with smooth uniformly
convex boundaries. Regard $N$ also as the point $(N,0)\in\bS^{d+1}\subset\bR^{d+1}\times\bR$,
and set $\cH_N:=\{Y\in\bS^{d+1}:Y\cdot N>0\}.$ Identify
\begin{align}
        T_N\bS^{d+1}
        =
        T_N\bS^d\oplus\operatorname{span}\{\partial_z\}
        \simeq
        \bR^d\times\bR.
\end{align}
Define the gnomonic projection based at $N$ by
\begin{align}
        \Gamma:\cH_N
        \rightarrow
        T_N\bS^{d+1},\quad 
        \Gamma(Y)
        :=
        \frac{Y}{Y\cdot N}-N.
\end{align}
Its inverse is
\begin{align}
        G:T_N\bS^{d+1}
        \rightarrow
        \cH_N,\quad 
        G(V)
        :=
        \frac{N+V}{\sqrt{1+|V|^2}}.
\end{align}
The restriction of $\Gamma$ to
$\cH_N\cap\{z=0\}$ takes values in
$T_N\bS^d\simeq\bR^d$. Since
\begin{align}
        K_j
        \subset
        \{x\in\bS^d:x\cdot N\geq\delta_j\}
        \subset
        \cH_N,
\end{align}
we may define
\begin{align}
        A_j
        :=
        \Gamma(K_j)
        \subset
        T_N\bS^d
        \simeq
        \bR^d.
\end{align}
Because the gnomonic projection maps geodesic segments contained in
$\cH_N$ onto Euclidean line segments, $A_j$ is compact and convex.
Furthermore, since $\Gamma$ is a diffeomorphism and $K_j$ has nonempty
interior, $A_j$ has nonempty interior. We next smooth $A_j$. Let
\begin{align}
        h_j(u)
        :=
        \sup_{x\in A_j}x\cdot u,
        \quad 
        f_j(x)
        :=
        \max_{u\in\bS^{d-1}}\bigl(x\cdot u-h_j(u)\bigr).
\end{align}
Then $f_j$ is convex and $1$-Lipschitz,
\begin{align}
        A_j
        =
        \{f_j\leq0\},
        \qquad
        \max\{f_j,0\}
        =
        \dist(\cdot,A_j).
\end{align}
Fix a nonnegative standard mollifier $\varphi_\eta$ supported in
$B_\eta(0)$ and satisfying $ \int_{\bR^d}\varphi_\eta\,\ud x = 1$. Choose $R_j$ such that $A_j\subset B_{R_j}(0)$, and define
\begin{align}
        F_{j,\eta}
        :=(f_j*\varphi_\eta)+\eta|x|^2,
        \quad 
        c_{j,\eta}
        :=2\eta(1+R_j^2),
        \\
        A_{j,\eta}
        :=\{F_{j,\eta}\leq c_{j,\eta}\},
        \quad 
        \rho_{j,\eta}
        :=F_{j,\eta}-c_{j,\eta}.
\end{align}
Since $\|f_j*\varphi_\eta-f_j\|_{C^0}\leq\eta$, we have
\begin{align}\label{eq:hemisphere-pancake-smooth-approximation}
        A_j
        \subset
        A_{j,\eta}
        \subset
        A_j+(c_{j,\eta}+\eta)\overline B_1(0),
        \qquad
        D^2\rho_{j,\eta}
        \geq
        2\eta I.
\end{align}
Since $f_j<0$ at every interior point of $A_j$, for all sufficiently small
$\eta>0$ the level $c_{j,\eta}$ lies strictly above the unique minimum of
the strictly convex function $F_{j,\eta}$. Hence $\nabla\rho_{j,\eta}\neq0$ on $\partial A_{j,\eta}$. Therefore, $A_{j,\eta}$ is a smooth uniformly convex body and
$d_H(A_{j,\eta},A_j)\to0$ as $\eta\to0$. Set
\begin{align}
        S_{j,\eta}
        :=
        G(A_{j,\eta}\times\{0\}).
\end{align}
Let $G_0(v)=G(v,0)$ for $v\in T_N \bS^{d}.$ For $v,\xi\in T_N\bS^d$, we have
\begin{align}
        d(G_0)_v(\xi)
        =
        \frac{(\xi,0)}{\sqrt{1+|v|^2}}
        -
        \frac{\la v,\xi\ra(N+v,0)}
        {(1+|v|^2)^{3/2}}.
\end{align}
It follows that
\begin{align}
        (G_0^*g_{\can})_v(\xi,\zeta)
        =
        \frac{\la\xi,\zeta\ra}{1+|v|^2}
        -
        \frac{\la v,\xi\ra\la v,\zeta\ra}
        {(1+|v|^2)^2}.
\end{align}
Therefore,
\begin{align}\label{eq:hemisphere-pancake-gnomonic-Jacobian}
        G_0^*\dvol_{g_{\can}}
        =
        J_{G_0}(v)\,\ud v,
        \qquad
        J_{G_0}(v)
        :=
        (1+|v|^2)^{-\frac{d+1}{2}}.
\end{align}

Fix $j$. Since
$d_H(A_{j,\eta},A_j)\to0$, all the sets $A_{j,\eta}$ with $\eta>0$
sufficiently small are contained in a fixed compact subset of
$T_N\bS^d$. The map $G_0$ is Lipschitz on this compact set, and hence
\begin{align}
        d_H(S_{j,\eta},K_j)
        \leq
        C_jd_H(A_{j,\eta},A_j)
        \rightarrow
        0,
\end{align}
as $\eta\to 0$. Moreover, by
\eqref{eq:hemisphere-pancake-smooth-approximation}, we have
\begin{align}
        A_j
        \subset
        A_{j,\eta}
        \subset
        A_j+r_{j,\eta}\overline B_1(0),
        \qquad
        r_{j,\eta}
        :=
        c_{j,\eta}+\eta
        \rightarrow
        0,
\end{align}
as $\eta\to 0.$ Consequently, as $\eta \to 0$
\begin{align}
        \mathbf{1}_{A_{j,\eta}}(v)
        \rightarrow
        \mathbf{1}_{A_j}(v)
        \quad
        \text{for every }v\in T_N\bS^d.
\end{align}
Let $f\in C^0(\bS^{d+1})$. Using
\eqref{eq:hemisphere-pancake-gnomonic-Jacobian}, we obtain
\begin{align}
        \int_{S_{j,\eta}}f\,\dvol_{g_{\can}}
        =
        \int_{T_N\bS^d}
        f(G_0(v))
        \mathbf{1}_{A_{j,\eta}}(v)
        J_{G_0}(v)\,\ud v.
\end{align}
Thus, the dominated convergence theorem gives
\begin{align}
        \int_{S_{j,\eta}}f\,\dvol_{g_{\can}}
        \rightarrow
        \int_{T_N\bS^d}
        f(G_0(v))
        \mathbf{1}_{A_j}(v)
        J_{G_0}(v)\,\ud v
        =
        \int_{K_j}f\,\dvol_{g_{\can}}.
\end{align}
Therefore, $ \mathbf{1}_{S_{j,\eta}}\dvol_{g_{\can}} \wto \mathbf{1}_{K_j}\dvol_{g_{\can}}$ as $\eta \to 0$. Choose $0<\eta_j<1/j$ diagonally so that, setting
\begin{align}
        A'_j
        :=
        A_{j,\eta_j},
        \quad 
        \rho_j
        :=
        \rho_{j,\eta_j},
        \quad 
        S_j
        :=
        S_{j,\eta_j},
\end{align}
we have as $j\to \infty$
\begin{align}\label{eq:hemisphere-pancake-smooth-limit}
        d_H(S_j,K)
        \rightarrow
        0,
        \quad
        \mathbf{1}_{S_j}\dvol_{g_{\can}}
        \wto
        \mathbf{1}_K\dvol_{g_{\can}}.
\end{align}
We now thicken each $S_j$ symmetrically in the additional $z$-direction
to obtain a smooth closed strictly convex hypersurface in $\bS^{d+1}$. Here the reader should think about the simple case of the unit interval $[-1,1]$ which can be fattened to obtain an ellipse that collapses onto the plane from upper and lower side. Thus, for $\varepsilon>0$, define
\begin{align}
        B_{j,\varepsilon}
        :=
        \left\{
        (x,z)\in\bR^d\times\bR:
        \rho_j(x)+\frac{z^2}{\varepsilon^2}\leq0
        \right\}.
\end{align}
The Hessian of the function $\rho_j+z^2/\veps^2$ is
\begin{align}
        \begin{pmatrix}
        D^2\rho_j&0\\
        0&2\varepsilon^{-2}
        \end{pmatrix}
        >
        0,
\end{align}
and its gradient does not vanish on the zero level. Thus
$\widehat\Sigma_{j,\varepsilon}:=\partial B_{j,\varepsilon}$ is a smooth
closed strictly convex Euclidean hypersurface invariant under
$(x,z)\mapsto(x,-z)$. Now observe that if $\zeta$ is the Euclidean unit normal to
$\widehat\Sigma_{j,\varepsilon}$ at $y$, then for tangent vectors $v,w$,
\begin{align}
        \operatorname{II}_{\bS^{d+1}}(dG_y(v),dG_y(w))
        =
        \frac{
        \operatorname{II}_{\bR^{d+1}}(v,w)
        }{
        \sqrt{1+|y|^2}
        \sqrt{1+(y\cdot\zeta)^2}
        }.
\end{align}
Thus, the map $G$ preserves positivity of the second fundamental form and therefore, $\Sigma_{j,\varepsilon} := G(\widehat\Sigma_{j,\varepsilon})$ is a smooth closed strictly convex hypersurface in an open hemisphere and is
invariant under $R$.

The upper half of $\widehat\Sigma_{j,\varepsilon}$ is the graph of $ z_{j,\varepsilon}(x)= \varepsilon\sqrt{-\rho_j(x)}$ for $x\in A_j'$ and hence $\Sigma_{j,\varepsilon}^+$ is parametrized by
\begin{align}
        \Psi_{j,\varepsilon}:\operatorname{int}(A'_j)
        \rightarrow
        \Sigma_{j,\varepsilon}\cap\{z>0\},
        \quad 
        \Psi_{j,\varepsilon}(x)
        :=
        G\bigl(x,z_{j,\varepsilon}(x)\bigr).
\end{align}
The same formula extends $\Psi_{j,\varepsilon}$ continuously to $A'_j$;
we continue to denote this extension by $\Psi_{j,\varepsilon}$. Since
$G$ preserves the sign of the last coordinate, we have
\begin{align}
        \Psi_{j,\varepsilon}
        \left(
        \operatorname{int}(A'_j)
        \right)
        =
        \Sigma_{j,\varepsilon}\cap\{z>0\},
        \quad
        \Psi_{j,\varepsilon}(A'_j)
        =
        \Sigma_{j,\varepsilon}^+.
\end{align}
 We have
\begin{align}
        \widehat\Sigma_{j,\varepsilon}\cap\{z>0\}
        =
        \{
        (x,\varepsilon\sqrt{-\rho_j(x)}):
        x\in\operatorname{int}(A'_j)
        \}.
\end{align}
Since $A'_j$ is convex, its interior is connected. The inverse gnomonic
projection $G$ preserves the sign of the $z$-coordinate, and hence
$\Sigma_{j,\varepsilon}\cap\{z>0\}$ is connected. Since $\rho_j=0$ and $\nabla\rho_j\neq0$ on $\partial A'_j$, there exist
constants $c_j,C_j,r_j>0$ such that
\begin{align}\label{eq:hemisphere-pancake-defining-function-distance}
        c_j\dist(x,\partial A'_j)
        \leq
        -\rho_j(x)
        \leq
        C_j\dist(x,\partial A'_j)
\end{align}
whenever $x\in A'_j$ and
$\dist(x,\partial A'_j)<r_j$. Moreover,
for $x\in\operatorname{int}(A'_j)$,
\begin{align}
        \nabla z_{j,\varepsilon}(x)
        =
        -\frac{\varepsilon\nabla\rho_j(x)}
        {2\sqrt{-\rho_j(x)}}.
\end{align}
Consequently, after increasing $C_j$ if necessary,
\begin{align}\label{eq:hemisphere-pancake-graph-gradient-bound}
        |\nabla z_{j,\varepsilon}(x)|
        \leq
        \frac{C_j\varepsilon}
        {\sqrt{\dist(x,\partial A'_j)}}
        \leq
        \frac{C_j}
        {\sqrt{\dist(x,\partial A'_j)}}
\end{align}
for $0<\varepsilon\leq1$. In particular, for every fixed
$x\in\operatorname{int}(A'_j)$, $z_{j,\varepsilon}(x)
        \rightarrow
        0,
        |\nabla z_{j,\varepsilon}(x)|
        \rightarrow
        0$ as $\varepsilon\to0$. Since $A'_j$ is compact, $\sup_{x\in A'_j}|z_{j,\varepsilon}(x)| \leq C_j\varepsilon.$ The map $G$ is Lipschitz on a fixed compact set containing all the graphs
for $0<\varepsilon\leq1$. Recalling that $ S_j
        =
        G(A'_j\times\{0\}),$ we obtain
\begin{align}\label{eq:hemisphere-pancake-upper-sheet-Hausdorff}
        d_H(\Sigma_{j,\varepsilon}^+,S_j)
        \leq
        \sup_{x\in A'_j}
        d_{\overline g_{\can}}
        \left(
        G\bigl(x,z_{j,\varepsilon}(x)\bigr),
        G(x,0)
        \right)
        \leq
        C_j\varepsilon
        \rightarrow
        0,
\end{align}
as $\varepsilon\to 0.$ We next prove convergence of the volume measures. First note that, for $y\in T_N\bS^{d+1}$ and $v,w\in T_N\bS^{d+1}$
\begin{align}
        (G^*\overline g_{\can})_y(v,w)
        =
        \frac{\la v,w\ra}{1+|y|^2}
        -
        \frac{\la y,v\ra\la y,w\ra}
        {(1+|y|^2)^2}.
\end{align}
Applying this formula with
\begin{align}
        y
        =
        \bigl(x,z_{j,\varepsilon}(x)\bigr),
        \qquad
        v
        =
        \bigl(\xi,\nabla z_{j,\varepsilon}(x)\cdot\xi\bigr),
\end{align}
and computing the determinant of the induced metric gives
\begin{align}\label{eq:hemisphere-pancake-spherical-graph-Jacobian}
        J_{j,\varepsilon}(x)
        :=
        \sqrt{
        \det\bigl(\Psi_{j,\varepsilon}^*
        \overline g_{\can}\bigr)
        }
        =
        \frac{
        \sqrt{
        1+|\nabla z_{j,\varepsilon}(x)|^2
        +
        \left(
        z_{j,\varepsilon}(x)
        -
        x\cdot\nabla z_{j,\varepsilon}(x)
        \right)^2
        }
        }{
        \left(
        1+|x|^2+z_{j,\varepsilon}(x)^2
        \right)^{\frac{d+1}{2}}
        }.
\end{align}
Thus, for every $x\in\operatorname{int}(A'_j)$, $ J_{j,\varepsilon}(x) \to \frac{1}{(1+|x|^2)^{\frac{d+1}{2}}}$ as $\varepsilon \to 0.$ Since $A'_j$ is bounded, it follows from \eqref{eq:hemisphere-pancake-graph-gradient-bound} that
\begin{align}
        J_{j,\varepsilon}(x)
        \leq
        C_j ( 1+
        \dist(x,\partial A'_j)^{-1/2}).
\end{align}
The right-hand side is integrable on $A'_j$.  Indeed, in a smooth collar neighborhood of the boundary we have
\begin{align}
        \int_0^{r_j}s^{-1/2}\,\ud s
        <
        \infty.
\end{align}
Since $\partial A'_j$ has zero $d$-dimensional measure, the area formula
on $\operatorname{int}(A'_j)$ gives for any $f\in C^0(\bS^{d+1})$
\begin{align}
        \int_{\Sigma_{j,\varepsilon}^+}
        f\,\dvol_{g_{\Sigma_{j,\varepsilon}}}
        =
        \int_{A'_j}
        f\bigl(\Psi_{j,\varepsilon}(x)\bigr)
        J_{j,\varepsilon}(x)\,\ud x.
\end{align}
Thus the dominated convergence theorem implies
\begin{align}
        \int_{\Sigma_{j,\varepsilon}^+}
        f\,\dvol_{g_{\Sigma_{j,\varepsilon}}}
        \rightarrow
        \int_{A'_j}
        f(G(x,0))
        \frac{\ud x}{(1+|x|^2)^{\frac{d+1}{2}}}
        =
        \int_{S_j}f\,\dvol_{g_{\can}}.
\end{align}
Therefore, $ \dvol_{g_{\Sigma_{j,\varepsilon}}} |_{\Sigma_{j,\varepsilon}^+} \wto \mathbf{1}_{S_j}\dvol_{g_{\can}}$ as $\veps\to 0.$ Choosing $0<\varepsilon_j<1/j$ diagonally and setting $\Sigma_j:=\Sigma_{j,\varepsilon_j}$, the conclusion follows from
\eqref{eq:hemisphere-pancake-smooth-limit}.
\end{proof}

\begin{proof}[Proof of Theorem~\ref{thm:hemisphere-transport}]
Apply Lemma~\ref{lem:hemisphere-pancake-approximation} to $K$. For each $j$,
apply Lemma~\ref{lem:hemisphere-smooth-symmetric-transport} to $\Sigma_j$, and
choose a round isometry from $\bS^d_+$ onto the resulting source hemisphere.
We obtain continuous surjective maps $T_j:\bS^d_+\to\Sigma_j^+$ such that
\begin{align}\label{eq:hemisphere-prelimit-map}
        d_{\overline g_{\can}}(T_j(x),T_j(y))
        \leq
        d_{g_{\can}}(x,y),  \quad 
        (T_j)_\#\sigma_+
        =
        \frac{1}{\vol_{g_{\Sigma_j}}(\Sigma_j^+)}
        \dvol_{g_{\Sigma_j}}|_{\Sigma_j^+}.
\end{align}
Here the first inequality follows from
\eqref{eq:hemisphere-smooth-half-contraction}, since ambient distance is
bounded by intrinsic distance on $\Sigma_j$.

By the Arzel\`a--Ascoli theorem, after passing to a subsequence,
$T_j\to T$ uniformly on $\bS^d_+$. The Hausdorff convergence in
\eqref{eq:hemisphere-pancake-Hausdorff-limit} gives $ T(\bS^d_+) \subset K$. Passing to the limit in the first line of \eqref{eq:hemisphere-prelimit-map}, and using that ambient distance agrees with round distance on the equatorial copy of $\bS^d$, proves
\eqref{eq:hemisphere-transport-contraction}. 

Testing \eqref{eq:hemisphere-pancake-measure-limit} against the constant
function $1$, we obtain as $j\to \infty$
\begin{align}\label{eq:hemisphere-pancake-mass-limit}
        \vol_{g_{\Sigma_j}}(\Sigma_j^+)
        \rightarrow
        \vol_{g_{\can}}(K).
\end{align}
Let $f\in C^0(\bS^{d+1})$. By
\eqref{eq:hemisphere-prelimit-map}, we have
\begin{align}
        \int_{\bS^d_+}
        f(T_j(x))\,\ud\sigma_+(x)
        =
        \frac{1}{\vol_{g_{\Sigma_j}}(\Sigma_j^+)}
        \int_{\Sigma_j^+}
        f\,\dvol_{g_{\Sigma_j}}.
\end{align}
Using the uniform convergence $T_j\to T$,
\eqref{eq:hemisphere-pancake-measure-limit}, and
\eqref{eq:hemisphere-pancake-mass-limit}, we may let $j\to\infty$ to obtain
\begin{align}
        \int_{\bS^d_+}
        f(T(x))\,\ud\sigma_+(x)
        =
        \frac{1}{\vol_{g_{\can}}(K)}
        \int_K f\,\dvol_{g_{\can}}.
\end{align}
By the definition of $\sigma_+$, this is \eqref{eq:hemisphere-transport-pushforward}. Finally, since $K$ is geodesically convex and has nonempty interior, we have $K=\overline{\operatorname{int}(K)}.$ Hence the measure on the right-hand side of \eqref{eq:hemisphere-transport-pushforward} has support equal to $K$. Since $T(\bS^d_+)$ is compact, it follows that $T(\bS^d_+)=K.$
\end{proof}
\section{Proof of Theorem~\ref{thm:gaussian-main}}
The idea behind Theorem~\ref{thm:gaussian-main} is to deform the Euclidean
metric in polar coordinates.  Recall that the standard metric on \(\mathbb R^d\)
is
\[
        \ud r^2+r^2g_{\bbS^{d-1}},
\]
so the spherical slices \(\{r=\mathrm{constant}\}\) have radius \(r\) and keep
expanding as \(r\) increases.  In our construction, we instead choose a
rotationally symmetric metric for which these slices expand in the usual way
near the origin, so that the metric closes smoothly, but then their radius
becomes constant and equal to \(\sqrt{d-2}\).  Thus the end of the manifold is
cylindrical, with cross-section \(\bbS^{d-1}(\sqrt{d-2})\).  This cylinder is the
source of the spectral obstruction: its cross-section has first nonzero
eigenvalue
\begin{align}
        \frac{d-1}{d-2}<2,
\end{align}
attained by the \(d\) coordinate functions restricted to
\(\bbS^{d-1}(\sqrt{d-2})\), while the cylindrical direction produces one
additional low-energy radial mode.  Finally, we put a quadratic potential in the
cylindrical direction, so that the end carries a Gaussian weight.  The weight makes the total measure of the space finite and provides enough
curvature for the curvature--dimension condition
\(\operatorname{CD}(1,\infty)\). Thus, the constructed weighted
manifold satisfies the curvature-dimension lower bound, but its cylindrical end
produces \(d+1\) nonzero test functions with Rayleigh quotient below the
Gaussian threshold \(2\), forcing
\begin{align}
        \lambda_{d+2}(\mathbb R^d,g,\mu)
        <
        2
        =
        \lambda_{d+2}(\mathbb R^d,|\cdot|,\gamma^d).
\end{align}
Before proceeding with the proof of Theorem~\ref{thm:gaussian-main}, we record a
technical smoothing lemma used to construct the torpedo function (cf. \cite{botvinnik2010homotopy}).
\begin{lem}\label{lem:theta-eps}
Let $\delta>0$, $B=\frac{\pi\delta}{2}$, and $0<\veps<B/2$. There exists a smooth function
$\te_\veps:[0,\infty)\to[0,\pi/2]$ such that
\begin{align}
        \te_\veps(r)
        &=
        \frac{r}{\delta}
        \qquad\text{when }0\leq r\leq B-\veps,\\
        \te_\veps(r)
        &=
        \frac{\pi}{2}
        \qquad\text{when }r\geq B+\veps,
\end{align}
with $\te_\veps'\geq0$, $\te_\veps'\leq\delta^{-1}$, and $\te_\veps''\leq0$.
\end{lem}
\begin{proof}
Define
\[
        \psi(s)
        =
        \begin{cases}
        \exp\!\left(-\dfrac{1}{1-4s^2}\right),& |s|<\dfrac12,\\[6pt]
        0,& |s|\geq \dfrac12.
        \end{cases}
\]
Then \(\psi\in C_c^\infty((-1,1))\), \(\psi\geq0\), \(\psi\) is even, and
\(\psi\) is not identically zero. Set
\[
        D:=\int_{-1}^1\psi(u)\,\ud u>0
\]
and define
\begin{align}
        \eta(s)
        =
        \frac{1}{D}\int_s^1\psi(u)\,\ud u .
\end{align}
Since \(\psi\geq0\), we have \(0\leq\eta\leq1\), and
\[
        \eta'(s)=-\frac{\psi(s)}{D}\leq0.
\]
Moreover, since \(\psi\) is supported in \([-1/2,1/2]\), we have
\[
        \eta(s)=1\quad\text{for }s\leq -\frac12,
        \qquad
        \eta(s)=0\quad\text{for }s\geq \frac12.
\]
Thus \(\eta\) is constant near \(s=-1\) and near \(s=1\). Finally, since
\(\psi\) is even, we have
\[
        \eta(-s)=1-\eta(s).
\]
Therefore,
\[
        \int_{-1}^1\eta(s)\,\ud s
        =
        \frac12\int_{-1}^1\bigl(\eta(s)+\eta(-s)\bigr)\,\ud s
        =
        1.
\]

Define
\begin{align}
        \te_\veps(r)=
        \begin{cases}
        \dfrac{r}{\delta},&0\leq r\leq B-\veps,\\[6pt]
        \dfrac{B-\veps}{\delta}
        +\dfrac{\veps}{\delta}
        \displaystyle\int_{-1}^{(r-B)/\veps}\eta(s)\,\ud s,
        &B-\veps\leq r\leq B+\veps,\\[12pt]
        \dfrac{\pi}{2},&r\geq B+\veps.
        \end{cases}
\end{align}
When \(B-\veps\leq r\leq B+\veps\), we have
\[
        \te_\veps'(r)
        =
        \frac{1}{\delta}\eta\!\left(\frac{r-B}{\veps}\right),
        \qquad
        \te_\veps''(r)
        =
        \frac{1}{\delta\veps}
        \eta'\!\left(\frac{r-B}{\veps}\right)
        \leq0.
\]
Since \(0\leq\eta\leq1\), this gives
\[
        0\leq \te_\veps'(r)\leq \delta^{-1}.
\]
The same inequalities are immediate on the two outer regions, where
\(\te_\veps(r)=r/\delta\) and \(\te_\veps(r)=\pi/2\), respectively. Thus
\(\te_\veps'\geq0\), \(\te_\veps'\leq\delta^{-1}\), and
\(\te_\veps''\leq0\) on all of \([0,\infty)\).
\end{proof}

\begin{proof}[Proof of Theorem~\ref{thm:gaussian-main}]
Let \(d\geq4\) and let \(0<\veps\ll1\). The argument proceeds in a similar style as in the proof of Theorem~\ref{thm:sphere-main}. 
\paragraph{$\mathrm{(i)}$ \textit{Construction of the metric}} Set
\begin{align}
        \delta=\sqrt{d-2},
        \qquad
        B=\frac{\pi\delta}{2}.
\end{align}
Following the notation of \cite[Definition 2.4]{botvinnik2010homotopy}, we use a
torpedo function \(\rho_\delta\). Since we require more precise estimates on
\(\rho_{\delta,\veps}\), we define it explicitly. Let
\(\te_\veps:[0,\infty)\to[0,\pi/2]\) be the function from Lemma
\ref{lem:theta-eps}, and define
\begin{align}
        \rho_{\delta,\veps}(r)=\delta\sin\te_\veps(r).
\end{align}
Then
\begin{align}
        \rho_{\delta,\veps}(r)
        &=
        \delta\sin(r/\delta)
        \qquad\text{when }0\leq r\leq B-\veps,\\
        \rho_{\delta,\veps}(r)
        &=
        \delta
        \qquad\qquad\qquad\text{when }r\geq B+\veps.
\end{align}
Moreover, since
\begin{align}
        \rho_{\delta,\veps}''
        =
        \delta\cos\te_\veps\,\te_\veps''
        -
        \delta\sin\te_\veps\,(\te_\veps')^2,
\end{align}
and $0\leq\te_\veps\leq\pi/2$, we have \(\rho_{\delta,\veps}''\leq0\). Thus
\(\rho_{\delta,\veps}\) is a torpedo function in the sense of
\cite[Definition 2.4]{botvinnik2010homotopy}, except that we keep track of the
transition scale \(\veps\).
We now define the metric directly on \(\bbR^d\). On \(\bbR^d\setminus\{0\}\),
write
\begin{align}
        x=r\omega,
        \qquad
        r=|x|,
        \qquad
        \omega\in\bbS^{d-1}.
\end{align}
Define
\begin{align}\label{eq:noncompact-torpedo-metric}
        g_\veps
        =
        \ud r^2+\rho_{\delta,\veps}(r)^2g_{\bbS^{d-1}},
\end{align}
where \(g_{\bbS^{d-1}}\) is the canonical round metric on the unit sphere
\(\bbS^{d-1}\). As \(r\to0\) we have
\begin{align}
        \rho_{\delta,\veps}(r)
        =
        \delta\sin(r/\delta)
        =
        r-\frac{r^3}{6\delta^2}+O(r^5),
\end{align}
the metric smoothly extends across the origin as a rotationally symmetric metric
on \(\bbR^d\). Moreover,
\begin{align}
        g_\veps
        =
        \ud r^2+\delta^2g_{\bbS^{d-1}}
        \qquad\text{when }r\geq B+\veps.
\end{align}
This also shows that \(g_\veps\) is cylindrical at infinity and hence complete.

\paragraph{$\mathrm{(ii)}$ \textit{Construction of the measure}} Next, we will define a suitable weight which will be quadratic on the cylindrical end. To this end, choose a smooth function \(q_\veps:[0,\infty)\to[0,1]\) such
that
\begin{align}
        q_\veps=0\quad\text{on }[0,B-2\veps],
        \qquad
        q_\veps=1\quad\text{on }[B-\veps,\infty).
\end{align}
Then define \(V_\veps\) as
\begin{align}
        V_\veps(r)=\int_0^r(r-u)q_\veps(u)\,\ud u,
        \qquad
        V_\veps\geq0,
        \qquad
        V_\veps'\geq0.
\end{align}
Set
\begin{align}
        \nu_\veps=e^{-V_\veps}\dvol_{g_\veps},
        \qquad
        Z_\veps=\int_{\bbR^d}\ud\nu_\veps,
        \qquad
        \mu_\veps=Z_\veps^{-1}\nu_\veps.
\end{align}
Since \(V_\veps\) is bounded from below by a quadratic function in \(r\) for
\(r\geq B-\veps\), the normalizing constant \(Z_\veps<\infty\).

\paragraph{$\mathrm{(iii)}$ \textit{Verification of the $\CD(1,\infty)$ condition}} We now show that \((\bbR^d,g_\veps,\mu_\veps)\) is \(\CD(1,\infty)\). Let
\(\partial_r\) denote the unit vector in the radial direction and let
\(X\perp\partial_r\) be a \(g_\veps\)-unit vector. Since the potential
\(V_\veps=V_\veps(r)\) is radial, the Bakry--\'Emery tensor
\begin{align}
        \Ric_{g_\veps,V_\veps}:=\Ric_{g_\veps}+\nabla^2V_\veps
\end{align}
is given by (cf. proof of Proposition 2 in \cite{wylie2016some})
\begin{align}
        \Ric_{g_\veps,V_\veps}(\partial_r,\partial_r)
        &=
        -(d-1)\frac{\rho_{\delta,\veps}''}{\rho_{\delta,\veps}}
        +
        V_\veps'',\\
        \Ric_{g_\veps,V_\veps}(\partial_r,X)
        &=
        0,\\
        \Ric_{g_\veps,V_\veps}(X,X)
        &=
        -\frac{\rho_{\delta,\veps}''}{\rho_{\delta,\veps}}
        +
        (d-2)\frac{(1-(\rho_{\delta,\veps}')^2)}{\rho_{\delta,\veps}^2}
        +
        \frac{V_\veps'\rho_{\delta,\veps}'}{\rho_{\delta,\veps}}.
\end{align}
Since
\begin{align}
        \rho_{\delta,\veps}'\geq0,
        \qquad
        \rho_{\delta,\veps}''\leq0,
\end{align}
and since
\begin{align}
        \rho_{\delta,\veps}'
        =
        \delta\cos\te_\veps\,\te_\veps',
\end{align}
we have from Lemma \ref{lem:theta-eps} that
\begin{align}
        (\rho_{\delta,\veps}')^2
        \leq
        \cos^2\te_\veps
        =
        1-\frac{\rho_{\delta,\veps}^2}{\delta^2}.
\end{align}
Since \(\delta^2=d-2\), we get
\begin{align}
        (d-2)\bigl(1-(\rho_{\delta,\veps}')^2\bigr)
        \geq
        \rho_{\delta,\veps}^2.
\end{align}
Since also \(-\rho_{\delta,\veps}\rho_{\delta,\veps}''\geq0\), we get
\begin{align}
        -\frac{\rho_{\delta,\veps}''}{\rho_{\delta,\veps}}
        +
        (d-2)\frac{1-(\rho_{\delta,\veps}')^2}{\rho_{\delta,\veps}^2}
        &=
        \frac{(d-2)(1-(\rho_{\delta,\veps}')^2)
        -\rho_{\delta,\veps}\rho_{\delta,\veps}''}
        {\rho_{\delta,\veps}^2}\geq1.
\end{align}
Moreover,
\begin{align}
        V_\veps'\frac{\rho_{\delta,\veps}'}{\rho_{\delta,\veps}}\geq0,
\end{align}
because \(V_\veps'\geq0\) and \(\rho_{\delta,\veps}'\geq0\). Thus
\begin{align}
        \Ric_{g_\veps,V_\veps}(X,X)\geq1.
\end{align}
It remains to check the radial direction. If \(0<r\leq B-\veps\), then
\begin{align}
        \rho_{\delta,\veps}(r)
        =
        \delta\sin(r/\delta),
        \qquad
        \rho_{\delta,\veps}''
        =
        -\frac{1}{\delta^2}\rho_{\delta,\veps}.
\end{align}
Hence
\begin{align}
        \Ric_{g_\veps,V_\veps}(\partial_r,\partial_r)
        &=
        -(d-1)\frac{\rho_{\delta,\veps}''}{\rho_{\delta,\veps}}
        +
        V_\veps''\geq
        \frac{d-1}{\delta^2}=
        \frac{d-1}{d-2}
        >
        1.
\end{align}
If \(r\geq B-\veps\), then since
\begin{align}
        V_\veps'(r)
        =
        \int_0^r q_\veps(u)\,\ud u
        \implies
        V_\veps''(r)=q_\veps(r),
\end{align}
and the identity $q_\veps=1$ on $[B-\veps,\infty)$ gives \(V_\veps''(r)=1\). Combining this fact with
\(\rho_{\delta,\veps}''\leq0\), we get
\begin{align}
        \Ric_{g_\veps,V_\veps}(\partial_r,\partial_r)
        =
        -(d-1)\frac{\rho_{\delta,\veps}''}{\rho_{\delta,\veps}}
        +
        1
        \geq1.
\end{align}
Thus, on \(\bbR^d\setminus\{0\}\),
\begin{align}
        \Ric_{g_\veps}+\nabla^2V_\veps\geq g_\veps,
\end{align}
which by continuity extends to $r=0.$ Writing \(W_\veps=V_\veps+\log Z_\veps\), we have
\begin{align}
        \mu_\veps=e^{-W_\veps}\dvol_{g_\veps},
        \qquad
        \Ric_{g_\veps}+\nabla^2W_\veps
        =
        \Ric_{g_\veps}+\nabla^2V_\veps
        \geq g_\veps.
\end{align}
Hence \((\bbR^d,g_\veps,\mu_\veps)\) satisfies \(\CD(1,\infty)\).
\paragraph{$\mathrm{(iv)}$ \textit{Construction of test functions}}
We finally show the desired spectral inequality. To achieve this goal we will
construct suitable test functions. Since the normalization constants cancel in the
Rayleigh quotient, it suffices to use \(\nu_\veps\) instead of \(\mu_\veps\). Thus,
for any nonzero test function \(u\in W^{1,2}(\bbR^d,g_\veps,\nu_\veps)\), set
\begin{align}
        \cR_\veps(u)
        =
        \frac{\int_{\bbR^d}|\nabla u|_{g_\veps}^2\,\ud\nu_\veps}
        {\int_{\bbR^d}u^2\,\ud\nu_\veps}.
\end{align}
We would like to compare \(\cR_\veps\) with a simpler model Rayleigh quotient. To
this end define
\begin{align}
        \rho_0(r)=
        \begin{cases}
        \delta\sin(r/\delta),&0\leq r\leq B,\\
        \delta,&r\geq B,
        \end{cases}
        \qquad
        V_0(r)=
        \begin{cases}
        0,&0\leq r\leq B,\\
        \frac{(r-B)^2}{2},&r\geq B.
        \end{cases}
\end{align}
Let
\begin{align}
        \nu_0=e^{-V_0}\rho_0(r)^{d-1}\,\ud r\,\dvol_{\bbS^{d-1}}.
\end{align}
By construction, \(\rho_{\delta,\veps}=\rho_0\) on
\([0,B-\veps]\cup [B+\veps,\infty)\), while \(V_\veps=V_0\) on
\([0,B-2\veps]\). Next, observe that since \(q_\veps(u)=0\) on
\([0,B-2\veps]\) and \(0\leq q_\veps\leq1\), we have
\begin{align}
        V_\veps(B)
        &=
        \int_{B-2\veps}^{B}(B-u)q_\veps(u)\,\ud u
        \leq 2\veps^2,\\
        V_\veps'(B)
        &=
        \int_{B-2\veps}^{B}q_\veps(u)\,\ud u
        \leq 2\veps.
\end{align}
Furthermore, since \(V_\veps''(r)=1\) for all \(r\geq B\), we get
\begin{align}
        V_\veps'(r)-V_\veps'(B)
        &=
        \int_B^r V_\veps''(t)\,\ud t
        =
        r-B,\\
        V_\veps(r)-V_\veps(B)
        &=
        \int_B^r V_\veps'(t)\,\ud t
        =
        V_\veps'(B)(r-B)+\frac{1}{2}(r-B)^2.
\end{align}
Therefore, using the fact \(V_0(r)=\frac{1}{2}(r-B)^2\) for \(r\geq B\), we get
for \(\veps\leq1\),
\begin{align}
        0\leq V_\veps(r)-V_0(r)
        \leq
        2\veps(1+r-B).
\end{align}
Thus, using the elementary inequality \(1-e^{-x}\leq x\) for all \(x\geq0\), we
get for \(r\geq B\),
\begin{align}\label{eq:end-density-error}
        \left|e^{-V_\veps(r)}-e^{-V_0(r)}\right|
        \leq e^{-V_0(r)} |1-e^{-(V_\veps(r)-V_0(r))}|
        \leq
        2\veps(1+r-B)e^{-(r-B)^2/2}.
\end{align}
Let \(x_1,\dots,x_d\) be the coordinate functions on \(\bbR^d\) restricted to the
unit sphere \(\bbS^{d-1}\). In particular, they satisfy
\begin{align}
        -\Delta_{\bbS^{d-1}}x_i=(d-1)x_i,
        \qquad
        \int_{\bbS^{d-1}}x_i\,\dvol_{\bbS^{d-1}}=0,
\end{align}
and
\begin{align}
        \int_{\bbS^{d-1}}x_ix_j\,\dvol_{\bbS^{d-1}}=0
        \qquad(i\neq j).
\end{align}
Set
\begin{align}
        a_\veps(r)=\frac{\rho_{\delta,\veps}(r)}{\delta},
        \qquad
        U_{\veps,i}(r,\omega)=a_\veps(r)x_i(\omega).
\end{align}
Then \(|a_\veps|\leq1\), \(|a_\veps'|\leq\frac1\delta\), and \(U_{\veps,i}\)
are smooth as \(r\to0\). For \(r>0\),
\begin{align}
        |\nabla U_{\veps,i}|_{g_\veps}^2
        &=
        (a_\veps')^2x_i^2
        +
        \frac{a_\veps^2}{\rho_{\delta,\veps}^2}
        |\nabla_{\bbS^{d-1}}x_i|^2=
        (a_\veps')^2x_i^2
        +
        \frac{1}{\delta^2}|\nabla_{\bbS^{d-1}}x_i|^2.
\end{align}
Define
\begin{align}
        a_0(r)=\frac{\rho_0(r)}{\delta},
        \qquad
        U_{0,i}(r,\omega)=a_0(r)x_i(\omega),
\end{align}
and
\begin{align}
        \calD_U^0=\int U_{0,i}^2\,\ud\nu_0,
        \qquad
        \calN_U^0=\int |\nabla U_{0,i}|^2\,\ud\nu_0,
\end{align}
where \(a_0'\) is understood in the weak sense in \(\calN_U^0\). Denote \(I_\veps:=[B-2\veps,B+\veps]\). Then 
\begin{align}\label{eq:angular-error-estimate-new}
        &\left|
        \int U_{\veps,i}^2\,\ud\nu_\veps-\calD_U^0
        \right|
        +
        \left|
        \int |\nabla U_{\veps,i}|_{g_\veps}^2\,\ud\nu_\veps-\calN_U^0
        \right| \notag\\
        &\leq
        \int_0^\infty\int_{\bbS^{d-1}}
        \left|
        a_\veps^2x_i^2 e^{-V_\veps}\rho_{\delta,\veps}^{d-1}
        -
        a_0^2x_i^2 e^{-V_0}\rho_0^{d-1}
        \right|
        \,\dvol_{\bbS^{d-1}}\ud r \notag\\
        &\quad+
        \int_0^\infty\int_{\bbS^{d-1}}
        \left|
        \left((a_\veps')^2x_i^2+\frac1{\delta^2}|\nabla_{\bbS^{d-1}}x_i|^2\right)
        e^{-V_\veps}\rho_{\delta,\veps}^{d-1}
        -
        \left((a_0')^2x_i^2+\frac1{\delta^2}|\nabla_{\bbS^{d-1}}x_i|^2\right)
        e^{-V_0}\rho_0^{d-1}
        \right|
        \,\dvol_{\bbS^{d-1}}\ud r \notag\\
        &\leq
        C_d |I_\veps|
        +
        C_d\int_B^\infty
        \left|e^{-V_\veps(r)}-e^{-V_0(r)}\right|\,\ud r \notag\\
        &\leq
        C_d\veps
        +
        C_d\veps
        \int_B^\infty
        (1+r-B)e^{-(r-B)^2/2}\,\ud r \notag\\
        &\leq
        C_d\veps,
\end{align}
with constants $C_d>0$ depending on the dimension $d\geq 4$ and possibly different on each line. In the above display, we used that \(x_i\), \(\nabla_{\bbS^{d-1}}x_i\), and the
quantities
\[
        |a_\veps|\leq 1,\qquad |a_\veps'|\leq \frac1\delta,\qquad
        0\leq \rho_{\delta,\veps}\leq \delta,\qquad e^{-V_\veps}\leq 1
\]
are uniformly bounded on \(I_\veps\). Define \(\delta_U:=2\calD_U^0-\calN_U^0\), then 
\begin{align}\label{eqn:notation}
        \omega_{d-1}:=|\mathbb S^{d-1}|,
        \qquad
        I_m:=\int_0^{\pi/2}\sin^m t\,\ud t,
        \qquad
        J_0:=\int_0^\infty e^{-\tau^2/2}\,\ud\tau.
\end{align}
Using
\begin{align}
        \int_{\mathbb S^{d-1}}x_i^2\,\dvol_{\mathbb S^{d-1}}
        =
        \frac{\omega_{d-1}}{d},
        \qquad
        \int_{\mathbb S^{d-1}}|\nabla_{\mathbb S^{d-1}}x_i|^2\,\dvol_{\mathbb S^{d-1}}
        =
        \frac{(d-1)\omega_{d-1}}{d},
\end{align}
and \(\delta^2=d-2\), we compute
\begin{align}
        \delta_U
        &=
        \frac{\omega_{d-1}}{d}
        \left[
        2\delta^d I_{d+1}
        +
        2\delta^{d-1}J_0
        -
        \delta^{d-2}
        \int_0^{\pi/2}
        \bigl(\cos^2t+d-1\bigr)\sin^{d-1}t\,\ud t
        -
        (d-1)\delta^{d-3}J_0
        \right]  \\
        &=
        \frac{\omega_{d-1}}{d}
        \left[
        \delta^{d-2}\bigl((2d-3)I_{d+1}-dI_{d-1}\bigr)
        +
        \delta^{d-3}(d-3)J_0
        \right]  \\
        &=
        \frac{\omega_{d-1}}{d}
        \left[
        \delta^{d-2}\frac{d(d-4)}{d+1}I_{d-1}
        +
        \delta^{d-3}(d-3)J_0
        \right]
        >0,
\end{align}
where in the last line we used the recurrence $I_{d+1}=\frac{d}{d+1}I_{d-1}$ and $d\geq 4.$ Next, we choose \(\veps>0\) so small that
\(3C_d\veps<\delta_U\). Then, using \eqref{eq:angular-error-estimate-new}, we get
\begin{align}
        \int |\nabla U_{\veps,i}|_{g_\veps}^2\,\ud\nu_\veps \leq
        \calN_U^0+C_d\veps
        =
        2\calD_U^0-\delta_U+C_d\veps<
        2(\calD_U^0-C_d\veps)\leq  2\int U_{\veps,i}^2\,\ud\nu_\veps.
\end{align}
Hence,
\begin{align}\label{eq:angular-strict-new}
        \cR_\veps(U_{\veps,i})<2
        \qquad(i=1,\dots,d).
\end{align}
We next define another test function coming from the radial direction along the cylindrical end. Denote
\begin{align}
        b_A(r)=
        \begin{cases}
        \cos(r/\delta),&0\leq r\leq B,\\
        -A(r-B),&r\geq B,
        \end{cases}
\end{align}
where \(A>0\) is chosen so that \(\int b_A\,\ud\nu_0=0\). This can be done since
\begin{align}
        \int b_A\,\ud\nu_0
        &=
        \omega_{d-1}
        \int_0^B
        \cos(r/\delta)\bigl(\delta\sin(r/\delta)\bigr)^{d-1}\,\ud r
        -
        A\omega_{d-1}\delta^{d-1}
        \int_0^\infty
        \ta e^{-\ta^2/2}\,\ud\ta\\
        &=
        \omega_{d-1}
        \int_0^B
        \cos(r/\delta)\bigl(\delta\sin(r/\delta)\bigr)^{d-1}\,\ud r
        -
        A\omega_{d-1}\delta^{d-1}.
\end{align}
The first term is strictly positive, since \(\cos(r/\delta)>0\) and
\(\sin(r/\delta)>0\) for \(0<r<B\). Hence there is a unique choice
\begin{align}
        A
        =
        \frac{
        \int_0^B
        \cos(r/\delta)\bigl(\delta\sin(r/\delta)\bigr)^{d-1}\,\ud r
        }{
        \delta^{d-1}
        }
        >0
\end{align}
for which \(\int b_A\,\ud\nu_0=0\). Our test function will be $B_\veps:=b_A-c_\veps,$ where we choose $c_\veps:=Z_\veps^{-1} \int b_A\,\ud\nu_\veps$ so that \(\int B_\veps\,\ud\nu_\veps=0\). Although \(b_A\) is generally not \(C^1\) at \(r=B\), it is continuous and piecewise smooth. Hence its weak derivative is the piecewise derivative
\begin{align}
        b_A'(r)=
        \begin{cases}
        -\delta^{-1}\sin(r/\delta),&0< r<B,\\
        -A,&r>B.
        \end{cases}
\end{align}
The function $b_A$ has at most linear growth and $b_A'$ is bounded. Since
$\nu_\veps$ has finite mass and has a Gaussian tail on the cylindrical end,
$b_A,B_\veps\in W^{1,2}(\mathbb R^d,g_\veps,\nu_\veps)$. Thus $B_\veps$ is an admissible test function, and all appearances of $B_\veps'=b_A'$ below are understood in the weak sense. We first note that
\(c_\veps=O_d(\veps)\). This is because \(\int b_A\,\ud\nu_0=0\),
\eqref{eq:end-density-error}, and
\begin{align}
        &\left|\int b_A\,\ud\nu_\veps\right|
        +
        \left|Z_\veps-\int1\,\ud\nu_0\right|
        =
        \left|\int b_A\,\ud\nu_\veps-\int b_A\,\ud\nu_0\right|
        +
        \left|\int1\,\ud\nu_\veps-\int1\,\ud\nu_0\right| \notag\\
        &\leq
        \omega_{d-1}
        \int_0^\infty
        (1+|b_A(r)|)\left|e^{-V_\veps(r)}\rho_{\delta,\veps}(r)^{d-1}
        -
        e^{-V_0(r)}\rho_0(r)^{d-1}\right|\,\ud r \notag\\
        &\leq
        C_d\int_{I_\veps}1\,\ud r
        +
        C_d\int_{B+\veps}^\infty
        (1+r-B)\delta^{d-1}
        \left|e^{-V_\veps(r)}-e^{-V_0(r)}\right|
        \,\ud r \notag\\
        &\leq
        C_d\veps
        +
        C_d\veps
        \int_{B+\veps}^\infty
        (1+r-B)^2e^{-(r-B)^2/2}\,\ud r
        \leq
        C_d\veps.
\end{align}
Moreover, choosing $\veps\leq B/4$, we have that $B/2\leq B-2\veps$ and therefore $V_\veps\equiv 0$ on $[0,B/2]$ which implies
\begin{align}
        Z_\veps
        =
        \int_{\bbR^d}1\,\ud\nu_\veps
        &\geq
        \omega_{d-1}
        \int_0^{B/2}
        e^{-V_\veps(r)}
        \rho_{\delta,\veps}(r)^{d-1}\,\ud r=
        \omega_{d-1}
        \int_0^{B/2}
        \bigl(\delta\sin(r/\delta)\bigr)^{d-1}\,\ud r
        =
        c_d>0.
\end{align}
Thus,
\begin{align}\label{eqn:c_eps-bound}
        |c_\veps|
        =
        \left|
        \frac{\int b_A\,\ud\nu_\veps}{Z_\veps}
        \right|
        \leq
        \frac{C_d\veps}{c_d}
        \leq
        C_d\veps.
\end{align}
Denote
\begin{align}
        \calD_{B_\veps}^0
        =
        \int B_\veps^2\,\ud\nu_0,
        \qquad
        \calN_{B_\veps}^0
        =
        \int |B_\veps'|^2\,\ud\nu_0.
\end{align}
Since \(B_\veps=b_A-c_\veps\), \(\int b_A\,\ud\nu_0=0\), and
\(B_\veps'=b_A'\), we have
\begin{align}
        \calD_{B_\veps}^0
        =
        \int b_A^2\,\ud\nu_0
        +
        c_\veps^2\int1\,\ud\nu_0,\quad \calN_{B_\veps}^0
        =
        \int |b_A'|^2\,\ud\nu_0.
\end{align}
We next claim that
\[
        \delta_b
        :=
        2\int b_A^2\,\ud\nu_0
        -
        \int |b_A'|^2\,\ud\nu_0
        >0.
\]
First, the condition \(\int b_A\,\ud\nu_0=0\) gives the explicit value
\[
        A
        =
        \frac{
        \int_0^B
        \cos(r/\delta)(\delta\sin(r/\delta))^{d-1}\,\ud r
        }{\delta^{d-1}}
        =
        \frac{\delta}{d}.
\]
% Indeed, after the change of variables \(t=r/\delta\),
% \[
%         \int_0^B
%         \cos(r/\delta)(\delta\sin(r/\delta))^{d-1}\,\ud r
%         =
%         \delta^d\int_0^{\pi/2}\cos t\,\sin^{d-1}t\,\ud t
%         =
%         \frac{\delta^d}{d}.
% \]
Also, by integration by parts,
\begin{align}
        \int_0^\infty \tau^2e^{-\tau^2/2}\,\ud\tau
        =
        \int_0^\infty e^{-\tau^2/2}\,\ud\tau
        =J_0.
\end{align}
Now using the same notation as in \eqref{eqn:notation} we get
\begin{align}
        \delta_b
        &=
        \omega_{d-1}
        \left[
        2\delta^d\int_0^{\pi/2}\cos^2t\,\sin^{d-1}t\,\ud t
        +
        2\delta^{d-1}A^2J_0
        -
        \delta^{d-2}I_{d+1}
        -
        \delta^{d-1}A^2J_0
        \right] \\
        &=
        \omega_{d-1}
        \left[
        \delta^{d-2}
        \left(
        2(d-2)(I_{d-1}-I_{d+1})-I_{d+1}
        \right)
        +
        \delta^{d-1}A^2J_0
        \right] \\
        &=
        \omega_{d-1}
        \left[
        \delta^{d-2}\frac{(d-4)}{d+1}I_{d-1}
        +
        \frac{\delta^{d+1}}{d^2}J_0
        \right]
        >0.
\end{align}
Here we used \(I_{d+1}=\frac{d}{d+1}I_{d-1}\), \(\delta^2=d-2\), \(A=\delta/d\), and  \(d\geq4\). Then
\begin{align}
        \delta_{B_\veps}^0
        &:=
        2\calD_{B_\veps}^0-\calN_{B_\veps}^0=
        \delta_b+2c_\veps^2\int1\,\ud\nu_0
        \geq
        \delta_b>0.
\end{align}
We next compare the \(\nu_\veps\)-mass and Dirichlet energy of \(B_\veps\) to these limiting
quantities. Since \(|c_\veps|\leq C_d\veps\), on the cylindrical end
\begin{align}
        |B_\veps|^2+|B_\veps'|^2
        \leq C_d(1+(r-B)^2).
\end{align}
Thus, recalling that \(I_\veps=[B-2\veps,B+\veps]\) and using \eqref{eq:end-density-error}, we get
\begin{align}\label{eq:Beps-error-estimate}
        &\left|
        \int B_\veps^2\,\ud\nu_\veps-\calD_{B_\veps}^0
        \right|
        +
        \left|
        \int |B_\veps'|^2\,\ud\nu_\veps-\calN_{B_\veps}^0
        \right| \notag\\
        &\leq
        C_d\int_{I_\veps}1\,\ud r
        +
        C_d\int_B^\infty
        \bigl(1+(r-B)^2\bigr)
        \left|e^{-V_\veps(r)}-e^{-V_0(r)}\right|
        \,\ud r \notag\\
        &\leq
        C_d\veps
        +
        C_d\veps
        \int_B^\infty
        \bigl(1+(r-B)^2\bigr)(1+r-B)e^{-(r-B)^2/2}
        \,\ud r
        \leq
        C_d\veps.
\end{align}
Moreover,
\begin{align}
        \int b_A^2\,\ud\nu_0
        &\geq
        \omega_{d-1}
        \delta^d
        \int_0^{\pi/2}\cos^2t\,\sin^{d-1}t\,\ud t  
        =:m_d>0.
\end{align}
Choose \(\veps>0\) so small that \(C_d\veps<\frac12\calD_{B_\veps}^0\), which is possible since $\calD_{B_\veps}^0\geq \int b^2_A \ud\nu_0\geq m_d>0$, and
\(3C_d\veps<\delta_b\). Then, using \eqref{eq:Beps-error-estimate} and
\(\delta_{B_\veps}^0\geq\delta_b\), we obtain
\begin{align}
        \int |\nabla B_\veps|_{g_\veps}^2\,\ud\nu_\veps
        &=
        \int |B_\veps'|^2\,\ud\nu_\veps
        \leq
        \calN_{B_\veps}^0+C_d\veps\\
        &=
        2\calD_{B_\veps}^0-\delta_{B_\veps}^0+C_d\veps
        <
        2(\calD_{B_\veps}^0-C_d\veps)\\
        &\leq
        2\int B_\veps^2\,\ud\nu_\veps.
\end{align}
Hence
\begin{align}\label{eq:radial-strict-new}
        \cR_\veps(B_\veps)<2.
\end{align}
Let
\begin{align}
        F_\veps
        =
        \spn\{1,B_\veps,U_{\veps,1},\dots,U_{\veps,d}\}.
\end{align}
We show that the functions spanning $F_\veps$ are pairwise \(L^2(\nu_\veps)\)-orthogonal. Since
\begin{align}
        \ud\nu_\veps
        =
        e^{-V_\veps(r)}\rho_{\delta,\veps}(r)^{d-1}\,\ud r\,\dvol_{\bbS^{d-1}},
\end{align}
we have
\begin{align}
        \int_{\bbR^d} B_\veps\,\ud\nu_\veps
        &=0,\\
        \int_{\bbR^d} U_{\veps,i}\,\ud\nu_\veps
        &=
        \left(\int_0^\infty
        a_\veps(r)e^{-V_\veps(r)}\rho_{\delta,\veps}(r)^{d-1}\,\ud r
        \right)
        \left(\int_{\bbS^{d-1}}x_i\,\dvol_{\bbS^{d-1}}\right)
        =
        0,\\
        \int_{\bbR^d} B_\veps U_{\veps,i}\,\ud\nu_\veps
        &=
        \left(\int_0^\infty
        B_\veps(r)a_\veps(r)e^{-V_\veps(r)}
        \rho_{\delta,\veps}(r)^{d-1}\,\ud r
        \right)
        \left(\int_{\bbS^{d-1}}x_i\,\dvol_{\bbS^{d-1}}\right)
        =
        0,\\
        \int_{\bbR^d} U_{\veps,i}U_{\veps,j}\,\ud\nu_\veps
        &=
        \left(\int_0^\infty
        a_\veps(r)^2e^{-V_\veps(r)}
        \rho_{\delta,\veps}(r)^{d-1}\,\ud r
        \right)
        \left(\int_{\bbS^{d-1}}x_ix_j\,\dvol_{\bbS^{d-1}}\right)
        =
        0
        \quad(i\neq j).
\end{align}
Thus the functions in the spanning set are \(L^2(\nu_\veps)\)-orthogonal. In
particular they are linearly independent, and hence \(\dim F_\veps=d+2\). We next show orthogonality with respect to the Dirichlet energy inner product. To this end, we first compute
the gradients,
\begin{align}
        \nabla B_\veps=B_\veps'\partial_r,
        \qquad
        \nabla U_{\veps,i}
        =
        a_\veps'x_i\partial_r
        +
        \frac{a_\veps}{\rho_{\delta,\veps}^2}
        \nabla_{\bbS^{d-1}}x_i.
\end{align}
As a result, we get
\begin{align}
        \int_{\bbR^d}\la\nabla B_\veps,\nabla U_{\veps,i}\ra_{g_\veps}\,\ud\nu_\veps
        &=
        \left(\int_0^\infty
        B_\veps'(r)a_\veps'(r)e^{-V_\veps(r)}
        \rho_{\delta,\veps}(r)^{d-1}\,\ud r
        \right)
        \left(\int_{\bbS^{d-1}}x_i\,\dvol_{\bbS^{d-1}}\right)
        =
        0,
\end{align}
for all \(1\leq i\leq d\). Moreover,
\begin{align}
        \int_{\bbR^d}\la\nabla U_{\veps,i},\nabla U_{\veps,j}\ra_{g_\veps}\,\ud\nu_\veps
        &=
        \left(\int_0^\infty
        (a_\veps')^2e^{-V_\veps}
        \rho_{\delta,\veps}^{d-1}\,\ud r
        \right)
        \left(\int_{\bbS^{d-1}}x_ix_j\,\dvol_{\bbS^{d-1}}\right)\\
        &\quad+
        \left(\int_0^\infty
        a_\veps^2e^{-V_\veps}
        \rho_{\delta,\veps}^{d-3}\,\ud r
        \right)
        \left(\int_{\bbS^{d-1}}
        \la\nabla_{\bbS^{d-1}}x_i,\nabla_{\bbS^{d-1}}x_j\ra
        \,\dvol_{\bbS^{d-1}}\right)\\
        &=
        0,
\end{align}
for all \(1\leq i,j\leq d\), \(i\neq j\). Here the last equality follows because
\begin{align}
        \int_{\bbS^{d-1}}x_ix_j\,\dvol_{\bbS^{d-1}}=0
        \quad(i\neq j),
\end{align}
and, using \(-\Delta_{\bbS^{d-1}}x_j=(d-1)x_j\),
\begin{align}
        \int_{\bbS^{d-1}}
        \la\nabla_{\bbS^{d-1}}x_i,\nabla_{\bbS^{d-1}}x_j\ra
        \,\dvol_{\bbS^{d-1}}
        =
        (d-1)\int_{\bbS^{d-1}}x_ix_j\,\dvol_{\bbS^{d-1}}
        =
        0.
\end{align}
Thus the spanning functions are pairwise orthogonal both in \(L^2(\nu_\veps)\)
and with respect to the Dirichlet energy inner product. Then, for every nonzero \(u\in F_\veps\),
\begin{align}
        \cR_\veps(u)
        \leq
        \max\{\cR_\veps(B_\veps),
              \cR_\veps(U_{\veps,1}),\dots,\cR_\veps(U_{\veps,d})\}
        <2.
\end{align}
By the min-max characterization of the variational eigenvalues,
\begin{align}
        \lambda_{d+2}(\bbR^d,g_\veps,\mu_\veps)
        \leq
        \sup_{0\neq u\in F_\veps}\cR_\veps(u)
        <2.
\end{align}
Recall that \eqref{eqn:gaussian-spectrum} implies that
\(\lambda_{d+2}(\bbR^d,|\cdot|,\gamma^d)=2\). Thus
\begin{align}
        \lambda_{d+2}(\bbR^d,g_\veps,\mu_\veps)
        <
        \lambda_{d+2}(\bbR^d,|\cdot|,\gamma^d),
\end{align}
and the theorem follows.
\end{proof}
\begin{cor}\label{cor:gaussian}
Let \((\R^d,g,\mu)\) be the weighted manifold constructed in Theorem
\ref{thm:gaussian-main}, which satisfies \(\CD(1,\infty)\). Then there does not exist a $1$-Lipschitz map
\[
        T:(\R^d,|\cdot|,\gamma^d)\to (\R^d,g,\mu)
\]
pushing forward $\gamma^d$ onto $\mu$ up to a finite constant.
\end{cor}
\begin{proof}
If such a map existed, then Theorem~\ref{thm:contraction} would imply that for all $k\geq 1$ we have
\begin{align}
   \lam_k(\R^d,g,\mu) \geq  \lam_k(\R^d,|\cdot|,\gamma^d).
\end{align}
This contradicts \eqref{eqn:index-gaussian-K}.
\end{proof}

\section{Proof of Theorem~\ref{thm:gaussian-split-main}}
The proof is based on a volume-preserving perturbation of the round metric on
$\bbS^{d-1}$. We first use this perturbation to construct an explicit
finite-dimensional space of ordinary polynomials with small Rayleigh quotient
on a weighted metric cone. We then smooth the vertex of the cone by a
trace-free slow interpolation.  For a metric $g$ and a smooth function $V$, we write
\begin{align}
        \cR_{g,V}(u)
        :=
        \frac{\int |\nabla u|_g^2e^{-V}\dvol_g}
        {\int u^2e^{-V}\dvol_g}.
\end{align}
When $V=0$, we simply write $\cR_g(u)$. Let $g_{\can}$ denote the
canonical round metric on $\bbS^n$. For a metric $\bar{g}$ on $\bbS^n$ and a number
$s>0$, define
\begin{align}\label{eq:gaussian-split-Ahs}
        A(\bar{g},s)
        :=
        \max\left\{
        1,
        \sup_{\substack{\omega\in\bbS^n\\0\neq\xi\in T_\omega^*\bbS^n}}
        \frac{s^{-2}|\xi|_{\bar{g}}^2}{|\xi|_{g_\can}^2}
        \right\}.
\end{align}
This constant arises naturally when comparing the Dirichlet energy of functions written over cones in either the spherical metric or the perturbed metric $\bar{g}.$ For $\vartheta\in\bbR$, define the rotation
$R_\vartheta:\bbS^n\to\bbS^n$ by $R_\vartheta(z,y):=(e^{\iii\vartheta}z,y)$, where we use $\bbS^n\subset \bbC\times \R^{n-1}$ when $n\geq 4.$

We begin with the required perturbation of the spherical link.
\begin{lem}\label{lem:gaussian-split-link}
Let $n\geq4$. There exist an integer $q\geq2$, parameters $0<s<\rho<1,$ a smooth metric $\bar g$ on $\bbS^n$, and a smooth family of metrics
$\{g_\ta\}_{\ta\in[0,1]}$ on $\bbS^n$ satisfying $R_\vartheta^*g_\ta=g_\ta$ for every $\vartheta\in\bbR$ and $\ta\in[0,1]$ on $\bbS^n$ such that ${g}_0=g_{\can},$ ${g}_1=\bar{g}$ and for every $\ta\in [0,1]$ we have $\dvol_{{g}_\ta}=\dvol_{g_\can}$, $\operatorname{tr}_{g_\ta}(\partial_\ta g_\ta)=0$ and $\Ric_{{g}_\ta}\geq(n-1)\rho^2{g}_\ta$. Moreover, if
\begin{align}
        f_q(z,y):=\Re(z^q),
        \qquad
        (z,y)\in\bbS^n\subset\bbC\times\bbR^{n-1},
\end{align}
then
\begin{align}\label{eq:gaussian-split-link-Rayleigh}
                \cR_{\bar g}(f_q)<s^2q(q+n-1),
\end{align}
and
\begin{align}\label{eq:gaussian-split-link-comparison}
        A(\bar{g},s)(q-1)<q.
\end{align}
\end{lem}

\begin{proof}
Using the notation from Section~\ref{sec:sphere-thm-main}, we recall that 
\begin{align}
        g_\can 
        =
        \ud t^2+\sin^2t\,\ud\theta^2
        +\cos^2t\,g_{\bbS^{n-2}}.
\end{align}
For $0\leq \ta\leq1$ and $\veps>0$, define
\begin{align}\label{eq:gaussian-split-link-metric}
        g_\ta
        :=
        e^{-2 \ta \veps\sin^2t/(n-1)}
        \left(
        \ud t^2+\cos^2t\,g_{\bbS^{n-2}}
        \right)
        +
        e^{2 \ta \veps\sin^2t}\sin^2t\,\ud\theta^2.
\end{align}
Write $z=x_1+\iii x_2$, and set
\begin{align}
        \upsilon
        &:=
        |z|^2,
        \qquad
        \beta_0
        :=
        x_1\,\ud x_2-x_2\,\ud x_1.
\end{align}
Define
\begin{align}
        b_\ta(\upsilon)
        &:=
        \begin{cases}
        \displaystyle
        \frac{
        e^{2\ta\veps\upsilon}
        -
        e^{-\frac{2\ta\veps}{n-1}\upsilon}
        }{\upsilon},
        &
        \upsilon>0,
        \\[2mm]
        \displaystyle
        \frac{2n\ta\veps}{n-1},
        &
        \upsilon=0.
        \end{cases}
\end{align}
The function $b_\ta$ is smooth in $(\ta,\upsilon)$ and nonnegative.
Since
\begin{align}
        |z|^2
        =
        \sin^2t,
        \quad
        \beta_0
        =|z|^2\,\ud\theta
\end{align}
on the set $U=\{(z,y)\in \bS^n:z\neq 0,y\neq 0\}$ (equivalently when $0<t<\pi/2$), we may rewrite the metric on this set as
\begin{align}
        g_\ta
        =
        e^{-\frac{2\ta\veps}{n-1}|z|^2}g_{\can}
        +
        b_\ta(|z|^2)\beta_0\otimes\beta_0.
\end{align}
Hence $\{g_\ta\}_{\ta\in[0,1]}$ extends to a smooth family of
Riemannian metrics on all of $\bbS^n$. Since the coefficients in \eqref{eq:gaussian-split-link-metric} are independent of $\theta$, we have $R_\vartheta^*g_\ta=g_\ta$ for every $\vartheta\in\bbR$ and $\ta\in[0,1]$. Furthermore, $g_0= g_\can$ and
% We first verify that $g_\ta$ extends smoothly across the two singular orbits.
% Near $t=0$, introduce the radial arclength coordinate, 
% \begin{align}
%         r(t)
%         :=
%         \int_0^t
%         e^{-\tau \veps\sin^2\sigma/(n-1)}\,\ud\sigma.
% \end{align}
% Then
% \begin{align}
%         r=t+O(t^3),
% \end{align}
% and the warping function of the collapsing circle satisfies
% \begin{align}
%         e^{u\veps\sin^2t}\sin t=t+O(t^3).
% \end{align}
% Expressed as a function of $r$, this warping function is smooth and odd, with
% value zero and first derivative one at $r=0$. The remaining coefficients are
% smooth even functions of $r$. Thus the circle closes smoothly at $t=0$.
% Near $t=\pi/2$, the collapsing block
% \begin{align}
%         \ud t^2+\cos^2t\,g_{\bbS^{n-2}}
% \end{align}
% is multiplied by one smooth positive even function of $\pi/2-t$, while the
% coefficient of $\ud\theta^2$ is also a smooth even function. Thus the
% $\bbS^{n-2}$-factor closes smoothly at $t=\pi/2$.
this family of metrics is volume preserving since
\begin{align}
        \dvol_{g_\ta}
        =
        \left(
        e^{-\frac{2\ta\veps\sin^2t}{n-1}}
        \right)^{\frac{n-1}{2}}
        \left(
        e^{2\ta\veps\sin^2t}
        \right)^{\frac12}
        \dvol_{g_0}
        =
        e^{-\ta\veps\sin^2t}
        e^{\ta\veps\sin^2t}
        \dvol_{g_0}
        =
        \dvol_{g_0}.
        \label{eq:gaussian-split-volume-preserving}
\end{align}
Since
\begin{align}
        0=\partial_\tau \dvol_{g_\ta}=\frac{1}{2}\tr_{g_\ta} (\partial_\ta g_\ta) \dvol_{g_\ta}
        =
        \frac12\operatorname{tr}
        \bigl(g_\ta^{-1}\partial_\ta g_\ta\bigr)\dvol_{g_\ta},
\end{align}
we also see that $\tr(g_\ta^{-1}\partial_\ta g_\ta)=0.$ We next estimate the Ricci curvature of the metric $g_\ta$. Recall that on the
coordinate set under consideration, a point of $\bbS^n$ is written as
\begin{align}
        (z,y)
        =
        \bigl(e^{\iii\theta}\sin t,(\cos t)\omega\bigr),
        \qquad
        t\in(0,\pi/2),
        \quad
        \theta\in [0,2\pi),
        \quad
        \omega\in\bbS^{n-2}.
\end{align}
Thus $\theta$ is the angular coordinate on the $\bbS^1$-factor, while
$\omega$ denotes a point of the $\bbS^{n-2}$-factor.

Fix $\ta\in[0,1]$. We define the arclength coordinate in the $t$-direction by
\begin{align}\label{eq:gaussian-split-link-arclength}
        r(t)
        :=
        \int_0^t
        e^{-\frac{\ta\veps}{n-1}\sin^2\sigma}\,\ud\sigma.
\end{align}
Then
\begin{align}\label{eq:gaussian-split-r-t-derivative}
        \ud r 
        =
        e^{-\frac{\ta\veps}{n-1}\sin^2t}
        \ud t.
\end{align}
Define the functions appearing in the warped product by
\begin{align}\label{eq:gaussian-split-link-warping-functions}
        a_\ta(r)
        :=
        e^{\ta\veps\sin^2t}\sin t,\quad 
        c_\ta(r)
        :=
        e^{-\frac{\ta\veps}{n-1}\sin^2t}\cos t,
\end{align}
Then \eqref{eq:gaussian-split-link-metric} is simply $ g_\ta =\ud r^2+ a_\ta(r)^2\,\ud\theta^2+c_\ta(r)^2g_{\bbS^{n-2}}.$  Let $\{v_j\}_{j=1}^{n-2}$ be a local orthonormal frame with respect to the round metric on $\bbS^{n-2}$. Then
\begin{align}
        e_r
        :=
        \partial_r,\quad
        e_\theta
        :=
        \frac1{a_\ta}\partial_\theta,\quad 
        e_j
        :=
        \frac1{c_\ta}v_j,
        \qquad
        1\leq j\leq n-2,
\end{align}
form a local orthonormal frame with respect to the metric $g_\ta$. In these coordinates, we define the following components of the Ricci curvature as
\begin{align}\label{eq:gaussian-split-link-Ricci-eigenvalue-definitions}
        R_r
        :=
        \Ric_{g_\ta}(e_r,e_r),\quad 
        R_\theta
        :=
        \Ric_{g_\ta}(e_\theta,e_\theta),\quad 
        R_S
        :=
        \Ric_{g_\ta}(e_j,e_j),
        \quad
        1\leq j\leq n-2,
\end{align}
where we note that $R_S$ is independent of the index $1\leq j\leq n-2$. We now apply Lemma \ref{lem:multiply-warped-ricci} to estimate the Ricci curvature of each component in the above display. First, we record some derivatives of the warping functions. Denote $\partial_r a_\ta =a'_\ta$, $\partial_r c_\ta=c'_\ta$ and so on for the second derivative. Then, we have 
\begin{align}
        \frac{a_\ta'}{a_\ta}
        &=
        e^{\frac{\ta\veps}{n-1}\sin^2t}
        \left(
        \cot t
        +
        2\ta\veps\sin t\cos t
        \right),\quad 
        \frac{c_\ta'}{c_\ta}
        =
        -
        e^{\frac{\ta\veps}{n-1}\sin^2t}
        \left(
        \tan t
        +
        \frac{2\ta\veps}{n-1}\sin t\cos t
        \right), \label{eq:gaussian-split-link-first-log-derivatives}\\
        \frac{a_\ta'c_\ta'}{a_\ta c_\ta}
        &=
        -
        e^{\frac{2\ta\veps}{n-1}\sin^2t}
        \left(
        1
        +
        2\ta\veps\sin^2t
        +
        \frac{2\ta\veps}{n-1}\cos^2t
        +
        \frac{4(\ta\veps)^2}{n-1}
        \sin^2t\cos^2t
        \right).\label{eq:gaussian-split-link-cross-log-derivative}
\end{align}
By differentiating \eqref{eq:gaussian-split-link-first-log-derivatives} we obtain
\begin{align}\label{eq:gaussian-split-link-second-log-derivatives}
        \frac{a_\ta''}{a_\ta}
        &=
        e^{\frac{2\ta\veps}{n-1}\sin^2t}
        \left(
        -1
        +
        \frac{\ta\veps}{n-1}
        \left(
        6n-4-(8n-6)\sin^2t
        \right)
        +
        \frac{4n(\ta\veps)^2}{n-1}
        \sin^2t\cos^2t
        \right),\\
        \frac{c_\ta''}{c_\ta}
        &=
        e^{\frac{2\ta\veps}{n-1}\sin^2t}
        \left(
        -1
        +
        \frac{2\ta\veps}{n-1}
        \left(
        3\sin^2t-1
        \right)
        \right).
\end{align}
To estimate some numerical constants arising in \eqref{eq:gaussian-split-link-first-log-derivatives},\eqref{eq:gaussian-split-link-cross-log-derivative} and \eqref{eq:gaussian-split-link-second-log-derivatives} we define $C_n:=n-1+\frac{2n}{n-1}.$ Since $n\geq4$, we have $\frac{4n}{n-1}<C_n<2(n-1).$ We choose $\veps>0$ sufficiently small that
\begin{align}\label{eq:gaussian-split-epsilon-first}
        \frac{4n}{n-1}
        +
        \frac{n}{n-1}\veps
        &\leq
        C_n,
        \quad
        C_n\veps<n-1,
\end{align}
and set $\rho^2:= 1-\frac{C_n\veps}{n-1}.$  In particular, $0<\rho<1$. Since $0\leq\ta\leq1$, we have $0\leq\ta\veps\leq\veps$, and therefore
\begin{align}\label{eq:gaussian-split-Ricci-error-bound}
        \frac{4n}{n-1}\ta\veps
        +
        \frac{n}{n-1}(\ta\veps)^2
        =
        \ta\veps
        \left(
        \frac{4n}{n-1}
        +
        \frac{n}{n-1}\ta\veps
        \right)\leq
        \ta\veps
        \left(
        \frac{4n}{n-1}
        +
        \frac{n}{n-1}\veps
        \right)\leq
        C_n\ta\veps
        \leq
        C_n\veps.
\end{align}
Moreover, since $C_n>2$, we have $ 2\ta\veps \leq C_n\veps.$  Now applying Lemma~\ref{lem:multiply-warped-ricci} and
\eqref{eq:gaussian-split-link-second-log-derivatives}, we get
\begin{align}
        R_r
        &=
        -\frac{a_\ta''}{a_\ta}
        -(n-2)\frac{c_\ta''}{c_\ta}\\
        &=
        e^{\frac{2\ta\veps}{n-1}\sin^2t}
        \left(
        n-1
        +
        \frac{\ta\veps}{n-1}
        \left(
        (2n+6)\sin^2t-4n
        \right)
        -
        \frac{4n(\ta\veps)^2}{n-1}
        \sin^2t\cos^2t
        \right)\\
        &\geq
        e^{\frac{2\ta\veps}{n-1}\sin^2t}
        \left(
        n-1
        -
        \frac{4n}{n-1}\ta\veps
        -
        \frac{n}{n-1}(\ta\veps)^2
        \right)\\
        &\geq
        n-1-C_n\veps = (n-1)\rho^2.
        \label{eq:gaussian-split-link-Rr}
\end{align}
Similarly, using Lemma~\ref{lem:multiply-warped-ricci},
\eqref{eq:gaussian-split-link-second-log-derivatives}, and
\eqref{eq:gaussian-split-link-cross-log-derivative}, we obtain
\begin{align}
        R_\theta
        &=
        -\frac{a_\ta''}{a_\ta}
        -(n-2)\frac{a_\ta'c_\ta'}{a_\ta c_\ta}\\
        &=
        e^{\frac{2\ta\veps}{n-1}\sin^2t}
        \left(
        n-1
        +
        \frac{\ta\veps}{n-1}
        \left(
        2(n^2+1)\sin^2t-4n
        \right)
        -
        \frac{8(\ta\veps)^2}{n-1}
        \sin^2t\cos^2t
        \right)\\
        &\geq
        e^{\frac{2\ta\veps}{n-1}\sin^2t}
        \left(
        n-1
        -
        \frac{4n}{n-1}\ta\veps
        -
        \frac{2}{n-1}(\ta\veps)^2
        \right)\\
        &\geq
        e^{\frac{2\ta\veps}{n-1}\sin^2t}
        \left(
        n-1
        -
        \frac{4n}{n-1}\ta\veps
        -
        \frac{n}{n-1}(\ta\veps)^2
        \right)\\
        &\geq
        n-1-C_n\veps=
        (n-1)\rho^2.
        \label{eq:gaussian-split-link-Rtheta}
\end{align}
Finally, since $\Ric_{\bbS^{n-2}}= (n-3)g_{\bbS^{n-2}},$ Lemma~\ref{lem:multiply-warped-ricci},
\eqref{eq:gaussian-split-link-first-log-derivatives}, and
\eqref{eq:gaussian-split-link-cross-log-derivative} imply that
\begin{align}
        R_S
        &=
        \frac{n-3}{c_\ta^2}
        -
        \frac{c_\ta''}{c_\ta}
        -
        (n-3)
        \left(
        \frac{c_\ta'}{c_\ta}
        \right)^2
        -
        \frac{a_\ta'c_\ta'}{a_\ta c_\ta}\\
        &=
        e^{\frac{2\ta\veps}{n-1}\sin^2t}
        \left(
        n-1
        +
        \ta\veps
        \left(
        \frac4{n-1}-2\sin^2t
        \right)
        +
        \frac{8(\ta\veps)^2}{(n-1)^2}
        \sin^2t\cos^2t
        \right)\\
        &\geq
        e^{\frac{2\ta\veps}{n-1}\sin^2t}
        \left(
        n-1-2\ta\veps
        \right)\\
        &\geq
        n-1-C_n\veps= (n-1)\rho^2.
        \label{eq:gaussian-split-link-RS}
\end{align}
Since $g_\ta$ is smooth on $\bbS^n$, the preceding Ricci curvature
inequalities extend to the two singular orbits by continuity. Hence
\begin{align}
        \Ric_{g_\ta}
        \geq
        (n-1)\rho^2g_\ta
\end{align}
on all of $\bbS^n$. Note that when $\ta=0$, then $g_\ta=g_\can$ and the above expressions reduce to
\begin{align}
        R_r=R_\theta=R_S=n-1
\end{align}
as expected. 

Set $\bar g:=g_1$ and
$f_q(z,y):=\Re(z^q)=\sin^qt\cos(q\theta)$. For every $j\geq0$, define
\begin{align}
 Q_j:=\int_{\bbS^n}\sin^{2j}t\,\dvol_{g_\can},\quad D_n:=2+\frac{2\eee}{(n-1)^2}.
\end{align}
Since $\frac{C_n}{n-1}<2$, we may choose an integer $q\geq2$ sufficiently large such that
\[
        \eta_{n,q}
        :=
        2q^2
        \left(
        1-\frac1{2q+n+1}
        \right)
        -
        \frac{C_n}{n-1}q(q+n-1)
        >
        0.
\]
Fix this value of $q$. For every sufficiently small $\veps>0$ satisfying all
the preceding conditions, we may assume that
$D_nq^2\veps<\frac12\eta_{n,q}$. Define
\begin{align}\label{eq:gaussian-split-s-definition}
        s^2
        :=
        \rho^2
        -
        \frac{\eta_{n,q}}{4q(q+n-1)}\veps=
        1-
        \left(
        \frac{C_n}{n-1}
        +
        \frac{\eta_{n,q}}{4q(q+n-1)}
        \right)\veps.
\end{align}
After decreasing $\veps>0$ if necessary, the right-hand side is positive.
Moreover, since $\eta_{n,q}>0$, we have $0<s<\rho$.

We now estimate the Rayleigh quotient \eqref{eq:gaussian-split-link-Rayleigh}. To this end, first observe that 
\begin{align}\label{eqn:rayleigh-estimate-helper}
  \int_{\bbS^n}f_q^2\,\dvol_{\bar g} &= \frac12Q_q  \\
  \int_{\bbS^n}|\nabla f_q|_{\bar g}^2\,\dvol_{\bar g}
        &=
        \frac{q^2}{2}
        \int_{\bbS^n}
        \sin^{2q-2}t
        \left(
        \cos^2t\,
        e^{\frac{2\veps}{n-1}\sin^2t}
        +
        e^{-2\veps\sin^2t}
        \right)
        \dvol_{g_\can}\\
        &\leq  \frac{q^2}{2}\int_{\bbS^n} \sin^{2q-2} t \left(2-\sin^2t
        +
        2\veps\sin^2t
        \left(
        \frac{\cos^2t}{n-1}-1
        \right)
        +
        D_n\veps^2\sin^4t\right) \dvol_{g_\can}\\
        &= \frac{q^2}{2}
        \left[
        2Q_{q-1}-Q_q
        +
        2\veps
        \left(
        \frac{Q_q-Q_{q+1}}{n-1}-Q_q
        \right)
        +
        D_n\veps^2Q_{q+1}
        \right],
\end{align}
where in the inequality above we used the pointwise inequality
$\eee^y\leq1+y+\frac \eee 2y^2$ and $\eee^{-y}\leq1-y+\frac12y^2$ for $0\leq y\leq1$, together with
$0<\veps\leq1/2$ and $D_n=2+\frac{2\eee}{(n-1)^2}$. Also, in the last equality above we used
\[
        \partial_tf_q
        =
        q\sin^{q-1}t\cos t\cos(q\theta),
        \qquad
        \partial_\theta f_q
        =
        -q\sin^qt\sin(q\theta),
\]
and
\[
        \int_0^{2\pi}\cos^2(q\theta)\,\ud\theta
        =
        \int_0^{2\pi}\sin^2(q\theta)\,\ud\theta
        =
        \pi.
\]
Dividing the preceding numerator estimate by
$\frac12Q_q$ gives
\begin{align}\label{eqn:R-estimate}
        \cR_{\bar g}(f_q)
        &\leq
        \frac{q^2}{Q_q}
        \left[
        2Q_{q-1}-Q_q
        +
        2\veps
        \left(
        \frac{Q_q-Q_{q+1}}{n-1}-Q_q
        \right)
        +
        D_n\veps^2Q_{q+1}
        \right]
        \notag\\
        &\leq
        q(q+n-1)
        -
        2\veps q^2
        \left(
        1-\frac1{2q+n+1}
        \right)
        +
        D_nq^2\veps^2\\
        &=
        \rho^2q(q+n-1)
        -
        \eta_{n,q}\veps
        +
        D_nq^2\veps^2\\
        &<
        \rho^2q(q+n-1)
        -
        \frac12\eta_{n,q}\veps\\
        &<
        \rho^2q(q+n-1)
        -
        \frac14\eta_{n,q}\veps=
        s^2q(q+n-1),
\end{align}


For the third line in the Rayleigh quotient estimate, we used
\begin{align}
        2Q_{q-1}-Q_q
        =
        \frac{q+n-1}{q}Q_q, \quad 
        \frac{Q_q-Q_{q+1}}{(n-1)Q_q}
        =
        \frac1{2q+n+1},\quad 
        Q_{q+1}
        \leq
        Q_q.
\end{align}
The first two identities in the above display follow from
$\frac{Q_j}{Q_{j-1}}=\frac{2j}{2j+n-1}$, while the last inequality in the above display follows
from $0\leq\sin^2t\leq1$. We next estimate
\begin{align}
        A(\bar g,s)(q-1)
        &=
        \max\left\{
        1,
        \sup_{\substack{\omega\in\bbS^n\\
        0\neq\xi\in T_\omega^*\bbS^n}}
        \frac{s^{-2}|\xi|_{\bar g}^2}
        {|\xi|_{g_\can}^2}
        \right\}
        (q-1)
        \leq
        \max\left\{
        1,
        s^{-2}e^{2\veps/(n-1)}
        \right\}
        (q-1)
        <
        q.
\end{align}
To justify the preceding estimate, write
\begin{align}
        \xi
        =
        \xi_t\,\ud t+\xi_\theta\,\ud\theta+\xi_S,
        \qquad
        \xi_S\in T_\omega^*\bS^{n-2}.
\end{align}
Then since  $\bar g
        =
        e^{-\frac{2\veps}{n-1}\sin^2t}
        \left(
        \ud t^2+\cos^2t\,g_{\bbS^{n-2}}
        \right)
        +
        e^{2\veps\sin^2t}
        \sin^2t\,\ud\theta^2,$
we have
\begin{align}
        |\xi|_{\bar g}^2
        =
        e^{\frac{2\veps}{n-1}\sin^2t}
        \left(
        \xi_t^2+
        \frac{|\xi_S|_{g_{\bbS^{n-2}}}^2}{\cos^2t}
        \right)
        +
        e^{-2\veps\sin^2t}
        \frac{\xi_\theta^2}{\sin^2t}
        &\leq
        e^{2\veps/(n-1)}
        \left(
        \xi_t^2+
        \frac{|\xi_S|_{g_{\bbS^{n-2}}}^2}{\cos^2t}
        +
        \frac{\xi_\theta^2}{\sin^2t}
        \right)
        \notag\\
        &=
        e^{2\veps/(n-1)}
        |\xi|_{g_\can}^2.
\end{align}
By \eqref{eq:gaussian-split-s-definition}, $s\to1$ as
$\veps\to 0 $. Therefore, as $\veps \to 0$
\begin{align}
        \max\left\{
        1,
        s^{-2}e^{2\veps/(n-1)}
        \right\}
        (q-1)
        \rightarrow
        q-1
        <
        q,
\end{align}
which proves
\begin{align}
        A(\bar g,s)(q-1)
        <
        q
\end{align}
for all sufficiently small $\veps>0$.
\end{proof}
Now, we will use the metric $\bar{g}$ constructed in the previous lemma and construct a cone over that metric with a suitable potential. The key point is to obtain a new weighted space with a suitable spectral property that will be needed for our counterexample. 
\begin{lem}\label{lem:gaussian-split-polynomial-cone}
Let $n\geq1$ and $q\geq2$. Let $g_0=g_\can$ be the canonical round metric
on $\bbS^n$, and let $\bar g$ be a metric on $\bbS^n$ satisfying
\begin{align}
        R_\vartheta^*\bar g
        =
        \bar g
        \qquad
        \text{for every }\vartheta\in\bbR,
        \qquad
        \dvol_{\bar g}
        =
        \dvol_{g_0}.
\end{align} 
Suppose that $s>0$ satisfies
\begin{align}\label{eq:gaussian-split-polynomial-assumptions}
        A(\bar g,s)(q-1)<q,
        \qquad
        \cR_{\bar g}(f_q)<s^2q(q+n-1),
\end{align}
where $f_q(z,y)=\Re(z^q)$. Then there exist $\alpha>1$ and a vector space
$E_C$ of ordinary polynomials on $\bbR^{n+1}$ such that, upon identifying
$\bbR^{n+1}\setminus\{0\}$ with $(0,\infty)\times\bbS^n$ by $x=r\omega$, we
have
\begin{align}\label{eq:gaussian-split-polynomial-dimension}
        \dim E_C=N_q+1,
        \qquad
        N_q:=\binom{n+q}{n+1},
\end{align}
and
\begin{align}\label{eq:gaussian-split-polynomial-Rayleigh}
        \sup_{0\neq p\in E_C}
        \cR_{g_C,V_C}(p)<q,
\end{align}
where
\begin{align}\label{eq:gaussian-split-polynomial-cone}
        C=(0,\infty)\times\bbS^n,
        \qquad
        g_C=\ud r^2+s^2r^2\bar g,
        \qquad
        V_C=\frac{\alpha r^2}{2}.
\end{align}
Moreover,
\begin{align}\label{eq:gaussian-split-Gaussian-index}
        \lam_{N_q+1}
        (\bbR^{n+1},|\cdot|,\gamma^{n+1})=q.
\end{align}
\end{lem}

\begin{proof}
Let
\begin{align}
        \cP_{<q}
        :=
        \left\{
        p:\bbR^{n+1}\to\bbR:
        p\text{ is a polynomial of degree strictly less than }q
        \right\}.
\end{align}
Since the number of monomials in $n+1$ variables of degree $j$ is
$\binom{n+j}{j}$, we have
\begin{align}\label{eq:gaussian-split-polynomial-count}
        \dim\cP_{<q}
        =
        \sum_{j=0}^{q-1}\binom{n+j}{j}
        =
        \binom{n+q}{n+1}
        =
        N_q.
\end{align}

We first estimate the Gaussian Rayleigh quotient on $\cP_{<q}$. Let
$L_0:=-\Delta+\la x,\nabla\cdot\ra$ denote the Ornstein--Uhlenbeck operator. The multivariable Hermite polynomials
\begin{align}
        H_\beta(x)
        :=
        \prod_{j=1}^{n+1}H_{\beta_j}(x_j),
        \qquad
        \beta\in\bbN_0^{n+1},
\end{align}
form an orthogonal basis of $L^2(\gamma^{n+1})$ and satisfy
$L_0H_\beta=|\beta|H_\beta$ (cf. Section~2.1 in
\cite{milman2018spectral}). Moreover, the Hermite polynomials with
$|\beta|<q$ span $\cP_{<q}$. Therefore, if
$p=\sum_{|\beta|<q}c_\beta H_\beta$, then
\begin{align}
        \int_{\bbR^{n+1}}|\nabla p|^2\,\ud\gamma^{n+1}
        =
        \int_{\bbR^{n+1}}pL_0p\,\ud\gamma^{n+1}
        &=
        \sum_{|\beta|<q}
        |\beta||c_\beta|^2
        \|H_\beta\|_{L^2(\gamma^{n+1})}^2\\
        &\leq
        (q-1)
        \sum_{|\beta|<q}
        |c_\beta|^2
        \|H_\beta\|_{L^2(\gamma^{n+1})}^2\\
        &= 
        (q-1)
        \int_{\bbR^{n+1}}p^2\,\ud\gamma^{n+1}.
\end{align}
Consequently, for every nonzero polynomial $p\in \calP_{<q}$ the Rayleigh quotient satisfies
\begin{align}\label{eq:gaussian-split-lower-polynomial-Gaussian}
        \frac{
        \int_{\bbR^{n+1}}|\nabla p|^2\,\ud\gamma^{n+1}
        }{
        \int_{\bbR^{n+1}}p^2\,\ud\gamma^{n+1}
        }
        \leq
        q-1.
\end{align}
Since the multiplicity of the Gaussian eigenvalue $j$ is
$\binom{n+j}{j}$, exactly $\sum_{j=0}^{q-1}\binom{n+j}{j}=N_q$ eigenvalues, including the zero
eigenvalue and counting multiplicity, lie strictly below $q$. This proves
\eqref{eq:gaussian-split-Gaussian-index}.

We next estimate the Rayleigh quotient of the same polynomials on the cone
with potential $r^2/2$. For every $0\neq p\in\cP_{<q}$, we have
\begin{align}\label{eq:gaussian-split-lower-polynomial-cone-one}
        \cR_{g_C,r^2/2}(p)
        \leq
        A(\bar g,s)
        \frac{
        \int_{\bbR^{n+1}}|\nabla p|^2\,\ud\gamma^{n+1}
        }{
        \int_{\bbR^{n+1}}p^2\,\ud\gamma^{n+1}
        }
        \leq
        A(\bar g,s)(q-1)
        <
        q,
\end{align}
where for the first inequality in the above display, we used
\begin{align}
        e^{-r^2/2}\dvol_{g_C}
        =
        s^ne^{-r^2/2}r^n\,\ud r\,\dvol_{\bar g}
        =
        s^ne^{-r^2/2}r^n\,\ud r\,\dvol_{g_0}
        =
        s^n(2\pi)^{\frac{n+1}{2}}\,\ud\gamma^{n+1},
\end{align}
and
\begin{align}
        |\nabla p|_{g_C}^2
        =
        |\partial_rp|^2
        +
        \frac1{s^2r^2}|\nabla_{\bbS^n}p|_{\bar g}^2
        \leq
        A(\bar g,s)
        \left(
        |\partial_rp|^2
        +
        \frac1{r^2}|\nabla_{\bbS^n}p|_{g_0}^2
        \right)
        =
        A(\bar g,s)|\nabla p|^2,
\end{align}
where the final norm is the Euclidean norm written in polar coordinates. The second inequality in \eqref{eq:gaussian-split-lower-polynomial-cone-one} follows from \eqref{eq:gaussian-split-lower-polynomial-Gaussian}, while the last inequality follows from \eqref{eq:gaussian-split-polynomial-assumptions}. 

Now observe that after rescaling the polynomials in $\cP_{<q}$ we obtain
\begin{align}\label{eq:gaussian-split-polynomial-alpha-scaling}
        \sup_{0\neq p\in\cP_{<q}}
        \cR_{g_C,\alpha r^2/2}(p)
        =
        \alpha
        \sup_{0\neq p\in\cP_{<q}}
        \cR_{g_C,r^2/2}(p).
\end{align}
Therefore, the estimate \eqref{eq:gaussian-split-lower-polynomial-cone-one} implies
\begin{align}\label{eq:gaussian-split-lower-polynomial-cone-alpha}
        \sup_{0\neq p\in\cP_{<q}}
        \cR_{g_C,\alpha r^2/2}(p)
        \leq
        \alpha A(\bar g,s)(q-1),
\end{align}
for any $\alp>0.$ Now we aim to construct the vector space of polynomials $E_C$ as stated in the lemma. This space will consist of $\cP_{<q}$ and the following polynomial written in polar coordinates
\begin{align}\label{eq:gaussian-split-Pq-polar}
        P_q(r,\omega)=r^qf_q(\omega),
\end{align}
where $f_q(z,y):=\operatorname{Re}(z^q)$ for $(z,y)\in \bbS^{n}\subset \bbC\times \bbR^{n-1}$. We will now estimate the Rayleigh quotient of $P_q$. To this end, define for $m>1$
\begin{align}
        I_m(\alpha)
        :=
        \int_0^\infty r^me^{-\alpha r^2/2}\,\ud r.
\end{align}
Integration by parts implies
\begin{align}\label{eq:gaussian-split-Im-recurrence}
        I_m(\alpha)
        =
        \frac{m-1}{\alpha}I_{m-2}(\alpha)
        \qquad
        \text{for every }m>1.
\end{align}
Since
\begin{align}
        |\nabla P_q|_{g_C}^2
        =
        r^{2q-2}
        \left(
        q^2f_q^2+s^{-2}|\nabla f_q|_{\bar g}^2
        \right),
\end{align}
we obtain
\begin{align}\label{eq:gaussian-split-Pq-Rayleigh}
        \cR_{g_C,\alpha r^2/2}(P_q)
        &=
        \frac{
        \left(
        q^2\int_{\bbS^n}f_q^2\dvol_{\bar g}
        +s^{-2}\int_{\bbS^n}|\nabla f_q|_{\bar g}^2\dvol_{\bar g}
        \right)
        I_{n+2q-2}(\alpha)
        }{
        \left(\int_{\bbS^n}f_q^2\dvol_{\bar g}\right)
        I_{n+2q}(\alpha)
        }
        \notag\\
        &=
        \frac{\alpha}{n+2q-1}
        \left(
        q^2+s^{-2}\cR_{\bar g}(f_q)
        \right),
\end{align}
where in the last equality  we used
$I_{n+2q}(\alpha)=\frac{n+2q-1}{\alpha}I_{n+2q-2}(\alpha)$.
By \eqref{eq:gaussian-split-polynomial-assumptions},
\begin{align}
        q^2+s^{-2}\cR_{\bar g}(f_q)
        <
        q^2+q(q+n-1)
        =
        q(n+2q-1).
\end{align}
Therefore, if
\begin{align}\label{eq:gaussian-split-Bq}
        B_q
        :=
        \frac{q^2+s^{-2}\cR_{\bar g}(f_q)}{n+2q-1},
\end{align}
then
\begin{align}
        0<B_q<q,
        \qquad
        \cR_{g_C,\alpha r^2/2}(P_q)=\alpha B_q.
\end{align}

Since $q\geq2$, the quantities
$q/[A(\bar g,s)(q-1)]$ and $q/B_q$ are well-defined and strictly larger than
one. We may therefore choose
\begin{align}\label{eq:gaussian-split-alpha-choice}
        1<\alpha<
        \min\left\{
        \frac{q}{A(\bar g,s)(q-1)},
        \frac{q}{B_q}
        \right\}.
\end{align}
For this choice of $\alpha$, we have
\begin{align}\label{eq:gaussian-split-two-polynomial-bounds}
        \sup_{0\neq p\in\cP_{<q}}
        \cR_{g_C,V_C}(p)
        &\leq
        \alpha A(\bar g,s)(q-1)
        <q,\notag\\
        \cR_{g_C,V_C}(P_q)
        &=
        \alpha B_q
        <q.
\end{align}
It remains to show that $P_q$ is orthogonal to $\cP_{<q}$ with respect to the Gaussian-weighted $L^2$ and Dirichlet inner products on the cone. For $u,v \in W^{1,2} \left(C,g_C,e^{-V_C}\dvol_{g_C}\right)$, define
\begin{align}
        \cB_0(u,v)
        :=
        \int_Cuv\,e^{-V_C}\dvol_{g_C},
        \quad 
        \cB_1(u,v)
        :=
        \int_C
        \la\nabla u,\nabla v\ra_{g_C}
        e^{-V_C}\dvol_{g_C}.
\end{align} 
Then, observe that the transformation $R_\vartheta(r,z,y):=(r,e^{\iii\vartheta}z,y)$ preserves the metric $g_C$ and the potential $V_C$. Then arguing as in \eqref{eqn:change-of-variables} earlier, we obtain that
\begin{align}\label{eq:gaussian-split-polynomial-orthogonality}
        \int_CP_qp\,e^{-V_C}\dvol_{g_C}
        =0,\quad 
        \int_C\la\nabla P_q,\nabla p\ra_{g_C}
        e^{-V_C}\dvol_{g_C}
        =0,
\end{align}
for every $p\in \cP_{<q}.$ Set $E_C :=\cP_{<q}\oplus\spn\{P_q\}.$ By \eqref{eq:gaussian-split-polynomial-count}, $\dim E_C=N_q+1$, which proves
\eqref{eq:gaussian-split-polynomial-dimension}. Finally, if
$0\neq u=p+aP_q\in E_C$, where $p\in\cP_{<q}$ and $a\in\bbR$, then
\begin{align}
        \cR_{g_C,V_C}(u)
        &=
        \frac{
        \int_C|\nabla p|_{g_C}^2e^{-V_C}\dvol_{g_C}
        +a^2\int_C|\nabla P_q|_{g_C}^2e^{-V_C}\dvol_{g_C}
        }{
        \int_Cp^2e^{-V_C}\dvol_{g_C}
        +a^2\int_CP_q^2e^{-V_C}\dvol_{g_C}
        }
        \notag\\
        &\leq
        \max\left\{
        \sup_{0\neq p\in\cP_{<q}}\cR_{g_C,V_C}(p),
        \cR_{g_C,V_C}(P_q)
        \right\}
        <q.
\end{align}
Here the first equality follows from
\eqref{eq:gaussian-split-polynomial-orthogonality}, and the last inequality
follows from \eqref{eq:gaussian-split-two-polynomial-bounds}. This proves
\eqref{eq:gaussian-split-polynomial-Rayleigh} and completes the proof.
\end{proof}

Unfortunately, the cone constructed above has a singularity at the origin. We now prove a lemma that will smooth this singularity without losing the curvature control on the metric and the potential and the spectral inequalities needed to prove Theorem~\ref{thm:gaussian-split-main}. We first recall a technical proposition, based on
\cite[Proposition~4.1]{barGauduchonMoroianu2005}, which will be useful for
computing the Ricci curvature of the smoothed metric.

\begin{prop}\label{prop:generalized-cylinder-curvature}
Let $M^n$ be a smooth manifold, let $I\subset\bbR$ be an interval, and let
$\{h_r\}_{r\in I}$ be a smooth family of Riemannian metrics on $M$. Set
\begin{align}
        Z
        :=
        I\times M,
        \qquad
        g^Z
        :=
        \ud r^2+h_r.
\end{align}
For $r\in I$, let $M_r:=\{r\}\times M$, let $\nu:=\partial_r$, and let
$W_r$ denote the Weingarten map of $M_r$ with respect to $\nu$, so that
\begin{align}\label{eq:generalized-cylinder-Weingarten}
        \la W_r(X),Y\ra_{h_r}
        =
        -\frac12\partial_rh_r(X,Y)
\end{align}
for $X,Y$ tangent to $M_r$. Then
\begin{align}\label{eq:generalized-cylinder-Ricci-radial}
        \Ric_{g^Z}(\partial_r,\partial_r)
        =
        \partial_r\Tr(W_r)-\Tr(W_r^2),
\end{align}
\begin{align}\label{eq:generalized-cylinder-Ricci-mixed}
        \Ric_{g^Z}(\partial_r,X)
        =
        X\bigl(\Tr(W_r)\bigr)
        -
        \la\div_{h_r}W_r,X\ra_{h_r},
\end{align}
and
\begin{align}\label{eq:generalized-cylinder-Ricci-tangential}
        \Ric_{g^Z}(X,Y)
        =
        \Ric_{h_r}(X,Y)
        +
        \la
        \left(
        \partial_rW_r-\Tr(W_r)W_r
        \right)X,
        Y
        \ra_{h_r}.
\end{align}
Here $\partial_rW_r$ is computed after identifying the tangent bundles of
the slices through the product structure.
\end{prop}

\begin{proof}
These identities follow from Proposition~4.1 in \cite{barGauduchonMoroianu2005}. To see this, first using
\eqref{eq:generalized-cylinder-Weingarten} we obtain
\begin{align}
        \partial_r\Tr(W_r)
        =
        2\Tr(W_r^2)
        -
        \frac12\Tr_{h_r}(\partial_r^2h_r),
\end{align}
which transforms equation~(18) in Proposition~4.1 in \cite{barGauduchonMoroianu2005} into
\eqref{eq:generalized-cylinder-Ricci-radial}. Moreover, differentiating \eqref{eq:generalized-cylinder-Weingarten} gives
\begin{align}
        \frac12\partial_r^2h_r(X,Y)
        =
        2\la W_r(X),W_r(Y)\ra_{h_r}
        -
        \la(\partial_rW_r)(X),Y\ra_{h_r}.
\end{align}
Substituting this identity into equation~(20) in Proposition~4.1 in \cite{barGauduchonMoroianu2005} gives \eqref{eq:generalized-cylinder-Ricci-tangential}, while equation~(19) in Proposition~4.1 in \cite{barGauduchonMoroianu2005} gives \eqref{eq:generalized-cylinder-Ricci-mixed}.
\end{proof}



\begin{lem}\label{lem:gaussian-split-smoothing}
Let $n,q,s,\rho,\bar g$, and $\{g_\ta\}_{\ta\in[0,1]}$ be the objects
constructed in Lemma~\ref{lem:gaussian-split-link}, and let
$\alpha$, $(C,g_C,V_C)$, $E_C$, and $N_q$ be the objects constructed in
Lemma~\ref{lem:gaussian-split-polynomial-cone}. Then, for every sufficiently
small $\eta>0$, there exist a smooth complete metric $g_\eta$ on
$\bbR^{n+1}$, a smooth function $V_\eta:\bbR^{n+1}\to\bbR$, and a
finite-dimensional vector space
\begin{align}
        E_\eta
        \subset
        C^\infty(\bbR^{n+1})
        \cap
        W^{1,2}
        \left(
        \bbR^{n+1},g_\eta,e^{-V_\eta}\dvol_{g_\eta}
        \right)
\end{align}
such that
\begin{align}\label{eq:gaussian-split-smoothed-curvature}
        \Ric_{g_\eta}\geq0,
        \quad
        \nabla_{g_\eta}^2V_\eta\geq g_\eta,
\end{align}
$\dim E_\eta=N_q+1,$ and
\begin{align}\label{eq:gaussian-split-smoothed-Rayleigh}
        \sup_{0\neq f\in E_\eta}
        \cR_{g_\eta,V_\eta}(f)
        <
        q.
\end{align}
Moreover,
\begin{align}\label{eq:gaussian-split-smoothed-finite-mass}
        \int_{\bbR^{n+1}}e^{-V_\eta}\dvol_{g_\eta}<\infty.
\end{align}
% and the embedding
% \begin{align}\label{eq:gaussian-split-smoothed-compact-embedding}
%         W^{1,2}
%         \left(
%         \bbR^{n+1},g_\eta,e^{-V_\eta}\dvol_{g_\eta}
%         \right)
%         \hookrightarrow
%         L^2
%         \left(
%         \bbR^{n+1},e^{-V_\eta}\dvol_{g_\eta}
%         \right)
% \end{align}
% is compact. In particular, the corresponding weighted Laplacian has compact
% resolvent.
\end{lem}

\begin{proof}
For notational convenience, set $E:=E_C$,  $\Lambda_E:=\sup_{0\neq p\in E_C} \cR_{g_C,V_C}(p)$ and $\delta_*:=\rho-s>0$. By \eqref{eq:gaussian-split-polynomial-Rayleigh}, we have $\Lambda_E<q.$ Set $P_\ta := g_\ta^{-1}\partial_\ta g_\ta.$ By Lemma~\ref{lem:gaussian-split-link}, for every $\ta\in[0,1]$ we have
\begin{align}\label{eq:gaussian-split-smoothing-Ricci-hypothesis}
        \Ric_{g_\ta}
        \geq
        (n-1)\rho^2g_\ta,\text{ and } \Tr(P_\ta)=0.
\end{align}
Our argument will proceed in five steps.
\paragraph{$\mathrm{(i)}$ \textit{Construction of a smooth metric}} Choose smooth functions $\psi:\bbR\to[s,1]$ and $\chi:\bbR\to[0,1]$ such that
\begin{align}
        &\psi(t)=1
        \text{ for }t\leq0,
        \quad
        \psi(t)=s
        \text{ for }t\geq4,
        \quad
        \psi'\leq0,\quad 
        \psi(t)
        \leq
        s+\frac{\delta_*}{4}
        \text{ for }t\geq2, \qquad \label{eq:gaussian-split-profile-small}\\
        &\chi(t)=0 \text{ for }t\leq2,\quad 
        \chi(t)=1 \text{ for }t\geq3,\quad 
        \chi'(t)>0 \text{ for }2<t<3,\quad 
        -\psi'(t)
        \geq
        c_0\delta_*
        \text{ for }2\leq t\leq 3 \qquad \label{eq:gaussian-split-profile-negative}
\end{align}
for some constant $c_0>0$.  For $T\geq1$, define $\psi_T(t):=\psi(t/T)$ and $\chi_T(t):=\chi(t/T).$ For $r>0$, set
\begin{align}\label{eq:gaussian-split-phi-u-definition}
        \phi(r)
        :=
        \int_0^r \psi_T(\log\sigma)\,\ud\sigma,
        \qquad
        u(r):=\chi_T(\log r).
\end{align}
Since $\psi_T(\log\sigma)=1$ and $\chi_T(\log\sigma)=0$ for all sufficiently small $\sigma>0$, we have $\phi(r)=r$ and $u(r)=0$ near $r=0$. Moreover,
\begin{align}\label{eq:gaussian-split-phi-bounds}
        s\leq\phi'\leq1,
        \qquad
        \phi''\leq0,
        \qquad
        sr\leq\phi(r)\leq r.
\end{align}
There exists $b\geq0$ such that, for every $r\geq e^{4T}$,
\begin{align}\label{eq:gaussian-split-phi-end}
        u(r)=1,
        \qquad
        \phi(r)=s(r+b).
\end{align}
Define $ \widetilde g :=\ud r^2+\phi(r)^2g_{u(r)}.$ Since $\phi(r)=r$ and $g_{u(r)}=g_0=g_\can$ near the origin, the metric
$\widetilde g$ extends smoothly across the origin. Moreover,
$\widetilde g\geq\ud r^2$, and hence $\widetilde g$ is complete.

\paragraph{$\mathrm{(ii)}$ \textit{Ricci curvature lower bound}} Set
$h_r:=\phi(r)^2g_{u(r)}$, so that
$\widetilde g=\ud r^2+h_r$, and set
\begin{align}
        F_r
        :=
        \frac12u'(r)P_{u(r)}.
\end{align}
Since $\partial_\ta g_\ta$ is symmetric, $P_\ta$ is self-adjoint with
respect to $g_\ta$. Consequently, $F_r$ is self-adjoint with respect to
$h_r$. By \eqref{eq:generalized-cylinder-Weingarten}, the Weingarten map of the slice $M_r=\{r\}\times\bbS^n$ satisfies
\begin{align}\label{eq:gaussian-split-shape-operator}
        -W_r
        =
        \frac{\phi'}{\phi}\Id+F_r.
\end{align}
In particular, if
$\mathrm{II}(X,Y):=-h_r(W_r(X),Y)$, then
\begin{align}
        \mathrm{II}(X,Y)
        =
        h_r
        \left(
        \left(
        \frac{\phi'}{\phi}\Id+F_r
        \right)X,Y
        \right).
\end{align}
Since $\Tr(P_\ta)=0$, we have
$\Tr(F_r)=0$ and
$\Tr(W_r)=-n\phi'/\phi$. Therefore,
Proposition~\ref{prop:generalized-cylinder-curvature} gives
\begin{align}\label{eq:gaussian-split-Ricci-rr-exact}
        \Ric_{\widetilde g}(\partial_r,\partial_r)
        &=
        -n\frac{\phi''}{\phi}
        -
        \Tr(F_r^2),\notag\\
        \Ric_{\widetilde g}(\partial_r,X)
        &=
        \la\div_{h_r}F_r,X\ra_{h_r},\notag\\
        \Ric_{\widetilde g}(X,Y)
        &=
        \Ric_{g_{u(r)}}(X,Y)
        -
        \left(
        \frac{\phi''}{\phi}
        +(n-1)\frac{(\phi')^2}{\phi^2}
        \right)h_r(X,Y)
        -
        h_r
        \left(
        \left(
        \partial_rF_r
        +
        n\frac{\phi'}{\phi}F_r
        \right)X,Y
        \right)
\end{align}
for vectors $X,Y$ tangent to $M_r$. Recall that $\delta_*:=\rho-s>0$ and set
\begin{align}
        M
        :=
        1+
        \sup_{\ta\in[0,1]}
        \left(
        |P_\ta|_{C^1(g_\ta)}
        +
        |\partial_\ta P_\ta|_{C^0(g_\ta)}
        \right).
\end{align}
Let $\cA_T:=\supp u'$. In the remainder of this step, $c>0$ and
$C<\infty$ denote constants independent of $T$. Since
$u(r)=\chi((\log r)/T)$, we have
\begin{align}
        |u'(r)|
        &\leq
        \frac{C}{Tr},
        \qquad
        |u''(r)|
        \leq
        \frac{C}{Tr^2}.
\end{align}
Moreover,
\begin{align}
        \partial_rF_r
        &=
        \frac12u''(r)P_{u(r)}
        +
        \frac12(u'(r))^2
        \left.
        \partial_\ta P_\ta
        \right|_{\ta=u(r)},
        \notag\\
        \div_{h_r}F_r
        &=
        \frac12u'(r)
        \div_{h_r}P_{u(r)}.
\end{align}
Using these identities, $\phi\leq r$, and the fact that
$h_r=\phi^2g_{u(r)}$ is a constant rescaling on each slice, we obtain
\begin{align}\label{eq:gaussian-split-F-bound}
        |F_r|_{h_r}
        \leq
        \frac{CM}{T\phi},\quad 
        |\div_{h_r}F_r|_{h_r}
        \leq
        \frac{CM}{T\phi^2},\quad 
        \left|
        \partial_rF_r + n\frac{\phi'}{\phi}F_r
        \right|_{h_r}
        \leq
        \frac{CM}{T\phi^2}.
\end{align}
On $\cA_T$, we have $(\log r)/T\in[2,3]$. Therefore,
\eqref{eq:gaussian-split-profile-negative},
\eqref{eq:gaussian-split-profile-small}, and
$sr\leq\phi(r)$ give
\begin{align}
        -\frac{\phi''}{\phi}
        \geq
        \frac{c\delta_*}{T\phi^2},
        \quad
        \rho^2-(\phi')^2
        \geq
        c\delta_*.
\end{align}
It follows from \eqref{eq:gaussian-split-Ricci-rr-exact} and
\eqref{eq:gaussian-split-F-bound} that
\begin{align}
        \Ric_{\widetilde g}(\partial_r,\partial_r)
        \geq
        \frac{c\delta_*}{T\phi^2}
        -
        \frac{CM^2}{T^2\phi^2},\quad 
        \left|
        \Ric_{\widetilde g}(\partial_r,X)
        \right|
        \leq
        \frac{CM}{T\phi^2}|X|_{h_r},\quad
        \Ric_{\widetilde g}(X,X)
        \geq
        \left(
        \frac{c\delta_*}{\phi^2}
        -
        \frac{CM}{T\phi^2}
        \right)|X|_{h_r}^2.
\end{align}
Let $Y=a\partial_r+X$, where $X$ is tangent to $M_r$. Young's inequality
gives
\begin{align}
        \Ric_{\widetilde g}(Y,Y)
        &\geq
        \frac{c\delta_*}{4T\phi^2}a^2
        +
        \frac{c\delta_*}{4\phi^2}|X|_{h_r}^2
        \geq
        0
\end{align}
provided that
\begin{align}\label{eq:gaussian-split-T-choice}
        T
        \geq
        \frac{CM^2}{\delta_*^2}.
\end{align}
Outside $\cA_T$, we have $F_r=0$ and $u(r)\in\{0,1\}$. Hence
\eqref{eq:gaussian-split-Ricci-rr-exact} reduces to
\begin{align}
        \Ric_{\widetilde g}(\partial_r,\partial_r)
        &=
        -n\frac{\phi''}{\phi}
        \geq
        0,
        \quad
        \Ric_{\widetilde g}(\partial_r,X)
        =
        0,
        \notag\\
        \Ric_{\widetilde g}(X,X)
        &=
        \Ric_{g_{u(r)}}(X,X)
        -
        \left(
        (n-1)(\phi')^2+\phi\phi''
        \right)g_{u(r)}(X,X).
\end{align}
If $u(r)=0$, then
$\Ric_{g_0}=(n-1)g_0$ and $\phi'\leq1$. If $u(r)=1$, then
$\Ric_{\bar g}\geq(n-1)\rho^2\bar g$ and
$\phi'\leq s+\delta_*/4<\rho$. Since $\phi''\leq0$, the tangential
component is nonnegative in either case. We conclude that $\Ric_{\widetilde g}\geq 0.$
\paragraph{$\mathrm{(iii)}$ \textit{Construction of the potential}} Define $\widetilde V$ by
\begin{align}\label{eq:gaussian-split-Vtilde-definition}
        \widetilde V'(r)
        :=
        \alpha\frac{\phi(r)}{\phi'(r)},
        \qquad
        \widetilde V(0):=0.
\end{align}
Since $\phi'\geq s>0$, the function $\widetilde V$ is smooth. Near the
origin, $\phi(r)=r$, and hence
$\widetilde V(r)=\alpha r^2/2$. In the radial direction,
\begin{align}\label{eq:gaussian-split-Vtilde-radial-Hessian}
        \nabla_{\widetilde g}^2\widetilde V(\partial_r,\partial_r)
        =
        \widetilde V''
        =
        \alpha
        \left(
        1-
        \frac{\phi\phi''}{(\phi')^2}
        \right)
        \geq
        \alpha
        >
        1.
\end{align}
For a vector $X$ tangent to a slice,
\begin{align}
        \nabla_{\widetilde g}^2\widetilde V(X,X)
        =
        \widetilde V'\mathrm{II}(X,X).
\end{align}
Using \eqref{eq:gaussian-split-shape-operator},
\eqref{eq:gaussian-split-F-bound}, and $\phi'\geq s$, we obtain
\begin{align}
        \mathrm{II}(X,X)
        \geq
        \left(
        \frac{\phi'}{\phi}
        -
        \frac{CM}{T\phi}
        \right)|X|_{h_r}^2
        \geq
        \left(
        1-
        \frac{CM}{sT}
        \right)
        \frac{\phi'}{\phi}|X|_{h_r}^2.
\end{align}
Therefore,
\begin{align}\label{eq:gaussian-split-Vtilde-tangential-Hessian}
        \nabla_{\widetilde g}^2\widetilde V(X,X)
        \geq
        \alpha
        \left(
        1-
        \frac{CM}{sT}
        \right)|X|_{h_r}^2.
\end{align}
Since $\alpha>1$, we may increase $T$, while preserving
\eqref{eq:gaussian-split-T-choice}, so that
\begin{align}
        \alpha
        \left(
        1-
        \frac{CM}{sT}
        \right)
        \geq
        1.
\end{align}
Since $\widetilde V=\widetilde V(r)$, we have
\begin{align}
        \nabla_{\widetilde g}\widetilde V
        =
        \widetilde V'\partial_r.
\end{align}
Moreover, since $\widetilde g=\ud r^2+h_r$, the curves
$r\mapsto(r,\omega)$ are unit-speed geodesics, and hence
$\nabla_{\partial_r}^{\widetilde g}\partial_r=0$. Therefore, for every
vector $X$ tangent to $M_r$,
\begin{align}
        \nabla_{\widetilde g}^2\widetilde V(\partial_r,X)
        =
        \la
        \nabla_{\partial_r}^{\widetilde g}
        (
        \widetilde V'\partial_r
        ),
        X
        \ra_{\widetilde g}=
        \widetilde V''
        \la\partial_r,X\ra_{\widetilde g}
        +
        \widetilde V'
        \la
        \nabla_{\partial_r}^{\widetilde g}\partial_r,
        X
        \ra_{\widetilde g}=
        0.
\end{align}
Now let $Y=a\partial_r+X$, where $X$ is tangent to $M_r$. Since
$\widetilde g(\partial_r,X)=0$, it follows from
\eqref{eq:gaussian-split-Vtilde-radial-Hessian} and
\eqref{eq:gaussian-split-Vtilde-tangential-Hessian} that
\begin{align}
        \nabla_{\widetilde g}^2\widetilde V(Y,Y)
        =
        a^2
        \nabla_{\widetilde g}^2\widetilde V
        (\partial_r,\partial_r)
        +
        2a
        \nabla_{\widetilde g}^2\widetilde V
        (\partial_r,X)
        +
        \nabla_{\widetilde g}^2\widetilde V(X,X)\geq
        \alpha a^2+|X|_{h_r}^2\geq
        a^2+|X|_{h_r}^2=
        |Y|_{\widetilde g}^2.
\end{align}
Consequently, $\nabla_{\widetilde g}^2\widetilde V \geq \widetilde g.$ By \eqref{eq:gaussian-split-phi-end}, after setting $R=r+b$, we have, outside
a compact set,
\begin{align}\label{eq:gaussian-split-gtilde-end}
        \widetilde g
        =
        \ud R^2+s^2R^2\bar g,
        \qquad
        \widetilde V
        =
        \frac{\alpha R^2}{2}+c_*
\end{align}
for some constant $c_*\in\bbR$.

\paragraph{$\mathrm{(iv)}$ \textit{Shrinking the cap and transferring the test space}} Let
$\Phi_\eta:\bbR^{n+1}\to\bbR^{n+1}$ be a radial dilation map $\Phi_\eta(r,\omega)=(r/\eta,\omega)$. Next, define the rescaled metric $g_\eta :=\eta^2\Phi_\eta^*\widetilde g$ and potential $V_\eta:=\eta^2\Phi_\eta^*\widetilde V.$
The metric $g_\eta$ is smooth and complete. Moreover,
\begin{align}
        \Ric_{g_\eta}
        \geq
        0,\quad 
        \nabla_{g_\eta}^2V_\eta
        =
        \eta^2\Phi_\eta^*
        (
        \nabla_{\widetilde g}^2\widetilde V
        )
        \geq
        \eta^2\Phi_\eta^*\widetilde g
        =
        g_\eta.
\end{align}
This proves \eqref{eq:gaussian-split-smoothed-curvature}. Choose $r_*=e^{4T}$, then $u(r)=1$ and $\phi(r)=s(r+b)$ for every $r\geq r_*$. Then when  $r>\eta r_*$,
\begin{align}
        g_\eta
        =
        \ud r^2+s^2(r+b\eta)^2\bar g,\quad 
        V_\eta
        =
        \frac{\alpha(r+b\eta)^2}{2}+c_\eta,
\end{align}
where $c_\eta:=\eta^2c_*$. Set
\begin{align}
        R_0
        :=
        r_*+b,
        \quad
        K_\eta
        :=
        \overline{B_{\eta r_*}(0)},
\end{align}
where the ball is taken with respect to the Euclidean metric. Define
\begin{align}
        \Psi_\eta:
        (R_0\eta,\infty)\times\bbS^n
        &\rightarrow
        \bbR^{n+1}\setminus K_\eta,
        \\
        \Psi_\eta(R,\omega)
        &:=
        (R-b\eta)\omega.
\end{align}
Then $\Psi_\eta$ is a diffeomorphism. It extends smoothly to
$[R_0\eta,\infty)\times\bbS^n$, and its extension maps
$\{R_0\eta\}\times\bbS^n$ onto $\partial K_\eta$. Finally,
\begin{align}\label{eq:gaussian-split-geta-end}
        \Psi_\eta^*g_\eta
        =
        \ud R^2+s^2R^2\bar g,\quad 
        \Psi_\eta^*V_\eta
        =
        \frac{\alpha R^2}{2}+c_\eta.
\end{align}
We now define the rescaled space of polynomials $E_\eta$. Choose a smooth function $\zeta_\eta:[0,\infty)\to[0,1]$ such that
$\zeta_\eta(R)=0$ for $R\leq\frac32R_0\eta$,
$\zeta_\eta(R)=1$ for $R\geq2R_0\eta$, and
$|\zeta_\eta'(R)|\leq C\eta^{-1}$ for every $R\geq0$. For $p\in E$ and
$(R,\omega)\in(0,\infty)\times\bbS^n$, set
\begin{align}
        p(R,\omega)
        :=
        p(R\omega),
\end{align}
where on the right-hand side $p$ is regarded as an ordinary polynomial
on $\bbR^{n+1}$. Define $T_\eta:E \rightarrow C^\infty(\bbR^{n+1})$ by
\begin{align}\label{eq:gaussian-split-Teta-definition}
        (T_\eta p)(x)
        :=
        \begin{cases}
        \zeta_\eta(R)p(R,\omega),
        &
        \text{if }
        x=\Psi_\eta(R,\omega)
        \in\bbR^{n+1}\setminus K_\eta,\\
        0,
        &
        \text{if }
        x\in K_\eta.
        \end{cases}
\end{align}
Note that since $\zeta_\eta(R)=0$ whenever $R_0\eta
        <
        R
        \leq
        \frac32R_0\eta$, $T_\eta p$ vanishes on an open collar of $\partial K_\eta$ and thus the transition between the two regions in \eqref{eq:gaussian-split-Teta-definition} is smooth. Moreover, $T_\eta p$ has polynomial growth and therefore, $T_\eta p \in W^{1,2}(\bbR^{n+1},g_\eta,e^{-V_\eta}\dvol_{g_\eta}).$ Define, $E_\eta:=T_\eta(E)$. If $T_\eta p=0$, then $p$ vanishes on the
nonempty open set $\{R>2R_0\eta\}$. Since $p$ is an ordinary polynomial,
it follows that $p\equiv0$. Thus $T_\eta:E_C\to E_\eta$ is injective, and
hence $\dim E_\eta = \dim E_C = N_q+1$. 

Let $p\in E_C$. Writing $p$ as a finite sum of homogeneous polynomials implies
\begin{align}\label{eq:gaussian-split-polynomial-tip-bound}
        |p|+|\nabla p|_{g_C}
        \leq
        C_p
        \qquad
        \text{for }0<R\leq1,
\end{align}
where $C_p>0$ is a constant depending on the polynomial $p.$ Using $\dvol_{g_C}=s^nR^n\,\ud R\,\dvol_{\bar g}$ and
$|\zeta_\eta'|\leq C\eta^{-1}$, we obtain
\begin{align}
        \int_{\{R\leq2R_0\eta\}}
        (p^2
        +
        |\nabla p|_{g_C}^2)e^{-V_C}\dvol_{g_C}
        =
        O(\eta^{n+1}),\quad 
        \int_{\{R_0\eta\leq R\leq2R_0\eta\}}
        |\zeta_\eta'|^2p^2e^{-V_C}\dvol_{g_C}
        =
        O(\eta^{n-1}).
\end{align}
The cross term arising from $\nabla(\zeta_\eta p)=\zeta_\eta\nabla p+p\nabla\zeta_\eta$ tends to zero by the Cauchy--Schwarz inequality. Since $n\geq4$, both displayed error terms also tend to zero as $\eta\to 0$. Choose a basis $p_1,\ldots,p_m$ of $E_C$, where $m=\dim E_C$. Define the matrices
\begin{align}
        G_\eta^{ij}
        :=
        e^{c_\eta}
        \int_{\bbR^{n+1}}
        T_\eta p_iT_\eta p_j
        e^{-V_\eta}\dvol_{g_\eta},\quad 
        D_\eta^{ij}
        :=
        e^{c_\eta}
        \int_{\bbR^{n+1}}
        \la\nabla T_\eta p_i,\nabla T_\eta p_j\ra_{g_\eta}
        e^{-V_\eta}\dvol_{g_\eta}.
\end{align}
The preceding estimates and \eqref{eq:gaussian-split-geta-end} imply that as $\eta \to 0$ we have
\begin{align}\label{eq:gaussian-split-L2-matrix-convergence}
        G_\eta^{ij}
        \rightarrow
        G_0^{ij}
        :=
        \int_Cp_ip_j e^{-V_C}\dvol_{g_C},\quad 
        D_\eta^{ij}
        \rightarrow
        D_0^{ij}
        :=
        \int_C
        \la\nabla p_i,\nabla p_j\ra_{g_C}
        e^{-V_C}\dvol_{g_C}.
\end{align}
For $a=(a_1,\ldots,a_m)\in\bbR^m$, set
$p_a:=\sum_{i=1}^ma_ip_i$ and
$p_{a,\eta}:=\sum_{i=1}^ma_iT_\eta p_i$. Since $G_0$ is positive definite,
there exists $c_E>0$ such that
\begin{align}
        \sum_{i,j=1}^mG_0^{ij}a_ia_j
        \geq
        c_E|a|^2
        \qquad
        \text{for every }a\in\bbR^m.
\end{align}
Since there are finitely many matrix entries, the corresponding quadratic
forms converge uniformly on $\bbS^{m-1}$. Moreover, since $G_0$ is
positive definite, $G_\eta$ is positive definite for every sufficiently
small $\eta>0$. Consequently, as $\eta \to 0$ we obtain
\begin{align}
        \sup_{0\neq f\in E_\eta}
        \cR_{g_\eta,V_\eta}(f)
        &=
        \sup_{a\in \bbS^{m-1}}
        \frac{
        \sum_{i,j=1}^mD_\eta^{ij}a_ia_j
        }{
        \sum_{i,j=1}^mG_\eta^{ij}a_ia_j
        }
        \rightarrow
        \sup_{a\in \bbS^{m-1}}
        \frac{
        \sum_{i,j=1}^mD_0^{ij}a_ia_j
        }{
        \sum_{i,j=1}^mG_0^{ij}a_ia_j
        }
        =
        \Lambda_E.
\end{align}
Since $\Lambda_E<q$, the above convergence implies that
\begin{align}
        \sup_{0\neq f\in E_\eta}
        \cR_{g_\eta,V_\eta}(f)
        <
        q
\end{align}
for every sufficiently small $\eta>0$. This proves
\eqref{eq:gaussian-split-smoothed-Rayleigh}.

\paragraph{$\mathrm{(v)}$ \textit{Integrability of the potential}} It remains to prove \eqref{eq:gaussian-split-smoothed-finite-mass}. On the exact conical end, \eqref{eq:gaussian-split-geta-end} gives
\[
        e^{-V_\eta}\dvol_{g_\eta}
        =
        e^{-c_\eta}s^n
        e^{-\alpha R^2/2}R^n
        \,\ud R\,\dvol_{\bar g}.
\]
The right-hand side is integrable. The complement of the conical end is compact, and $V_\eta$ and $g_\eta$ are smooth there. Consequently,
\begin{align}
        \int_{\bbR^{n+1}}
        e^{-V_\eta}\dvol_{g_\eta}
        <
        \infty.
\end{align}
This proves \eqref{eq:gaussian-split-smoothed-finite-mass}.
\end{proof}

\begin{proof}[Proof of Theorem~\ref{thm:gaussian-split-main}]
Let $d\geq5$ and set $n:=d-1\geq4$. From Lemma~\ref{lem:gaussian-split-link} we obtain an integer $q\geq2$, which is the degree of the test function $f_q(z,y)=\Re(z^q)$, a cone scaling
parameter $s$, a Ricci curvature lower-bound parameter $\rho$ satisfying $0<s<\rho<1$, a smooth metric $\bar g$ on $\bbS^n$, and
a smooth family of metrics $\{g_\ta\}_{\ta\in[0,1]}$ joining $g_0=g_\can$ to $g_1=\bar g$. Next, apply Lemma~\ref{lem:gaussian-split-polynomial-cone}. We obtain a coefficient $\alpha>1$ for the quadratic potential, the cone $C=(0,\infty)\times\bbS^n$ equipped with the metric $g_C=\ud r^2+s^2r^2\bar g$ and the potential $V_C=\alpha r^2/2$, and a
vector space $E_C$ of polynomial test functions satisfying $\dim E_C=N_q+1$ and
$\sup_{0\neq p\in E_C}\cR_{g_C,V_C}(p)<q$. Here
$N_q=\binom{n+q}{n+1}$ is the number of eigenvalues of the standard Gaussian space, counted with multiplicity, which are strictly less than $q$.
Consequently,
\begin{align}\label{eq:gaussian-split-final-Gaussian-index}
        \lam_{N_q+1}
        (\bbR^{n+1},|\cdot|,\gamma^{n+1})
        =
        q.
\end{align}
Now apply Lemma~\ref{lem:gaussian-split-smoothing} and fix $\eta>0$ sufficiently
small. Set
\begin{align}
        g
        &:=
        g_\eta,
        \qquad
        V
        :=
        V_\eta,
        \qquad
        Z
        :=
        \int_{\bbR^d}e^{-V}\dvol_g,
        \notag\\
        \mu
        &:=
        Z^{-1}e^{-V}\dvol_g,
        \qquad
        K
        :=
        N_q+1.
\end{align}
Since $d=n+1$, Lemma~\ref{lem:gaussian-split-smoothing} shows that $g$ is a
smooth complete metric on $\bbR^d$, $V$ is smooth, $0<Z<\infty$, and
\begin{align}
        \Ric_g
        &\geq
        0,
        \qquad
        \nabla_g^2V
        \geq
        g,
        \notag\\
        \dim E_\eta
        &=
        K,
        \qquad
        \sup_{0\neq f\in E_\eta}
        \cR_{g,V}(f)
        <
        q.
\end{align}
Since $\mu=e^{-W}\dvol_g$, where $W:=V+\log Z$, we have
\begin{align}
        \Ric_g+\nabla_g^2W
        =
        \Ric_g+\nabla_g^2V
        \geq
        g.
\end{align}
Thus $\mu$ is a probability measure and $(\bbR^d,g,\mu)$ satisfies
$\CD(1,\infty)$. Finally,
\begin{align}
        \lam_K(\bbR^d,g,\mu)
        \leq
        \sup_{0\neq f\in E_\eta}
        \frac{
        \int_{\bbR^d}|\nabla f|_g^2\,\ud\mu
        }{
        \int_{\bbR^d}f^2\,\ud\mu
        }
        =
        \sup_{0\neq f\in E_\eta}
        \cR_{g,V}(f)
        <
        q
        =
        \lam_K(\bbR^d,|\cdot|,\gamma^d).
\end{align}
The first inequality follows from the min--max characterization and
$\dim E_\eta=K$. The second equality follows because the normalization factor
$Z^{-1}$ cancels in the Rayleigh quotient, while the last equality follows
from \eqref{eq:gaussian-split-final-Gaussian-index}, $d=n+1$, and
$K=N_q+1$. This proves \eqref{eqn:index-gaussian-split-K} and completes the
proof.
\end{proof}
\begin{cor}\label{cor:gaussian-split-main}
Let $(\R^d,g,\mu)$ be the weighted manifold constructed in Theorem
\ref{thm:gaussian-split-main}, where
$\mu=Z^{-1}e^{-V}\dvol_g$, $\Ric_g\geq0$, and
$\nabla_g^2V\geq g$. Then there does not exist a $1$-Lipschitz map
\[
        T:(\R^d,|\cdot|,\gamma^d)\to(\R^d,g,\mu)
\]
pushing forward $\gamma^d$ onto $\mu$ up to a finite constant.
\end{cor}

\begin{proof}
If such a map existed, then Theorem~\ref{thm:contraction} would imply that
for all $k\geq1$ we have
\begin{align}
        \lam_k(\R^d,g,\mu)
        \geq
        \lam_k(\R^d,|\cdot|,\gamma^d).
\end{align}
This contradicts \eqref{eqn:index-gaussian-split-K}.
\end{proof}

\bibliographystyle{alpha}
\bibliography{refs}
\bigskip
\centerline{\scshape Shrey Aryan}
\smallskip
{\footnotesize
 \centerline{Department of Mathematics, Massachusetts Institute of Technology}
\centerline{77 Massachusetts Avenue
Cambridge, MA 02139-4307, USA}
\centerline{\email{shrey183@mit.edu}}
} 


\end{document}
